% ============================================================
% 微模块2：人工智能赋能科研
% ============================================================

\chapter{AI赋能科研}

%前四章讨论了人工智能的基本概念、技术原理与社会影响。从本章开始，视角从``认识AI''转向``使用AI''——进入微模块二：人工智能赋能科研。

科研活动范式的发展脉络可划分为三个演进阶段。

最早是\textbf{经验驱动}范式：研究者基于个人积累与理论推演提出假设，设计实验加以验证。随后进入\textbf{数据驱动}范式：高通量测序、工业传感器和卫星遥感使数据获取成本断崖式下降，配合统计学习与数据挖掘工具的普及，研究者开始从海量数据中直接发现规律。

当前正在经历的是第三个阶段：\textbf{AI驱动}范式。AI（Artificial Intelligence，人工智能）并非替代前两种范式，而是在这两者的基础之上叠加了一个新的能力层——它不直接提出科学问题，也不替代实验本身，而是加速了科研活动中那些高度依赖语言、推理和综合判断的环节：文献检索、方案设计、写作润色和审稿回复。这些曾经消耗研究者大量精力的工作，如今可以被显著压缩，使研究者能将更多时间投入到真正创造性的思考中。三大科研范式的演进关系如图\ref{fig:ch5_01}所示。

\begin{figure}[H]
\centering
\BookFigureImage{images_chapter_05/三大科研范式演进图.png}
\caption{三大科研范式演进图}
\label{fig:ch5_01}
\end{figure}
\par

本章的核心线索是：\textbf{AI在科研中的价值不在于替代研究者的判断，而在于加速那些重复性最高、最消耗时间的环节}。以下七节以硕士生从开题到答辩的典型路径为线索，逐步展示AI在各环节中的具体用法与能力边界：

\begin{itemize}[nosep]
\item \textbf{5.1 AI工作原理与领域导航}：从逐词预测到知识树构建，理解AI的能力边界并用它快速建立陌生领域的方向感；
\item \textbf{5.2 AI辅助文献检索与阅读}：从语义检索到六维度精读，用AI从几百篇论文中高效锁定核心文献；
\item \textbf{5.3 AI辅助课题设计与项目申报}：从空白识别到假设推导，厘清AI在选题中能做什么、不能做什么；
\item \textbf{5.4 AI辅助论文写作}：从引言到摘要，理解``人写AI改''为何是必须坚持的铁律；
\item \textbf{5.5 AI辅助语言润色与翻译}：从词汇对等到功能对等，识别学术翻译中最隐蔽的陷阱；
\item \textbf{5.6 AI辅助自我审稿与投稿准备}：让AI扮演审稿人，分辨哪些意见可信、哪些不可信；
\item \textbf{5.7 AI辅助多模态信息处理与内容创作}：从海报到基金申请书，掌握信息压缩的边界——什么可以省略、什么绝不能丢弃。
\end{itemize}

本章所有AI对话示例均基于真实交互模式重构，供教学参考。AI输出中出现的引用信息（作者、期刊、文献标题等）应视为需要核实的草稿，读者在使用提示词模板时务必在数据库中独立验证。

\subsubsection*{\textbf{•两条提醒}}

1. 研究者的数据会不会离开研究者可控的范围？

本章演示的诸多操作，包括上传论文全文让AI做结构化提取、把问卷数据交给AI做异常诊断、把未发表论文初稿交给AI审稿，都涉及将研究材料上传到第三方AI平台。在上传之前，研究者需要确认：
\begin{itemize}[nosep]
\item 所在机构是否有关于AI工具使用的明确规定？
\item 上传内容是否包含人类受试者的敏感数据，是否已做脱敏处理？
\item 未发表的论文或基金申请书，所用的AI平台是否承诺不将数据用于模型训练？
\end{itemize}
如果是敏感研究数据或未发表的核心成果，优先考虑在本地部署的开源模型上操作，或至少使用承诺数据不用于训练的API（Application Programming Interface，应用程序接口）服务\footnote{此处指通过云端API调用第三方AI模型的服务模式，数据会离开本地传输至服务商的服务器。}。

2. 不要因为``读起来很顺''就轻信其真实性。

AI的输出擅长模仿学术文本的形式感：术语准确、句式流畅、格式规范。这种``看起来可靠''的表面特征会使研究者不自觉地降低警惕。全章反复出现的一个原则就是：AI给研究者的每一条输出，不管是文献信息、统计数据还是方法建议，都先视为需要核实的草稿。核实的方法因信息类型而异：文献去PubMed\footnote{PubMed：美国国立医学图书馆（NLM）下属国家生物技术信息中心（NCBI）维护的生物医学文献免费检索系统，收录超过3600万篇论文。}/Web of Science\footnote{Web of Science：科睿唯安（Clarivate）维护的多学科引文数据库，核心合集覆盖自然科学、社会科学和艺术人文领域的权威期刊。}查，统计数据用G*Power\footnote{G*Power：免费统计功效分析软件，用于计算研究所需的最小样本量，支持t检验、F检验、卡方检验等多种统计检验类型。}或R重算，方法建议去近两年的方法学文献里交叉验证。


\medskip
\textbf{本章后续各节涉及上传论文全文、研究数据或未发表初稿的操作时，均适用上述两条提醒。这些原则适用于每一次上传，正文中不再逐一重复。}

两条提醒看似在划定约束，实质是在界定分工的边界：AI负责加速，研究者负责把关。带着这层意识，读者对每一个AI输出都会多一分审视，少一分盲信。

\medskip
学完七节之后，摆在研究者面前的问题不再是``AI能不能用''，而是\textbf{哪些环节交给AI加速，哪些判断必须留给自己}。这个选择本身，正是AI时代研究者核心竞争力的起点。
\newpage

% ============================================================
% 5.1 AI工作原理与领域导航
% ============================================================
\section{AI工作原理与领域导航}
\label{sec:5.1}

\begin{sectionkeypoints}
\item 理解大语言模型的逐词预测机制与文献幻觉的产生原理，据此判断AI在科研中哪些任务可靠、哪些不可靠
\item 用「知识树构建法」在几分钟内建立陌生领域的方向感
\item 用「领域资源三查法」锁定核心期刊、研究团队和学科争议
\item 用「综述对比四格法」从多篇综述中提取共识、争议和研究空白
\item 掌握提示词优化的核心循环：提需求→得输出→审视→改提示词再问
\end{sectionkeypoints}

在进入具体操作之前，需要先说明一个贯穿本章的工作习惯。初次使用者常见的误区是：输入一个问题便期待一个完美的答案。结果往往令人失望：AI要么回答得过于笼统，要么生成的内容看似合理实则虚假。

真正有效的用法是一个循环：研究者提需求→AI给输出→研究者审视→研究者改提示词再问→拿到更好的输出→再审视→再改。AI不是一个全知全能的学者，而是一个在海量文本上训练出来的统计模型——它的知识来自训练数据中的统计规律，而非对学术文献的真正理解。在未被提供充分上下文时，AI并不了解研究者的具体研究处境。研究者的任务不是给它下一道完美的命令，而是通过一轮一轮的追问，逐步把模型训练数据中对研究者有用的信息引导出来。

在这个循环里，有三件事在任何情况下都只属于研究者、不能交给AI。

第一，判断一个发现``对不对''、``值不值得做''。AI没有面向未来的知识，也理解不了研究者的学科共同体当下的评价标准，这个判断只能由研究者来做。

第二，决定审稿意见中``在哪让步、在哪坚持''。AI起草回复时习惯全盘接受所有意见，但这个战略性判断只有研究者能做。

第三，让未经核实的引用进入自己的论文。AI推荐的文献可能是编造的，AI审稿时指控研究者``漏引''的文献也可能根本不存在，每一条引用都必须由研究者亲手在数据库里确认。

\subsubsection*{\textbf{•大语言模型的工作机制与能力边界}}

读者需要理解以下问题：大语言模型到底怎么``回答''研究者的提问？

大语言模型本质上是一个逐词预测器。研究者输入一段话，它预测下一个最可能接什么词；预测完了拼接上去，再预测下一个。循环往复，直到生成一段完整回复。这并非某种神秘的``思考''，而是纯粹的统计预测。人们在手机上打字时，输入法猜测下一个可能出现的词，原理相同，只是大语言模型将这一机制扩展到了极大规模。\footnote{严格来说，大语言模型预测的基本单位是token（子词），而非完整的词——例如``running''可能被拆分为``run''和``ning''两个token分别预测。本书为便于非计算机专业读者理解，统一使用``逐词预测''这一简化表述，但读者需知晓其技术实质是逐token预测。}

仅凭这一机制，就能解释研究者在科研中使用AI时遇到的绝大多数困惑。

为什么AI能帮研究者理清领域脉络，但不能把AI生成的文献信息直接视为检索结果？因为在未启用检索增强功能的情况下，它并未实际检索PubMed，它只是算出了``这个问题后面最常接什么讨论''。它能告诉研究者这个领域大概长什么样，但它没见过训练截止之后发表的任何论文。

为什么AI给出的答案有时高度精准、有时又完全错误？因为它的知识停留在训练截止那天，之后的新发现它完全无法知晓。遇到知识空白，它不会用``我不知道''来回答：人类写下的文本里，给出答案的例子远多于承认无知的例子。于是模型会生成一个语法正确、措辞专业、逻辑通顺，但事实上并不存在的答案。这就是幻觉（hallucination），不是偶发故障，而是模型工作机制的必然结果。后文将专门讨论应对策略。

核心原则：把AI每次给出的答案视为需要核实的草稿，而非可以引用的定论。它是辅助工具，不是知识权威。

理解了上述工作机制之后，另一个问题便清晰起来：同样的模型，为什么有人用得得心应手，有人觉得它难以使用？差别不在模型，在提示词。提示词并不神秘：它的作用是为模型设定清晰的任务边界。边界明确，模型输出精准；边界模糊，模型输出散乱。

逐词预测、训练截止与幻觉、提示词设边界——这三者不是彼此独立的原理，而是贯穿本节所有操作的底层逻辑：知识树利用的是共现统计，资源列表依赖的是模式记忆，综述对比本质上是跨文本归纳。每当AI输出一份看似权威的答案，这三条原理就是研究者判断它可信度的第一道防线。带着这个框架进入接下来的操作，读者会反复看到同一个规律：AI能帮研究者加速，但判断必须由研究者亲手完成。

\begin{funbox}
\textbf{大语言模型的基石：Transformer架构的诞生。}2017年，谷歌研究团队的八位科学家在神经信息处理系统大会（NeurIPS）上发表了一篇题为``Attention Is All You Need''的论文，提出了一种全新的神经网络架构——Transformer。这篇论文最初的目标非常具体：改进机器翻译的质量。然而，Transformer中的核心创新——``自注意力机制''——使得模型能够在处理一个词时同时``关注''输入文本中的所有其他词，从而以前所未有的精度捕捉长距离的上下文依赖关系。

这一特性远超机器翻译的需求。短短数年间，基于Transformer的GPT系列（生成式预训练Transformer）和BERT（来自Transformer的双向编码器表示）等模型相继问世，将逐词预测这一简单原理扩展到千亿甚至万亿参数的规模，催生了今天研究者所使用的大语言模型。可以说，一个翻译任务的副产品，以无人预料的方式改变了人工智能的轨迹——这也恰好说明了一个道理：基础研究的突破往往不以应用为目标，但它们一旦发生，影响远超最初的想象。
\end{funbox}

% --- 领域入门：三个研究者的共同起点 ---
\subsection{陌生领域与入门困境}

进入一个新领域，第一步往往是理清这个领域的基本格局：有哪些子方向、大家在用什么方法、目前争论什么。但面对几百篇论文，从何入手？以下三个案例来自同一个困境的不同侧面。

小林本科专业是GIS\footnote{GIS（Geographic Information System，地理信息系统）：用于采集、存储、分析和管理地理空间数据的计算机系统。}和遥感：卫星影像、空间分析、地图制图。导师给的新课题却是``城市绿地与居民心理健康''。她打开Web of Science检索后发现，这个领域的论文使用了一系列她不熟悉的指标，包括NDVI\footnote{NDVI（归一化植被指数）：基于卫星遥感影像计算的植被覆盖程度指标，数值范围通常为-1到1，正值越大表示植被越茂密。}、绿视率、公园可达性，每个指标背后都有一套测量方法和文献传统。

健康结局方面同样复杂：有人用PHQ-9量表\footnote{PHQ-9（病人健康问卷-9）：国际通用的抑郁症状自评量表，共9个条目，总分0-27分，通常以$\geq$10分作为抑郁症状阳性的筛查阈值。}测抑郁，有人看抗抑郁药处方记录，还有人测唾液皮质醇。研究设计也五花八门：横截面的、队列的、自然实验的。她勉强读了几篇，读完反而更迷茫了：这些论文之间到底是什么关系？

王工的处境类似。他本科学机械电子工程，导师让他做``基于深度学习的旋转机械智能故障诊断''。这一方向看似具体，但他在IEEE Xplore\footnote{IEEE Xplore：电气电子工程师学会（IEEE）维护的科技文献数字图书馆，涵盖电气工程、计算机科学和电子工程等领域的期刊、会议论文和技术标准。}里搜``bearing fault diagnosis deep learning''，返回的论文铺天盖地：小波变换、经验模态分解、CNN（Convolutional Neural Network，卷积神经网络）\footnote{CNN：一种擅长提取局部空间特征的深度学习架构。}、LSTM（Long Short-Term Memory，长短期记忆网络）\footnote{LSTM：一种擅长建模时间序列依赖关系的循环神经网络变体。}、Transformer\footnote{Transformer：一种基于自注意力机制的神经网络架构，近年来在自然语言处理和故障诊断领域均有广泛应用。}、图神经网络……每种方法背后都有上百篇论文，彼此之间的关系他完全理不清。更为棘手的是，有些论文说CNN好，有些说Transformer才是未来，他一个刚入门的硕士生，该信谁的？

张博遇到的困难又有所不同。他本科学分子生物学，之前跟随师兄从事过lncRNA\footnote{lncRNA（Long Non-Coding RNA，长链非编码RNA）：长度超过200个核苷酸、不编码蛋白质的RNA分子，在基因表达调控中发挥多种功能。}相关研究，现在导师让他转向circRNA\footnote{circRNA（Circular RNA，环状RNA）：通过反向剪接（back-splicing）形成的共价闭合环状RNA分子，具有高稳定性和组织特异性，在基因调控中发挥重要作用。}。circRNA这个领域2015年之后随着高通量测序技术的成熟迅速发展：有人研究它怎么在细胞里生成的，有人研究它怎么像海绵一样吸附miRNA\footnote{miRNA（MicroRNA，微小RNA）：长度约22个核苷酸的小型非编码RNA，通过与靶mRNA的3$'$非翻译区碱基配对来调控基因表达。}，有人研究它跟胃癌增殖侵袭的关系，还有人已经在做液体活检的临床转化了。

PubMed里搜``circRNA+gastric cancer''，返回几百篇论文。每一篇都在讲一个不同的故事，但张博不知道这些故事之间是怎么连起来的。

三个人的问题本质完全相同：需要的不是把某一篇论文读透，而是先理清这个领域的整体格局。

要回答这个问题，第一个操作是让AI画一张领域草图。耗时约五分钟，不替代深度阅读，但能让研究者知道接下来要读的东西大概落在哪个位置。

% --- 知识树构建：方法论 + 对话演示 ---
\subsection{知识树构建法}

传统做法是找一篇高被引综述来研读。但综述从动笔到发表通常滞后两三年，且每篇综述只从特定角度切入。知识树的价值在于：在精读综述之前，先帮研究者建立对整个领域结构的初步把握。

\subsubsection*{\textbf{•知识树构建的操作逻辑}}用AI画知识树，本质上利用了训练数据中的共现统计：某个领域论文里反复一起出现的术语和概念，模型能自动归纳出主从关系。这不是理解，是统计。这决定了它的两个核心特征：主流方向（论文多、模式强）AI的归纳较为准确，冷门交叉方向（论文少、模式弱）则容易被遗漏。

因此，知识树不是学科地图，而是方向感工具。正确的使用方法分三步：①让AI构建毛坯框架（3-4个子领域+每个子领域的核心关键词）；②研究者自己对照一篇高被引综述检查框架是否合理；③重点关注AI漏掉的部分，那往往是研究者真正能做出贡献的交叉地带。

以下通过小林的案例展示具体操作方法。她把提示词写成``拆分为3-4个主要子领域、每个列出核心关键词和2-3个里程碑方向''，这是关键：不加数量限定，AI可能给研究者拆解出七八个粒度不均的子方向，快速建立方向感就无从谈起了。

\begin{userbox}
我正在入门``城市绿地与居民心理健康（抑郁、焦虑、幸福感）''这个研究领域。请帮我构建一个知识树，把这个领域拆分为3-4个主要子领域。每个子领域列出核心关键词和2-3个里程碑式的研究方向。输出用缩进列表形式，层级清晰，不要使用超长句子。
\end{userbox}

\begin{aibox}
好的，这个领域大致分四块，我给你搭个框架：

\begin{itemize}[nosep]
\item \textbf{绿地暴露怎么测}——主要有三种路径：卫星遥感做NDVI（最常用，数据好拿但粒度粗）、街景图像做绿视率（更接近人眼感受，但获取成本高）、GIS做公园可达性（缓冲区分析）
\item \textbf{健康结局用什么指标}——自评量表（PHQ-9、GAD-7，后者为广泛性焦虑障碍量表，$\geq$10分为临床阈值）最主流，也有用临床诊断数据和生理指标（皮质醇等）的，后两种门槛高、论文少
\item \textbf{中间机制是什么}——体力活动、心理恢复、社会交往、空气净化，这几条路径都有文献支持，但很少有研究放在同一个模型里同时比
\item \textbf{怎么处理混杂}——SES是最大的混淆因素，纵向设计和自然实验比横截面更有说服力但难度大很多
\end{itemize}

这是骨架，需要我展开哪个方向？
\end{aibox}

\begin{critiquebox}
不宜期待AI绘制的知识树具有很高的精确度，这个领域有经验的研究生都能画得更准。它的价值不在于精确，而在于速度。小林拿到输出后立刻可以判断，接下来要读的论文大致分为四类，可以分块阅读，无需同时消化所有内容。

但拿到之后不要全盘接受。小林迅速发现``中间机制''那部分遗漏了一个重要方向：空气污染作为竞争性中介。绿地对心理健康的影响，有一部分可能是通过降低空气污染来实现的，这和``多运动''是完全不同的两条路径。AI未能列出。
\end{critiquebox}


\subsubsection*{\textbf{•知识树在不同学科中的变体}}三个人面对的是同一个操作，但每个人的知识树形态完全不同，取决于学科的知识组织方式。以下通过王工和张博的案例展示两个要点：知识树在不同学科中的具体形态，以及各自最常见的误判陷阱。

王工（机械故障诊断）需要的不是``领域分几块''，而是一张\textbf{方法谱系图}，涵盖从数据获取到故障分类的完整技术链路。

工科领域用AI画知识树时，有一个常见陷阱：AI倾向于把不同时代的方法并列展示，而不标注演进关系。比如传统方法SVM（Support Vector Machine，支持向量机）、上一代深度学习方法（CNN）和当前热门方法（Transformer）可能被放在同一个层级，但读者需要自己补上``从SVM到CNN再到Transformer是一个逐步演进的过程''这个视角。

此外，工科领域里``特征提取''和``模型训练''的边界在实际研究中越来越模糊（很多端到端模型跳过了手工提特征这一步），AI画的知识树可能反映的是几年前的学术共识。

张博（circRNA与胃癌）需要的是\textbf{生物学研究链图谱}，涵盖从基因表达谱到分子机制到动物表型到临床转化的完整研究链路。获取AI输出的链路后，有两个方面需要格外警惕。一是AI举例时偏好引用最经典的案例（如ciRS-7吸附miR-7），但经典案例可能来自不同的生物系统：ciRS-7的发现是在脑组织中，miR-7在胃癌里的功能跟脑组织不完全一样，直接套用会出问题。二是AI容易混淆数据来源的性质：比如GEO\footnote{GEO（Gene Expression Omnibus，基因表达综合数据库）：美国国家生物技术信息中心（NCBI）维护的公共功能基因组学数据库，收录全球研究者提交的高通量基因表达和基因组数据。}是研究者各自上传的归档库（质量参差不齐），TCGA\footnote{TCGA（The Cancer Genome Atlas，癌症基因组图谱）：由美国国家癌症研究所（NCI）和国家人类基因组研究所（NHGRI）联合发起的 landmark 项目，提供33种癌症类型的多维基因组数据。}是标准化项目（数据可比性强），AI在列数据资源时可能把它们并列展示而不加区分。

\medskip
纵观三个学科，可以发现三位研究者做了同一件事：拿到AI的知识树之后，都在找AI漏了什么。小林补了``空气污染中介''，王工发现``端到端学习让特征提取边界模糊''，张博指出``GEO和TCGA的数据性质完全不同''。

三个案例指向同一个结论：AI搭建毛坯框架，研究者精加工。关注的不是AI覆盖了什么，而是它遗漏了什么——那才是真正的贡献所在。

\subsubsection*{\textbf{•从模糊兴趣到可检索的问题}}还有一种常见的入门困境：导师给了一个大方向（比如``数字化转型对就业的影响''），但研究者连这个方向在学术上叫什么名字都不确定，遑论去数据库检索文献。这时候可以先让AI帮研究者把模糊的兴趣方向翻译成学术语言：让它提3-5个可能的学术命名方案，选最接近研究者描述的那个，再用知识树展开。关于获得检索式之后如何在数据库中执行，详见下一节；把日常生活语言转换成学术关键词这一步，可以在本节的知识树阶段就完成。对应的提示词模板1-4见本节末尾。

% --- 资源锁定：期刊、团队与争议 ---
\subsection{核心期刊、团队与争议}

领域框架有了，下一步是填入具体信息，厘清该领域的参与者：哪些期刊具有权威性、哪些团队处于领先地位、学术界在争论什么。

\subsubsection*{\textbf{•领域资源三查法}}这个操作需要AI输出的信息分三类：期刊、团队、争议。每一类的可靠性不同，验证方法也不同。

期刊信息（名称、偏好）AI通常可靠，但影响因子和分区是快照不是实时数据，必须去JCR\footnote{JCR（Journal Citation Reports，期刊引证报告）：科睿唯安（Clarivate）每年发布的期刊评价工具，提供影响因子、分区、被引频次等期刊学术影响力指标。}核对。

团队信息（机构、方向）AI仅能提供大致信息，但学术人员流动频繁，带头人可能换了单位甚至退休，需要去Web of Science搜近期发表来确认。

争议信息（学科分歧）AI能归纳主要立场，但措辞往往较为粗糙，``自选择偏误''被简化成``富人选了绿地''这类表述，研究者需要把它润成学术讨论中更准确的表达。

记住这个优先级：\textbf{AI先列→研究者去核实→研究者整理成个人清单。}

以下通过小林的案例展示具体操作方法。

\begin{userbox}
研究领域：城市绿地暴露与居民心理健康（抑郁、焦虑方向）。请用表格列出以下信息：
1. 该领域TOP5国际期刊，标注各刊的风格偏好
2. 全球3-5个主要研究团队，标注机构、侧重方向和标志性贡献
3. 目前该领域3个核心争议或未解决的关键问题，列出争议双方的主要依据

重要限制：仅列出训练数据截止日期之前真实存在的期刊和团队。如果无法确定某个具体信息（如影响因子、带头人姓名），在该单元格标注``[待核实]''，不要编造。核心争议部分需列出各方的主要依据，而非只描述一方观点。
\end{userbox}

\begin{aibox}
期刊清单：

\begin{formalTable}[H]
\begin{tabular}{p{4.5cm} p{8cm}}
\toprule
\textbf{期刊名} & \textbf{风格偏好} \\
\midrule
\textit{Landscape and Urban Planning} & 偏好实证研究，重视城市规划和景观设计应用价值 \\
\textit{Environmental Health Perspectives} & 偏好流行病学设计，对暴露评估方法严格性要求高 \\
\textit{Health \& Place} & 偏好``地方与健康''交叉研究，接受混合方法 \\
\textit{Environment International} & 偏好大规模环境暴露与健康结局研究 \\
\textit{Urban Forestry \& Urban Greening} & 偏好城市绿化生态-社会效益研究，审稿相对友好 \\
\bottomrule
\end{tabular}
\end{formalTable}

\smallskip
核心研究团队（部分示例）：

\begin{formalTable}[H]
\begin{tabular}{p{4.5cm} p{8cm}}
\toprule
\textbf{机构} & \textbf{标志性贡献方向} \\
\midrule
哈佛T.H. Chan公共卫生学院 & 护士健康研究队列中的绿地-健康分析，纵向设计标杆 \\
巴塞罗那全球健康研究所（ISGlobal） & 欧洲多城市绿地-健康研究，空气污染中介效应验证 \\
墨尔本大学 & 澳大利亚纵向队列中绿地与心理健康的长期追踪 \\
\bottomrule
\end{tabular}
\end{formalTable}

三大核心争议：

\begin{enumerate}[label=\arabic*.]
\item \textbf{绿地暴露测量指标的标准化问题。}NDVI（基于卫星遥感）成本低但粒度粗，绿视率（基于街景图像，更贴近人眼感知）精度高但获取成本高。不同指标的研究经常得出不一致的结论，全领域尚无共识。
\item \textbf{中介机制的主次关系。}体力活动、心理恢复、社会交往分别有流行病学和实验证据支持，但鲜有研究在同一模型中同时检验多个竞争性中介路径。
\item \textbf{相关与因果的边界。}大量横截面研究发现绿地与健康的关联，但少数利用自然实验（如城市绿化改造）的因果推断研究得到的效应量通常小于横截面研究，自选择偏误可能被广泛低估。
\end{enumerate}
\end{aibox}

\begin{critiquebox}
小林拿到输出后，用半小时完成三项核实工作。她首先在JCR（期刊引证报告，科睿唯安发布，每年更新）中核实五本期刊的最新影响因子和分区：期刊属实，但影响因子数据可能不是最新的。随后在Web of Science中检索三个团队近两年的发表，确认带头人是否仍在原来机构——AI的信息停在训练截止那天，有人可能已变更单位甚至退休。最后将三个核心争议的表述逐一修正：AI说``自选择偏误''时只举了``高收入群体选了绿地''，但更准确的表述是``个体层面不可观测的特征（健康意识、社会经济资源）同时影响居住选择和健康结局''。经过半小时的核实，组会汇报时所用的期刊和团队信息便均为最新状态。

概而言之：AI提供的机构、人员、影响力数据都是快照，不是实时数据。上述核实工作均须由研究者自行完成。
\end{critiquebox}

对于王工和张博而言，同一操作的核心逻辑不变，但输出内容因学科而异。

王工的工程领域核心资源集中在\textit{Mechanical Systems and Signal Processing}、\textit{IEEE Transactions on Industrial Electronics}、\textit{ISA Transactions}等期刊，核心团队分布在西安交通大学、清华大学、德国帕德博恩大学等机构，核心争议集中在``变工况下模型泛化能力不足''和``深度学习模型的可解释性''两个方向。

张博的生物医学核心资源则是\textit{Molecular Cancer}、\textit{Gastric Cancer}、\textit{Oncogene}等期刊，核心争议集中在``circRNA作为miRNA海绵的功能是否有被高估的风险''和``circRNA生物标志物的组织特异性是否足以支持液体活检应用''。

三个学科的交叉验证方法是一样的：期刊去JCR核实，团队去数据库搜近期发表，争议表述跟最新综述对照。不论哪个学科，核实的原则不变。

期刊和团队理清了，争议方向也知道了，这个过程中研究者手里会积累三五篇该领域的综述论文。下一步不是逐篇研读，而是横向比较，从共识和分歧中找到研究者自己的研究切入点。

% --- 综述对比：横向比较三篇综述 ---
\subsection{综述对比与空白识别}

\subsubsection*{\textbf{•综述对比四格法}}这个操作的价值不在于AI的判断精确与否，而在于它帮研究者从逐篇精读的细节中抽离出来，看到更大尺度的信号。该方法的核心步骤为：把3-5篇综述的摘要和关键段落交给AI，让它填四个格子：

\textbf{可靠共识}——所有综述都同意的结论，写作时一笔带过即可，不用花大篇幅论证。

\textbf{领域争议}——各综述之间存在分歧的结论，正好放进引言里作为研究者自己的研究切入点。

\textbf{有潜力的切入点}——只有一篇综述提过的空白，需自行确认是否为真正的空白。

\textbf{综合判断}——不超过300字的整体评估。

需要警惕的是第四格：AI的综合判断本质上是统计归纳（``出现频率高=创新空间小''），不能替代研究者自己的学术判断。张博手里有3篇circRNA与胃癌综述，他让AI帮他做了一件事。

\begin{userbox}
以下是3篇关于``circRNA在胃癌中作用机制''的综述论文摘要（已上传全文摘要和关键段落）。请横向对比分析：
1. 所有综述都同意的结论是什么？标记为「可靠共识」
2. 各综述之间存在分歧的结论是什么？各方的依据是什么？标记为「领域争议」
3. 哪篇综述提出了其他综述未涉及的研究空白？标记为「有潜力的切入点」
4. 综合判断段落（不超过300字）：这些综述共同揭示的稳定结论和新机会

输出格式：对比表格，然后综合判断段落。
\end{userbox}

\begin{aibox}
【对比表格节选】

\textbf{可靠共识}（3/3篇综述同意）：
\begin{itemize}[nosep]
\item circRNA通过miRNA海绵机制调控胃癌增殖和侵袭行为 → 机制方向已确认，不是``是否''而是``哪个circRNA通过哪条miRNA-靶基因轴发挥作用''的问题
\item 胃癌组织与癌旁正常组织中存在大量差异表达的circRNA → 筛选层面共识稳固
\item circRNA在胃癌化疗耐药（顺铂/5-FU）中发挥功能 → 3篇均引用同一组关键实验证据
\end{itemize}

\textbf{领域争议}（存在分歧的方向）：
\begin{itemize}[nosep]
\item \textbf{circRNA作为液体活检生物标志物的可行性：}综述A认为``circRNA在血清中高度稳定，比线性RNA更适合作标志物''；综述C指出``血清circRNA的低丰度使检测重现性不足，需要更大的临床队列验证''。→ 两者不是``对vs错''，而是对同一证据链的不同评估标准
\end{itemize}

\textbf{有潜力的切入点}（仅1篇综述提出）：
\begin{itemize}[nosep]
\item 综述B在最后一段提出了一个其他2篇未涉及的方向——circRNA编码的短肽在胃癌微环境中的免疫调节功能。该方向2022年才出现首篇胃癌相关报道（文献量<5篇），但circRNA编码能力在其他癌种（胶质母细胞瘤、肝癌）已有初步验证。
\end{itemize}

\medskip
\textbf{综合判断}

circRNA在胃癌中的``miRNA海绵→调控靶基因→影响增殖/侵袭''主线已基本确认，继续在成熟方向上累积``第N个circRNA-第M个miRNA-第K个靶基因''的创新空间有限。真正的机会点可能出现在三个方向：(1) circRNA的非海绵功能（翻译短肽、RBP互作）在胃癌微环境中的角色；(2) 将circRNA从``功能描述''推进到``临床验证''——液体活检的大规模队列研究；(3) 靶向circRNA的核酸药物递送系统的技术突破。
\end{aibox}

\begin{critiquebox}
张博认为AI的总结方向是对的，但``创新空间有限''这个判断需要谨慎对待。AI之所以得出这个结论，可能是因为训练数据中``circRNA-miRNA-靶基因''轴的论文数量远多于其他方向，模型在统计上将``出现频率高''等同于``创新空间小''。但``成熟方向创新空间有限''是一个学术判断，它是否正确取决于该领域中究竟还有多少重要的miRNA-靶基因轴未被鉴定，这个问题AI无法回答，因为它没有面向未来的知识。张博把这个判断视为一个``值得与导师讨论的假设''，而非一个``可以写入论文引言作为立论基础的事实''。
\end{critiquebox}

以小林的环境健康领域为例，3篇综述都认为NDVI与抑郁症状存在负相关是可靠共识，但在``这种关系是否独立于SES\footnote{SES（Socioeconomic Status，社会经济地位）：综合衡量个体或家庭经济和社会地位的指标，通常基于收入、教育程度和职业等因素。}''这个问题上存在分歧。这个分歧比共识更有价值，因为它为后续研究提供了可供检验的假设。

前三个格子的价值排序是：争议 > 独有空白 > 共识。争议是研究者介入学术对话的入口，独有空白是研究者潜在的创新方向，共识是用来在引言里一笔带过的背景。

前面三个操作，知识树、锁定资源、综述对比，核心目的都是让AI向研究者提供信息。但在继续深入之前，必须先处理一个基础性问题：AI给研究者的这些信息，有多少是真实存在的，有多少是它虚构的？如果不在继续深入之前理清这个问题，前面三步的成果随时可能失去根基。

% --- 幻觉应对：机制演示与防范策略 ---
\subsection{文献幻觉与信息核验}

AI编造文献，这是用AI做科研时最令人困扰的问题。

研究者让它列几篇某个领域的核心论文，它返回的东西格式上无可挑剔：作者名是该领域的常见姓氏，期刊是该领域的主流刊物，年份是这个方向的高产期，标题由几个高频关键词拼起来。每一条单独看都像真的；综合来看，在数据库里根本不存在。

这一现象是如何产生的？人类学者推荐文献，是从真实记忆里提取：想起某篇读过的论文，知道它确实存在。LLM（Large Language Model，大语言模型）不这样做。

它只是逐词生成：根据上文``该领域有一些重要文献，例如''，模型判断下一个词应该是个作者名，于是从训练数据里该领域常见姓氏中采样一个；然后是年份，选该领域高产期的；然后是期刊名，选该领域出现频率最高的；最后是标题，用该领域高频关键词拼出一个``看起来像论文标题''的字符串。

每一步在概率上都是合理的，但拼在一起便是一个不存在的实体。可作如下类比：一个熟读药方但不通医术的人，开出了一张格式完美、药材名称规范、剂量合理，但治不了病的、也从来没有被任何医书收录过的处方。

了解了这个机制，不难预见什么时候最容易遇到问题。三种情况尤其危险：

\begin{itemize}[nosep]
\item \textbf{研究者第一次进入一个陌生领域，没有背景知识来识别虚假信息。}此时研究者对领域内的学者姓名、期刊名称、经典文献没有任何积累，AI编造的文献在形式上足以以假乱真。
\item \textbf{研究者要求AI给出精确的外部信息。}如``列出引用最多的10篇论文''——LLM在这种任务上的可靠性远低于概念解释和结构梳理，因为精确列表需要逐条对应事实，而模型只擅长统计归纳。
\item \textbf{研究者在提示词里没有加限制语。}模型面对``请列出核心文献''这种不加限定的要求，默认策略是``给出完整的答案''，而不是``承认自己不确定''。
\end{itemize}

下面通过张博的实际案例加以说明。他不加限定地要求AI推荐文献：

\begin{userbox}
请列出``circRNA在胃癌中作为miRNA海绵''领域的5篇核心文献。
\end{userbox}

\begin{lightquote}
AI返回了5条文献——作者名是该领域的常见姓氏，期刊是该领域的主流刊物，年份是这个方向的高产期，标题由几个高频关键词拼起来。单独看任何一条都像一篇真实存在的高被引论文。为了不给读者留下可能被误引用的假文献条目，此处不展示具体输出文本，但研究者可以自己做一次这个实验——不加任何限制让AI推荐研究者所在领域的核心文献，然后去数据库逐条检索，亲身体验幻觉的形态。

张博逐条在PubMed和Web of Science里检索后发现，5条中无一条能完全匹配真实文献——AI返回的是一份格式完美但内容虚构的清单。
\end{lightquote}

由此可得出明确的底线原则：AI推荐的每一篇文献，研究者都必须自己去数据库里核实。这不是``更保险的做法''，是唯一的正确做法。AI没有事实概念，它只有概率分布。引用需要的是事实。

一个实用经验：在提示词里加一句``不确定的标注[待核实]，不要编造''，能降低幻觉的概率（不能消除）。同时应养成习惯：AI提供的精确引用、统计数字、具体方法，这些是高风险输出，逐条核实；概念解释、结构梳理、语言润色，这些是低风险输出，快速浏览加抽样检查即可。

% --- 提示词模板 ---
\subsection*{提示词模板}

本节涉及的提示词已整理为模板，方括号\texttt{[ ]}里的内容替换为研究者自身的内容即可。

\textbf{模板1-1：构建领域知识树}

\begin{tcolorbox}[promptbox]
我正在入门 [领域名称] 这个研究领域。请帮我构建一个知识树，把这个领域拆分为3-4个主要子领域。每个子领域列出核心关键词和2-3个里程碑式的研究方向。输出用缩进列表形式，层级清晰，不要使用超长句子。如果对某个子领域是否属于该领域不确定，标注出来。
\end{tcolorbox}

适用场景：进入新课题的第一天，系统检索之前。当连领域的学术名称都拿不准时，改用以下版本：

\textbf{模板1-1B：不确定学术名称时的替代版}

\begin{tcolorbox}[promptbox]
我对 [粗糙的兴趣方向描述] 感兴趣，但不确定这个方向在学术上的正式命名。请帮我：1. 提出3-5个可能的学术命名方案；2. 选择最接近我描述的一个，基于它构建子领域知识树。
\end{tcolorbox}

\textbf{模板1-2：识别领域格局}

\begin{tcolorbox}[promptbox]
研究领域：[领域名称]。

请用表格列出以下三部分信息：
1. 该领域TOP5国际期刊，标注各刊的风格偏好和对不同类型数据的友好度
2. 全球3-5个主要研究团队，标注机构、侧重方向和标志性贡献
3. 目前该领域的3个核心争议或未解决的关键问题，列出争议双方各自的主要依据

限制要求：
- 仅列出在训练数据截止日期之前真实存在的期刊和团队
- 无法确定的信息标注``[待核实]''，不要编造
- 核心争议部分不可只描述一方观点
\end{tcolorbox}

使用前提：输出结果须经JCR、Web of Science独立核验后再整理成个人清单。

\textbf{模板1-3：多篇综述横向对比}

\begin{tcolorbox}[promptbox]
以下是 [N] 篇关于 [主题] 的综述论文（已上传/粘贴摘要和关键段落）。请进行横向对比分析：

1. [共识分析]：哪些结论是所有综述都同意的？标记为「可靠共识」
2. [分歧分析]：哪些结论存在分歧？各方依据是什么？标记为「领域争议」
3. [空白分析]：哪篇综述提出了其他综述未涉及的研究空白？标记为「有潜力的切入点」
4. [综合判断]：用不超过300字总结这些综述揭示的稳定结论和新机会

输出格式：先给对比表格，再给综合判断段落。
\end{tcolorbox}

使用前提：已精读3-5篇综述，准备撰写文献综述部分时使用。帮助研究者从逐篇笔记中抽离出来，看到结构性信号。

\textbf{模板1-4：将模糊方向转化为可检索的研究问题}

\begin{tcolorbox}[promptbox]
我对 [模糊的兴趣方向描述] 感兴趣，但不知如何将它转化为可以检索文献的研究问题。请按以下步骤处理：

步骤1：用3-5个该领域的专业关键词或短语重新表述我的兴趣方向
步骤2：用这些关键词组合出3个不同的核心研究问题（每个不超过30字），区分两种类型：``验证已有发现''和``探索潜在创新点''
步骤3：为每个研究问题推荐1-2个检索词组合，可直接用于Web of Science或PubMed
\end{tcolorbox}

使用前提：导师已给出大方向，但尚未确定具体研究问题的阶段。


% --- 常见错误与规避 ---

\begin{mistakebox}

\textbf{把AI推荐的文献直接放入引用列表。}AI生成的引用中可能包含格式完美但事实上并不存在的假文献，每一条都必须在PubMed或Web of Science中独立核实后再使用。

\smallskip

\textbf{把AI画的知识树当成学科分类标准。}知识树是方向感工具而非精确地图——主流方向AI归纳较准，冷门交叉方向容易被遗漏，拿到后需对照一篇高被引综述检查修正。

\smallskip

\textbf{只用一轮对话就接受AI的初次输出。}有效的AI使用是一个循环：提需求→得输出→审视→改提示词再问。反复追问才能逐步把模型训练数据中对研究者有用的信息引导出来。

\end{mistakebox}

\medskip
\subsection*{小结}
本节从大语言模型的逐词预测原理出发，展示了AI辅助科研的基本工作循环。知识树构建法帮助研究者在几分钟内建立方向感，领域资源三查法锁定核心期刊和团队，综述对比四格法从共识和分歧中识别研究空白。贯穿全节的核心原则是：AI输出均应视为需要核实的草稿，研究者的判断不可替代。

\bigskip
\subsection*{课后习题}
\begin{enumerate}
\setlength{\itemsep}{6pt}

\item 为什么说大语言模型是概率预测器而非知识库？这对科研中使用AI有何意味？

\item 选一个陌生领域，用AI构建知识树，审视并指出AI遗漏了哪些方向，以及这些遗漏对你发现研究空白有何启示。

\item 什么是文献幻觉？它的产生机制是什么？列举至少两种有效的防范策略。

\item 知识树构建法、领域资源三查法和综述对比四格法分别解决了领域入门的什么问题？三者之间的使用顺序是什么？为什么？

\end{enumerate}
\newpage

% ============================================================
% 5.2 AI辅助文献检索与阅读
% ============================================================
\section{AI辅助文献检索与阅读}
\label{sec:5.2}

\begin{sectionkeypoints}
\item 理解语义检索与关键词检索的技术原理和适用场景差异
\item 用PICOS框架将模糊的研究兴趣转化为可执行的数据库检索策略
\item 用「三级分类法」批量初筛与「六维度结构化提取法」单篇精读，高效完成从海量文献到核心篇目的筛选
\item 通过「方法学横向对比」判断多篇论文的证据质量差异，避免将不同质量的研究同等对待
\item 搭建个人论文库、数据库和代码库，使检索与精读的成果成为后续写作和审稿时可随时调用的知识体系
\end{sectionkeypoints}

\subsubsection*{\textbf{•语义检索与传统检索的本质区别}}

如5.1所述，LLM本质是逐词预测器——它并不实际查询数据库，直接让它``推荐几篇文献''极易产生幻觉。本节讨论的语义检索则不同：它依赖RAG（Retrieval-Augmented Generation，检索增强生成）机制\footnote{RAG：AI先到外部数据库进行语义检索，将搜到的文献作为参考资料填入上下文，再基于这些资料生成回答。即``先检索，后回答''。}，AI在回答前会先到文献数据库中实际执行一次检索，再将检索结果作为参考资料生成回答。使用内置了RAG的AI检索工具（如Semantic Scholar的AI助手、Elicit、Consensus等），检索结果是真实存在的论文。本节以下讨论均针对后者。

从经验驱动到AI驱动，文献检索的逻辑发生了根本变化：关键词检索依赖术语的精确匹配，语义检索则通过向量相似度发现内容关联。在传统检索流程中，研究者打开Web of Science，在搜索框中输入\texttt{``urban green space'' AND ``mental health''}，按回车键，等待数据库返回结果。这便是关键词检索，其原理并不复杂：将输入的关键词与论文元数据进行逐字比对，匹配成功即返回，匹配失败则不返回。

然而，学术写作中同一概念常有多种表述方式。一篇论文的标题使用residential greenness而非urban green space、psychological well-being而非mental health，关键词检索将无法命中这些文献，即便它们与研究者的关注点高度相关。更为隐蔽的情况是：当相关概念出现在正文而非标题或摘要中，或摘要采用了不同的术语表述，关键词检索同样无法匹配。

语义检索采用了不同的技术路径。它依赖文本嵌入（text embedding），即将任意一段文本（无论是论文标题摘要，还是研究者的查询语句）转换为一个高维向量\footnote{向量（vector）：一段文本被表示为一组数字，通常是几百到上千维。两段文本的向量在空间中距离越近，语义就越相似。}。在这个向量空间中，语义相近的文本会自动聚集，即便它们之间没有任何共享词汇。研究者输入``小区绿化对心情的影响''，与之距离最近的论文向量可能是标题中含有``residential greenness and depressive symptoms''的文献，尽管查询语句与这些标题之间不存在任何共同的关键词。

使用AI进行文献检索时，研究者无需在关键词和布尔运算符的精确选择上投入过多精力，用自然语言将研究问题描述清楚即可：描述越具体，检索结果越全面。但这并不意味着传统检索方法应被弃用。在撰写系统综述和PRISMA\footnote{PRISMA（系统综述和荟萃分析报告规范）：要求完整报告检索策略、筛选流程和排除原因的国际标准。}报告时，关键词检索的可复现性是语义检索无法替代的。两种方法互为补充（图\ref{fig:ch5_03}）。

\begin{figure}[H]
\centering
\BookFigureImage{images_chapter_05/语义检索vs关键词检索原理对比图.png}
\caption{语义检索与关键词检索原理对比图}
\label{fig:ch5_03}
\end{figure}
\par

\begin{funbox}
\textbf{从问答系统到检索增强：RAG的演进简史。}让机器自动回答学术问题，这一设想至少可以追溯到1960年代的早期问答系统——它们依赖手工编写的规则，只能处理极窄的领域。此后半个世纪，问答系统的主流思路始终是``先建知识库，再从中检索''，但手工建库的成本限制了其覆盖范围。

真正的转折发生在2020年。Facebook AI Research（现Meta AI）的Lewis等人在论文``Retrieval-Augmented Generation for Knowledge-Intensive NLP Tasks''中首次提出了RAG架构。其核心洞察简洁而深刻：大语言模型本身是一个``闭卷考试''的考生——它只能凭训练时记住的东西作答，遇到没见过的内容只能猜。RAG让它变成了``开卷考试''——在回答之前，系统先从外部知识库中检索相关文档，然后将检索到的内容作为参考资料一并送入模型。模型不再依赖记忆，而是基于眼前这份刚检索到的资料生成回答。这正是``先检索，后回答''的由来。

RAG的命名本身就揭示了它的两个核心组件：Retrieval（检索）从外部数据库中找到相关文档，Augmented Generation（增强生成）将这些文档作为上下文增强模型的回答能力。两者的结合在2023年后迅速成为AI检索工具的标准架构，催生了Semantic Scholar的AI助手、Elicit、Consensus等学术检索工具。2024年后，RAG进一步分化为更精细的形态：Agentic RAG赋予系统多步推理和自主决策的检索能力，Graph RAG则通过知识图谱增强对实体间关系的理解和检索精度。

对研究者而言，理解RAG的``开卷考试''隐喻意味着两条实用判断：其一，答案的上限取决于检索到的文献质量——如果数据库不完整或向量搜索未能命中关键论文，再强的模型也生成不出正确的答案；其二，传统的关键词检索仍有不可替代之处——它的可复现性和透明性是语义检索无法提供的。这也是为什么在系统综述等关键检索任务中，两种方法需要交叉验证。
\end{funbox}

% --- 检索准备：PICOS框架拆解与检索式生成 ---
\subsection{研究问题的PICOS拆解}

张博最初的操作是许多初学者的典型做法：直接将一个模糊的研究方向输入PubMed，查询``circRNA在胃癌增殖和侵袭中的作用''。系统返回了数百条记录。他随机浏览了其中几篇，发现多数论文仅进行了circRNA表达谱分析，在讨论中顺带提及``可能与胃癌有关''，并非他所需要的机制研究。

在打开数据库之前，需要先将模糊的研究兴趣转化为边界清晰的问题。流行病学和实证研究长期使用的PICOS\footnote{PICOS（Population, Intervention/Exposure, Comparison, Outcome, Study Design）：循证医学和流行病学中用于构建可检索研究问题的标准化框架，分别对应研究对象、干预/暴露、对照、结局和研究设计五个维度。}框架为此提供了有效的结构化工具：

\begin{tcolorbox}[promptbox]
\begin{itemize}
\setlength{\itemsep}{2pt}
\item \textbf{P（Population，研究对象）}：研究对象是谁？是否有年龄、地域、特征限制？
\item \textbf{I/E（Intervention/Exposure，干预或暴露）}：核心干预或暴露因素是什么？
\item \textbf{C（Comparison，对照）}：与什么条件进行比较？
\item \textbf{O（Outcome，结局）}：关注的结果指标是什么？
\item \textbf{S（Study Design，研究设计）}：限定哪些研究类型？
\end{itemize}
\end{tcolorbox}

\subsubsection*{\textbf{•PICOS拆解的操作要点}}使用AI拆解PICOS时，提示词中需明确三项内容：①完整的研究问题（不超过100字，问题越具体，拆解结果越精确）；②要求AI按P/I/C/O/S五个维度分别输出；③要求基于拆解结果为指定数据库生成检索式，并区分``高灵敏度版''（适用于初次海筛，命中量大但噪声较多）和``高特异度版''（适用于筛选核心文献，命中量小但精准度更高）。

获得AI生成的检索式后，不应直接套用。AI可能遗漏领域内的习惯用语（如circRNA领域的研究者常用``sponge''描述miRNA海绵机制，而AI仅给出了``miRNA''），也可能使用了过于宽泛的通配符（如proliferat*会将所有提及PCNA\footnote{PCNA（Proliferating Cell Nuclear Antigen，增殖细胞核抗原）：一种在细胞增殖中起关键作用的蛋白，常作为实验内参，其出现并不意味该研究以细胞增殖为核心问题。}的文献一并纳入）。建议在目标数据库中对检索式的每个组件单独进行测试，确认代表性论文是否被覆盖。

张博现在需要将前述模糊的研究方向转化为可在PubMed中直接执行的检索式。

\begin{userbox}
我的研究问题是：circRNA在胃癌中的异常表达是否通过miRNA海绵机制调控胃癌细胞的增殖和侵袭能力，如果是，关键的下游靶基因和信号通路是什么。

请用PICOS框架拆解这个研究问题：
P：研究对象  I/E：暴露或核心变量  C：对照条件  O：结局指标  S：研究设计类型

然后基于拆解结果，为PubMed生成布尔逻辑检索式。标注``高灵敏度版''（适合初次海筛）和``高特异度版''（适合筛选核心文献）。
\end{userbox}

\begin{aibox}
\textbf{P（研究对象）：}胃癌细胞系（AGS、MGC-803、BGC-823、SGC-7901等）和/或胃癌组织标本。排除：非胃部的消化道肿瘤（食管癌、结直肠癌）

\textbf{I/E（暴露/核心变量）：}circRNA过表达或敲低——通过siRNA/shRNA/CRISPR-Cas13等技术干预

\textbf{C（对照）：}正常胃上皮细胞系（GES-1等）或癌旁正常组织（\(\geq\)5cm距肿瘤边缘）

\textbf{O（结局）：}细胞增殖（CCK-8/EdU/集落形成）、侵袭迁移（Transwell/划痕实验）、凋亡（流式细胞术Annexin V/PI染色）、miRNA靶标及下游靶基因表达变化（实时荧光定量PCR/Western blot）

\textbf{S（研究设计）：}体外细胞功能实验 + 体内裸鼠成瘤实验 + 临床组织样本验证。排除：纯生物信息学预测（未经实验验证的）、仅报告circRNA表达谱而未做功能实验的

\smallskip
\noindent\rule{\linewidth}{0.4pt}
\smallskip

\textbf{PubMed检索式}

【高灵敏度版】（预期命中300--600条）：

{\ttfamily\small\setlength{\parskip}{0pt}
\#1\quad "RNA, Circular"[Mesh] OR circRNA[tiab] OR "circular RNA"[tiab]

\#2\quad "Stomach Neoplasms"[Mesh] OR "gastric cancer"[tiab] OR "gastric carcinoma"[tiab]

\#3\quad "Cell Proliferation"[Mesh] OR "Neoplasm Invasiveness"[Mesh] OR proliferat*[tiab] OR invas*[tiab] OR migrat*[tiab] OR metast*[tiab]

\#4\quad \#1 AND \#2 AND \#3

\#5\quad \#4 NOT ("esophageal"[tiab] OR "colorectal"[tiab])
}

【高特异度版】（预期命中80--150条）：

{\ttfamily\small\setlength{\parskip}{0pt}
\#6\quad "MicroRNAs"[Mesh] OR "miRNA"[tiab] OR "miR"[tiab] OR "ceRNA"[tiab] OR "competitive endogenous RNA"[tiab]

\#7\quad \#1 AND \#2 AND \#3 AND \#6
}
\end{aibox}

\begin{critiquebox}
张博并未直接使用AI生成的检索式。他做了两处修改。其一，高灵敏度版中\#3的proliferat*通配符过于宽泛：``proliferating cell nuclear antigen/PCNA作为内参''这类并非研究增殖的文献也会被命中，因此他将检索限定到了标题字段（``proliferat*[ti] OR invas*[ti]''）。其二，高特异度版中miRNA检索词遗漏了``sponge''``sponging''（该领域研究者常用此表述），他补上了``sponge*[tiab]''。

结论是：AI生成的检索式可作为合理的起点，但不应未经审阅直接使用。研究者对领域术语的习惯用法最为熟悉，经过微调后，命中率和精准度均有明显提升。
\end{critiquebox}

\smallskip
同样的PICOS框架，在小林和王工的研究中产生了完全不同的拆解结果。放在一起对比，即可看出这个框架的通用性。

小林的PICOS拆解：

\smallskip
\noindent\textbf{P：}城市社区居民，未限定年龄段，无临床诊断限制。排除住院病人和机构化护理群体

\noindent\textbf{I：}社区绿地暴露——包括NDVI、绿视率（街景图像）、公园可达性（GIS缓冲区分析）

\noindent\textbf{C：}低绿地暴露或无绿地暴露的社区居民

\noindent\textbf{O：}抑郁症状水平——通过标准化量表测量（PHQ-9、CES-D\footnote{CES-D（Center for Epidemiologic Studies Depression Scale，流调中心抑郁量表）：由20个条目组成的自评量表，总分0-60分，通常以\(\geq\)16分作为抑郁症状的筛查阈值。}），或临床诊断记录（ICD\footnote{ICD（International Classification of Diseases，国际疾病分类）：世界卫生组织（WHO）发布的疾病和健康相关问题的国际标准分类系统，用于临床诊断编码和流行病学统计。}编码）

\noindent\textbf{S：}横截面研究、纵向队列研究、准自然实验。排除纯质性研究、病例报告
\smallskip

王工的PICOS拆解：

\smallskip
\noindent\textbf{P：}旋转机械核心部件——滚动轴承（深沟球轴承、圆柱滚子轴承）、齿轮箱

\noindent\textbf{I：}深度学习故障诊断模型——1D-CNN/LSTM/Transformer，输入为原始振动信号或时频图

\noindent\textbf{C：}传统机器学习方法（SVM、随机森林）+ 手工特征（时域/频域/时频域特征）

\noindent\textbf{O：}故障分类准确率、F1-score\footnote{F1-score：精确率（Precision）与召回率（Recall）的调和平均数，是分类任务中兼顾查准与查全的综合评估指标，取值范围0-1，越接近1表示模型性能越好。}、混淆矩阵、跨工况泛化性能

\noindent\textbf{S：}实验台架数据验证（CWRU\footnote{CWRU（Case Western Reserve University Bearing Data Center，凯斯西储大学轴承数据中心）：旋转机械故障诊断领域使用最广泛的公开轴承数据集，包含不同故障类型和严重程度的振动信号，但采集条件为实验室环境。}/Paderborn\footnote{Paderborn轴承数据集：德国帕德博恩大学发布的轴承故障诊断数据集，包含多种轴承型号和更接近工业现场的运行条件，常作为CWRU数据集的补充验证数据。}等公开数据集）+ 变工况迁移实验（不同转速/负载下的泛化测试）。排除：仅仿真数据无实验验证、仅使用单一工况数据的研究
\smallskip

三份拆解放在一起，框架的通用性一目了然：无论是社区居民、细胞系还是滚动轴承，PICOS都能将模糊的研究兴趣拆解为五个可操作的检索维度。但具体到不同学科，框架的适配度存在差异。

\subsubsection*{\textbf{•一个提醒：PICOS框架在不同学科中的适用性存在差异}}在临床医学和流行病学中，PICOS是成熟的检索方法论，拆解后的每个组件都直接对应PubMed的专业检索字段。在环境健康（小林）和分子生物学（张博）中，套用PICOS需要做一些调整：``研究对象''变为``细胞系/组织样本''，但框架的整体逻辑仍然可用。

在工程和计算机科学（王工）中，PICOS的适配度更低：``研究对象''是轴承/数据集，``干预''是算法模型，``对照''是baseline方法，这些组件在技术文献数据库（如IEEE Xplore）中没有对应的受控词表，检索效果取决于研究者所用的数据库是否支持自然语言检索。

如果PICOS在研究者的学科中适用性不佳，可以考虑使用替代框架：工程领域常用的``任务-方法-数据-指标''四要素拆解，或者计算机科学中``问题定义→现有方法→本文方法→评估方式''的结构化检索。关键在于找到一个能够帮助研究者将模糊想法拆解为边界清晰的问题的框架，框架的名称并非关键。



% --- 批量初筛：AI驱动的三级分类 ---
\subsection{文献的批量初筛与分类}

检索式调整完毕后，张博执行了检索，获得500余条结果。其中一半以上属于``关键词匹配但实际内容不相关''，逐条审阅全部标题与摘要的工作量极为繁重。这类重复性筛选工作正是AI的适用场景。

\subsubsection*{\textbf{•三级分类法的操作要点}}批量文献筛选的核心在于将判断标准表述清楚。提示词中需定义三个级别：\textbf{高度相关}（核心研究问题直接匹配，需要精读全文）、\textbf{中度相关}（部分变量或方法相关，需要进一步查看正文后再决定）、\textbf{不相关}（主题不匹配）。

关键在于为每个级别配备\textbf{可操作的判断标准}，即不是``主题相关''这类模糊描述，而是``是否同时报告了A指标和B指标''``是否包含某种实验类型''等可从摘要中直接判断的条件。

同时必须附加一条硬性限制：\textbf{仅基于提供的摘要内容判断，不要推测摘要中未提到的信息；信息不足以判断时分类为「中度相关」。}这条限制是为了防止AI因识别到熟悉的术语组合而做出过度关联，将不相关的论文错误纳入精读范围。

以下通过张博的案例展示具体操作方法。他把500多篇的标题摘要打包给了AI。

\begin{userbox}
以下是我在一次文献检索中获得的50篇论文的标题与摘要（检索式见前文PICOS拆解部分）。我的研究问题是：circRNA通过miRNA海绵机制调控胃癌细胞增殖和侵袭的功能与分子机制。

请将每篇论文按以下三级分类：
-「高度相关」：研究报告了circRNA在胃癌中的功能实验（增殖/侵袭/迁移表型）且涉及miRNA海绵机制
-「中度相关」：在胃癌中研究了circRNA但未深入机制，或研究了circRNA-miRNA互作但限于其他癌种
-「不相关」：虽然摘要中包含circRNA和gastric cancer，但核心研究内容为纯表达谱分析、生物信息学预测无实验验证、或主要研究其他癌种

每篇论文给出：分类结果 + 一句判断依据（研究者依据摘要中的什么信息做出了这个分类）。
重要限制：仅基于我提供的摘要内容进行判断，不要推测摘要中没有提到的信息。如果摘要信息不足以做出判断，分类为「中度相关」。
\end{userbox}

\begin{aibox}
\begin{formalTable}[H]
\begin{tabular}{p{2.5cm} p{4.0cm} p{4.0cm} p{3.0cm}}
\toprule
\textbf{编号} & \textbf{论文信息} & \textbf{分类} & \textbf{判断依据} \\
\midrule
1 & ``CircRNA hsa\_circ\_0001234 promotes gastric cancer cell proliferation and invasion via sponging miR-xxx to upregulate YYY expression'' & 高度相关 & 胃癌+功能实验（增殖/侵袭）+miRNA海绵机制，三个维度全部匹配 \\
2 & ``Expression profile and bioinformatics analysis of circRNAs in gastric cancer tissues'' & 不相关 & 纯表达谱+生物信息学预测（GO/KEGG富集），摘要中无任何功能实验验证 \\
3 & ``CircHIPK3 acts as a miR-XXX sponge to promote glioblastoma progression'' & 中度相关 & miRNA海绵机制明确但因研究对象为胶质母细胞瘤而非胃癌。若其机制通路与胃癌共享可考虑作为引文参考 \\
4 & ``CircRNA\_001569 knockdown inhibits gastric cancer growth in vitro and in vivo'' & 中度相关 & 胃癌+功能表型明确，但摘要仅描述表型而未提及miRNA或其他分子机制，需查看正文机制部分 \\
5 & ``The role of non-coding RNAs in drug resistance of gastric cancer: a systematic review'' & 不相关 & 综述文章，非原创研究。可作为背景引文来源，但不属于本次筛选目标（原创机制研究） \\
\bottomrule
\end{tabular}
\end{formalTable}

（仅展示前5篇作为格式示例，完整输出为50篇的分类表）
\end{aibox}

\begin{critiquebox}
张博并未直接采信这份分类表。他从AI标记为``高度相关''的论文中随机抽取了10篇，前往PubMed核对全文。其中8篇确实高度相关，1篇实际为结直肠癌与胃癌的联合分析（AI因摘要中出现``胃癌''一词即将其归入高度相关，查看正文后方知胃癌样本仅占一小部分），另有1篇的miRNA海绵证据仅来自荧光素酶报告实验，未进行功能回复实验（证据强度不足）。他又从``不相关''中抽检了5篇，其中1篇正文中实际上进行了circRNA功能实验，但摘要未使用``proliferation''、``invasion''等词汇，AI在摘要层面出现了漏判。

总体而言，AI的三级分类作为初筛辅助工具具有较好的效果，但不应以之替代研究者做最终决定。高度相关类别中可能混入方法学不够严谨的论文，不相关类别中也可能遗漏摘要信息不完整的论文。抽取若干篇进行人工核对即可有效控制这一风险。
\end{critiquebox}

\smallskip
同样的三级分类操作，在小林和王工的研究中形成了截然不同的分类标准，学科差异在分类维度上体现得最为明显。

小林（环境健康）的分类维度是绿地暴露指标×抑郁症状指标×研究设计。她对300余条Web of Science检索结果运行了分类，在高度相关中抽检了10篇：8篇确实相关，但2篇AI看到摘要中``green space''和``depression''同时出现即归入高度相关，实际研究内容却是绿地与体力活动的关系。从``不相关''中另发现1篇漏检，正文确实研究了绿地与抑郁，但摘要仅使用了``built environment and health''这类笼统表述。

王工（机械故障诊断）的分类维度是深度学习方法×故障类型×变工况实验。他的400余条IEEE Xplore结果运行完毕后，发现AI对``故障诊断''（fault diagnosis，分类问题）与``状态监测''（condition monitoring）、``剩余寿命预测''\footnote{RUL（Remaining Useful Life，剩余使用寿命）：预测机械设备或部件在给定运行条件下还能正常工作的时间长度，属于回归问题，与故障诊断（分类问题）在任务定义上有本质区别。}的区分能力不足：AI识别到vibration signal和bearing等熟悉的术语组合即做出过度关联，但fault diagnosis与RUL prediction在任务定义上完全不同。如果研究者自身不熟悉子领域的任务边界，可能会将不相关的论文错误纳入精读范围。

三人的经验指向同一个结论：AI的分类作为初筛辅助工具效果良好，但不能让它替代研究者决定哪些论文精读、哪些排除。正确的做法是：根据文献总量确定抽检比例（如总数超过500篇，建议高度相关和不相关各抽检不少于20篇），从AI标记为``高度相关''和``不相关''的论文中随机抽取核对原文，确认分类边界不存在系统性偏差后，再参考其余的分类结果。若抽检中发现明显误判，应调整分类标准并重新筛选，不能直接批量信任。


% --- 引文扩展：顺藤摸瓜找关键文献 ---
\subsection{基于引文网络的文献扩展}

张博筛选出20余篇高度相关论文，这些是他接下来需要精读的核心文献。但在开始逐篇精读之前，还有一个值得关注的步骤：这些论文的参考文献中可能包含更具基础性价值的文献。被多篇高相关论文共同引用的文献，通常是该领域的奠基性研究；反过来追踪哪些文献引用了这些高相关论文，则可能发现在其基础上进一步发展的近期研究。两个方向分别检索，有助于减少关键文献的遗漏。

AI在此环节的价值不在于替代数据库的引文查询功能（数据库在这一任务上更为精确），而在于帮助研究者从大量被引文献中识别出应优先阅读的篇目。

\begin{userbox}
以下是我筛选出的20篇高度相关论文的参考文献列表（已合并去重，共287条）。请基于参考文献的标题和发表信息，帮我识别：
1. 被3篇及以上高相关论文共同引用的文献（这些很可能是该领域的奠基性研究）
2. 发表年份在最近3年内、被至少2篇论文引用的文献（这些可能是近年受到关注的新发现）
3. 如果研究者无法确定某篇文献的具体作者或年份，标注``[待核实]''
\end{userbox}

从500余条记录到35篇核心文献，张博依次完成了初筛、分类与引文扩展三个步骤。精读35篇论文大约需要一至两周，这一节奏对于进入一个新领域而言较为合理。

% --- 结构化精读：六维度提取阅读导航图 ---
\subsection{论文的六维度结构化精读}

阅读论文的效率低，原因通常不在于阅读速度本身，而在于读完一段后需要停下来思考：这段内容的含义是什么？与自身研究有何关联？这种在阅读与思考之间的来回切换最为消耗精力。在逐段精读之前，先让AI提取论文的结构化信息，相当于在正式阅读前获得一张导航图。它帮助研究者提前标注了重点信息（效应量的规模、方法的可靠性、局限所在、与自身研究的关联程度），研究者带着这些标注去阅读，可以将注意力集中在最关键的内容上。

\subsubsection*{\textbf{•六维度的设计逻辑}}六个维度的选取各有其针对性，各自解决阅读过程中的一个具体问题：

\textbf{研究问题}：帮助研究者判断该论文是否值得阅读（研究问题与自身研究是否相关？）

\textbf{研究设计与数据}：帮助研究者评估方法的可信度（样本量是否充足？实验体系是否完整？）

\textbf{核心变量与测量}：帮助研究者判断操作化定义的合理性（变量是如何测量的？）以及自身能否复用这套测量方法

\textbf{主要发现}：帮助研究者快速定位核心结论，引用原文数据可防止AI转述时出现数值转录错误或措辞从审慎变为确定

\textbf{关键局限}：即作者自己承认的问题，研究者在引言或讨论中可直接引用，作为自身研究切入点的论据

\textbf{与自身研究的相关性}：帮助研究者进行优先级排序，以1-5分评分决定该论文是精读还是浏览

提示词中必须附加一条关键限制：\textbf{找不到某个维度信息时标注``[论文未报告]''，不要推测。}这能防止AI``补充''出看似合理但论文中并不存在的信息。

以下通过张博的案例展示六维度提取的具体用法。注意：本案例中出现的论文标题、细胞系名称和实验数据为教学模拟示例，用于演示结构化提取的操作方法，不对应真实发表的论文。

\begin{userbox}
我已上传一篇论文。请帮我按以下维度进行结构化提取，每个维度的输出不超过5行：

1. 研究问题：该论文要回答的核心问题（用一句话概括）
2. 研究设计与数据：实验类型（体外/体内/临床样本）、细胞系/动物模型/样本量、数据来源
3. 核心变量与测量：circRNA的鉴定方法、miRNA靶标的验证方法（荧光素酶报告实验/RNA pull-down等）、功能表型的测量方法
4. 主要发现：核心表型效应和分子机制，引用原文关键数据
5. 关键局限：作者自己报告的局限有哪些（不超过3条）
6. 与我的研究的相关性：用1-5分评分，并给出理由

重要：如果在论文中找不到某个维度的信息，标注``[论文未报告]''而不是推测。
\end{userbox}

\begin{aibox}
\textbf{论文：}CircRNA hsa\_circ\_0008039 promotes gastric cancer proliferation via sponging miR-375 to upregulate YAP1 expression

\begin{enumerate}[nosep,leftmargin=*]
\item \textbf{研究问题：}hsa\_circ\_0008039是否通过吸附miR-375来解除miR-375对YAP1的抑制，从而促进胃癌细胞增殖？

\item \textbf{研究设计与数据：}
  \begin{itemize}[nosep]
  \item 体外：AGS、MGC-803胃癌细胞系，GES-1正常胃上皮细胞
  \item 体内：BALB/c裸鼠皮下成瘤模型（n=6/组）
  \item 临床：40对胃癌-癌旁组织（qPCR验证表达差异）
  \end{itemize}

\item \textbf{核心变量与测量：}
  \begin{itemize}[nosep]
  \item circRNA鉴定：RNase R消化+放线菌素D稳定性实验+Sanger测序验证环化位点
  \item miRNA靶标验证：荧光素酶报告实验（野生型/突变型）+ RNA pull-down + FISH（荧光原位杂交）共定位
  \item 功能表型：CCK-8增殖曲线、EdU染色、集落形成、流式细胞术细胞周期分析
  \end{itemize}

\item \textbf{主要发现：}
  \begin{itemize}[nosep]
  \item circ\_0008039在胃癌组织中显著上调（fold change=3.2, p<0.001 vs 癌旁）
  \item circ\_0008039 knockdown后AGS细胞增殖能力下降约45\%（CCK-8 OD值, p<0.01）
  \item 荧光素酶实验证实miR-375直接结合circ\_0008039（野生型组荧光强度\(\downarrow\)58\%, p<0.001）
  \item YAP1蛋白水平与circ\_0008039正相关（r=0.62, p<0.001），miR-375 mimic可逆转此相关性
  \end{itemize}

\item \textbf{关键局限：}
  \begin{itemize}[nosep]
  \item 临床样本量偏小（40对），未分析circ\_0008039表达与TNM分期的关联
  \item 未检测circ\_0008039是否还通过其他miRNA发挥作用（未做系统筛选）
  \item 裸鼠模型仅测量了肿瘤体积，未分析转移灶
  \end{itemize}

\item \textbf{与我的研究相关性：}4/5。circRNA→miRNA→靶基因的研究框架可作为方法学模板。局限（未做系统miRNA筛选、临床样本量偏小）为我的研究设计提供了改进方向。
\end{enumerate}
\end{aibox}

\smallskip
\subsubsection*{\textbf{•六维度在不同学科中的使用侧重}}同一套六维度框架，不同学科在使用中作用最大的维度各有侧重。

小林（环境健康）使用结构化提取处理一篇欧洲四城市公园可达性与老年人抑郁的纵向研究时，\textbf{关键局限}一栏最具参考价值：作者自己写明``结果可能不适用于亚洲高密度城市''，这恰好成为她可以在引言中直接引用的论据：``已有纵向证据来自欧洲，尚无亚洲高密度城市的研究''。此外，\textbf{核心变量与测量}中标注的``[论文未报告]中介/调节变量''也提醒她，自己的研究在这一点上恰好具有增量，她将心理恢复感作为中介变量纳入了研究设计。

王工（机械故障诊断）使用同一框架拆解一篇多尺度Transformer故障诊断论文时，\textbf{核心变量与测量}一栏信息密度最高：AI帮他提取了数据增强策略（高斯噪声+随机缩放）、完整超参数配置（d\_model=256、dropout=0.1）和消融实验结果。这些细节通常分散在方法部分的文字叙述中，有了这份结构化摘要，他在自己的实验中可直接对照参考，省去了在方法部分与附录之间反复检索的时间。

\subsubsection*{\textbf{•一个关键提醒：AI的结构化摘要不能替代原文}}AI在复述文献内容时容易产生三类错误：数值转录错误（效应量数值偏差）、语义强度偏移（原文中审慎的``may suggest''被改为确定的``demonstrates''）、前提条件遗漏（有条件的结论被表述为无条件结论）。凡计划引用的具体发现，必须从原文中亲自确认用词和数值。六维度摘要是阅读地图，不是引用来源。


% --- 方法对比：多篇论文方法学横向比较 ---
\subsection{多篇论文的方法学比较}

精读5-10篇高相关论文之后，研究者通常会遇到一个问题：这些论文使用的细胞系不同、检测方法各异、实验组合也存在差异，各项研究的结论之间难以直接比较。到了这个阶段，明确``哪项研究的方法更为可靠''比了解``哪项研究发现了什么''更具优先性。

\subsubsection*{\textbf{•方法学横向对比的核心价值}}不同论文报告的效果差异为何如此显著（例如跨工况准确率从60\%到85\%不等），方法学差异通常是主要原因，而非真实效果的差异。将每篇论文的实验体系、测量工具和统计方法进行系统性的横向比对，能帮助研究者回答三个关键问题：

哪些结论差异来源于方法差异（如使用了不同的量表、不同的``跨工况''操作化定义），而非真实效应差异；

哪些论文的方法学最为稳健，值得作为研究者自身实验设计的参照模板；

领域内是否存在系统性的方法学短板（如多数论文仅使用了一个数据集进行验证），研究者可在自身研究中加以弥补。

提示词中需定义3-5个对比维度，并明确要求AI``仅基于研究者提供的信息进行比较，不要填补论文中未报告的信息''。

以下通过张博的案例展示方法学横向对比的具体用法。注意：本案例中的论文编号、实验体系和对比数据为教学模拟示例，用于演示方法学比较的分析维度，不对应真实发表的论文。

\begin{userbox}
以下是我已精读的8篇关于circRNA在胃癌中作用机制的实证论文。请帮助我做一个方法学维度的横向比较：

比较维度：
1. 实验体系：使用的细胞系种类和数量、是否有体内实验（裸鼠模型类型和样本量）
2. circRNA鉴定：是否进行了RNase R消化验证、放线菌素D稳定性实验、Sanger测序
3. miRNA靶标验证：是否使用荧光素酶报告实验（是否包含突变型对照）、RNA pull-down、FISH
4. 功能表型覆盖：增殖/侵袭迁移/凋亡/细胞周期，分别使用了什么测量方法
5. 表型回复实验（rescue experiment）：是否通过共转染（circRNA抑制剂+miRNA inhibitor）验证了circRNA的效果确实通过该miRNA介导

输出为一个横向对比表格，每篇论文一列。表格之后，请选出在方法学上最为稳健的2篇论文，并为每篇给出2-3条选择理由。

注意：仅基于我提供的信息进行比较，不要填补论文中未报告的信息。
\end{userbox}

\begin{aibox}
\textbf{方法学横向对比表（节选前4篇）：}

\begin{formalTable}[H]
\small
\begin{tabular}{p{2.8cm} p{2.5cm} p{2.5cm} p{2.5cm} p{2.5cm}}
\toprule
\textbf{对比维度} & \textbf{论文1} & \textbf{论文2} & \textbf{论文3} & \textbf{论文4} \\
\midrule
细胞系数量 & 2种（AGS, MGC-803） & 1种（BGC-823） & 3种（AGS, MGC-803, BGC-823） & 2种（SGC-7901, AGS） \\
体内实验 & 裸鼠成瘤（n=6） & 无 & 裸鼠成瘤（n=8）+ 转移模型 & 裸鼠成瘤（n=5） \\
circRNA鉴定 & RNase R + 放线菌素D + Sanger & RNase R仅 & RNase R + Sanger & RNase R + 放线菌素D \\
miRNA验证 & 荧光素酶（含突变型）+ RNA pull-down + FISH & 荧光素酶（仅野生型） & 荧光素酶（含突变型）+ RNA pull-down & 荧光素酶（含突变型） \\
功能表型覆盖 & 增殖+侵袭+迁移+凋亡 & 仅增殖 & 增殖+侵袭+迁移+凋亡+细胞周期 & 增殖+侵袭 \\
回复实验 & 有（共转染circ siRNA + miR inhibitor） & 无 & 有（共转染+双荧光素酶回复） & 无 \\
\bottomrule
\end{tabular}
\end{formalTable}

\smallskip
\textbf{方法学最稳健的2篇论文：}

\textbf{论文3}（3种细胞系 + 体内成瘤及转移模型 + 完整circRNA鉴定流程 + 突变型对照 + 五类表型全覆盖 + 双回复实验）。选择理由：(1) 实验体系最完整——三种细胞系覆盖了胃癌细胞系的常见亚型，体内实验同时包含皮下成瘤和转移模型；(2) miRNA靶标验证达到了该领域的金标准——荧光素酶报告实验含突变型对照、RNA pull-down独立验证、FISH共定位，三者交叉印证。

\textbf{论文1}（2种细胞系 + 裸鼠成瘤 + 完整circRNA鉴定 + 含突变型对照的miRNA验证 + 回复实验）。选择理由：(1) 虽然细胞系数量少于论文3，但其circRNA鉴定和miRNA验证的完整度高于其余6篇；(2) 共转染回复实验直接验证了circRNA的效果确实通过该miRNA介导，这一实验设计在8篇论文中仅2篇完成。

注意：论文3和论文1使用了不同的细胞系（AGS/MGC-803/BGC-823 vs AGS/MGC-803），两者的结果不可直接比较。例如论文3报告的增殖抑制率更高，可能是因为BGC-823细胞系中该circRNA的本底表达更高，而非其方法本身更有效。
\end{aibox}

\begin{critiquebox}
张博拿到表格后，注意到一个此前忽略的维度：\textbf{回复实验}。8篇论文中仅2篇完成了共转染回复实验（circRNA抑制剂+miRNA inhibitor），其余6篇仅通过荧光素酶报告实验间接推断miRNA海绵机制。没有回复实验意味着无法排除一种可能性——circRNA敲低后观察到的表型变化，可能是通过研究者未检测的其他miRNA或其他非miRNA机制（如RBP互作）介导的。回复实验的缺失不是一个小问题，它直接影响了这部分论文关于``circRNA通过miRNA海绵机制调控胃癌增殖''这一结论的证据强度。

张博在自己的实验方案中因此增加了一项：在circRNA敲低的基础上共转染miRNA inhibitor，验证表型是否被回复。如果没有这张对比表，他很可能沿用大多数论文的做法——荧光素酶报告实验后直接下结论，而不会意识到回复实验的缺失是已有文献中一个系统性的方法学薄弱点。
\end{critiquebox}

这份对比的产出不是告诉研究者``谁的结果最可信''（那是研究者自己判断的事），而是让研究者看清楚方法学上的差异。例如，为什么A论文里circRNA敲低后增殖抑制了30\%，B论文里抑制了60\%？可能是因为使用的细胞系不同，或者增殖测量的方法不同。将这些方法学差异系统地呈现出来，研究者在综合解读已有证据的时候就能做出更准确的判断。

\smallskip
\subsubsection*{\textbf{•方法学对比在不同学科中的典型发现}}同一套横向对比逻辑，在不同学科中暴露出的方法学问题各不相同。

小林（环境健康）对比了8篇绿地-抑郁关联研究后，发现了一个容易被忽略的测量学问题：不同量表测量的构念并不相同。PHQ-9更侧重睡眠、食欲等躯体化症状，CES-D更侧重情绪体验，使用不同量表的研究所报告的绿地-抑郁效应量范围差异明显（PHQ-9组β=-0.12至-0.22，CES-D组β=-0.05至-0.14）。这种差异可能来源于量表本身的测量特性差异，而非绿地效应的真实差异。

若不厘清测量工具的差异而直接比较``绿地-抑郁效应量''，无异于将不同性质的结果混为一谈。她在自己的研究中沿用了PHQ-9，但在讨论部分特别说明：本研究的效应量与使用CES-D的研究不可直接比较。

王工（机械故障诊断）在对比10篇故障诊断论文时，则发现了一个更为根本的问题：不同论文报告的``跨工况准确率''跨度很大（部分论文低至60\%左右，高者可接近90\%），主要原因并非模型性能的实际差异如此之大，而是\textbf{``跨工况''的操作性定义不一致}。有的论文仅跨了转速（如1797→1772 rpm），有的仅跨了负载（如0→3 HP），还有的同时跨了两者。跨转速迁移的难度通常低于跨负载迁移，因为转速变化主要影响信号的时间尺度，而负载变化则会影响信号的幅值和冲击特性。

这一发现直接影响了王工的实验设计：他在自己的方案中明确写明需同时进行跨转速和跨负载两个方向的测试，并增加Paderborn作为第二验证数据集，以避免审稿人质疑``仅在一个数据集的一个迁移方向上进行了测试''。


% --- 三库构建：论文库、数据库、代码库 ---
\subsection{论文库、数据库与代码库}

前面几节讨论了如何检索、如何筛选、如何阅读。但有一个问题尚未回答：文献阅读完成之后，研究成果如何系统化存储与调用？许多人读完几十篇论文后，笔记分散在PDF批注、桌面便签和各种收藏夹中，数周后便难以检索。检索→筛选→精读这条流程的终点，应当是一个可随时检索与调用的结构化知识库。具体来说，是三个库。

\textbf{论文库}是地基。论文库不是简单地将PDF文件存入文件夹，而是需要让AI为每篇论文赋予结构化标签。张博的做法是：每精读完一篇论文，将六维度结构化摘要（详见本节六维度结构化精读部分）存入Zotero\footnote{Zotero：开源的文献管理软件，支持自动抓取论文元数据、PDF管理和标签分类，可一键生成参考文献列表。}，同时让AI根据摘要内容自动生成3-5个关键词标签和一句不超过50字的核心里程碑描述。数月之后，他的circRNA论文库积累了35篇论文，每篇均附有标签和一句话摘要。撰写引言时，若需查找``circRNA在胃癌中作为miRNA海绵的证据''，按标签筛选即可调出十余篇相关论文，无需重新检索PubMed。王工的论文库更偏工程，他按``数据集''``模型架构''``工况条件''三个维度打标签，每篇论文的代码链接和超参数配置也一并存入，在后续实验中可直接作为参考。

\textbf{数据库}管理的是实验和调查产生的一手数据。数据库的核心价值不在于``存储''，而在于存储时即附带完整的元数据描述。小林的那份问卷调查（n=680），AI帮她生成了一份数据字典，包含每个变量的名称、含义、取值范围、缺失值编码和数据来源。这份字典与清洗后的数据文件存放在同一目录下，半年后再次打开，无需逐一回忆每个变量的具体含义。王工的振动信号数据同理，AI帮他将每个数据文件的采样频率、工况参数和轴承型号写入一个README文件，撰写论文方法部分时可直接引用。

\textbf{代码库}是三个库中最容易被忽略的。硕士生的代码往往在实验完成后即被搁置，分散在多个Jupyter Notebook\footnote{Jupyter Notebook：一种交互式编程环境，支持在浏览器中编写和运行Python（及R、Julia等）代码，可混合代码、文本和可视化输出，广泛应用于数据科学和机器学习领域。}中，缺少注释和版本记录。当审稿人要求复现研究结果时，研究者自己可能已无法重新运行这些代码。AI在此能做两件事：一是补充注释，将一段没有注释的Python或R代码交给AI，让它逐行解释并生成规范的docstring\footnote{docstring（Documentation String，文档字符串）：Python代码中用于描述函数、类或模块功能的字符串，通常放在定义体的首行，是三引号包裹的多行文本。}；二是撰写实验日志，每次运行完一组实验，将参数配置、结果摘要和代码commit\footnote{commit：Git版本控制系统中记录代码修改快照的基本单位，每次commit包含修改内容的描述和唯一标识符（commit hash）。}号记录下来，AI帮助格式化为Markdown\footnote{Markdown：一种轻量级标记语言，使用纯文本格式化语法，支持标题、列表、表格和代码块等元素，可快速转换为HTML或PDF。}表格。张博的习惯是每完成一轮细胞实验，就让AI生成一段实验记录，存入代码库，撰写论文方法部分时可直接提取使用。

三个库的搭建贯穿整个研究过程，不需要一次性完成：读完一篇存入一篇，做完一轮实验记录一轮。建库本身不需要AI参与决策（标签如何设定、变量如何命名，最终由研究者自己确定），但AI能够帮助研究者将重复性的整理工作压缩至最低水平。

\smallskip
\begin{userbox}
我已经精读了35篇circRNA与胃癌相关的论文，每篇都有六维度结构化摘要。请帮我为每篇论文生成：(1) 3-5个关键词标签（从circRNA名称、机制类型、表型、临床意义四个维度各选至少一个）；(2) 一句不超过50字的核心里程碑描述。输出为一个Markdown表格，列名为：论文编号、标签、里程碑描述。
\end{userbox}

\begin{aibox}
AI输出了一个35行的Markdown表格。每篇论文的标签来自四个维度（circRNA名称、机制类型、表型、临床意义），里程碑描述控制在50字以内。张博把这个表格导入Zotero作为自定义字段，以后写引言和讨论时可以直接按标签筛选。
\end{aibox}

\begin{critiquebox}
标签体系应由研究者自己确定，因为AI不知道研究者接下来的研究方向是什么，它只能基于论文内容提供建议。张博根据自己的课题方向（circRNA编码短肽的免疫调节）对AI生成的标签做了调整，将``免疫微环境''和``翻译短肽''两个维度纳入其中。论文库的价值不在于AI节省了多少录入时间，而在于研究者将来能按照自己关心的维度快速检索到所需文献。
\end{critiquebox}

\begin{figure}[H]
\centering
\BookFigureImage{images_chapter_05/三库关系与调用图.png}
\caption{三库关系与调用图（图中``知识库、证据库、任务库''分别对应正文中的``论文库、数据库、代码库''）}
\label{fig:ch5_04}
\end{figure}
\par

检索→精读→建库，三库之间的关系如图\ref{fig:ch5_04}所示。至此，研究者在论文库中拥有了打好标签的核心文献，在数据库中拥有了清洗好的数据和完整的数据字典，在代码库中拥有了带注释和实验日志的代码。


% --- 提示词模板 ---
\subsection*{提示词模板}

\textbf{模板2-1：PICOS框架拆解与检索式生成}

\begin{tcolorbox}[promptbox]
我的研究问题是：[完整研究问题，限100字]。

请用PICOS框架拆解：
- P（研究对象）：人群/样本/细胞系/动物模型等
- I/E（暴露/干预）：核心自变量或干预条件
- C（对照）：对照条件
- O（结局）：结局指标
- S（研究设计）：设计类型限定

然后基于拆解结果为 [目标数据库] 生成检索式，标注``高灵敏度版''和``高特异度版''。
\end{tcolorbox}

\textbf{模板2-2：批量文献相关性分类}

\begin{tcolorbox}[promptbox]
以下是 [N] 篇论文的标题与摘要。我的研究问题是：[完整研究问题]。

请将每篇论文按三级分类：
-「高度相关」：核心研究问题直接相关，需要精读全文
-「中度相关」：部分变量或方法相关，需要进一步查看正文后再决定
-「不相关」：主题不匹配

每篇论文输出：分类结果 + 一句判断依据。
要求：仅基于我提供的摘要内容判断，不要推测摘要中未提到的信息。信息不足以判断时分类为「中度相关」。
\end{tcolorbox}

\textbf{模板2-3：单篇论文结构化精读}

\begin{tcolorbox}[promptbox]
我已上传一篇论文。请用以下维度进行结构化提取（每维度不超过5行）：

1. 研究问题：核心问题（一句话概括）
2. 研究设计与数据：类型、样本量/样本类型、数据来源
3. 核心变量与测量：自变量/因变量/关键实验的测量方式
4. 主要发现：核心效应量和统计指标（引用原文数值）
5. 关键局限：作者自己报告的局限（不超过3条）
6. 与我的研究相关性：1-5分评分 + 理由

如果在论文中找不到某个维度的信息，标注``[论文未报告]''。
\end{tcolorbox}

\textbf{模板2-4：多篇论文方法学横向对比}

\begin{tcolorbox}[promptbox]
以下是我已精读的 [N] 篇关于 [主题] 的论文。请做方法学维度横向比较：

对比维度：[根据学科自定义，如样本量/测量工具/实验条件/统计方法等]

输出为表格（每篇一列），附上研究者的方法学质量排序和理由。
\end{tcolorbox}

\textbf{模板2-5：构建个人论文库}

\begin{tcolorbox}[promptbox]
我已精读 [N] 篇关于 [主题] 的论文，每篇都有结构化摘要。
请帮我为每篇论文生成：
1. [N] 个关键词标签（从 [维度1]、[维度2]、[维度3] 等维度各选至少一个）
2. 一句不超过50字的核心里程碑描述
输出为Markdown表格，列名为：论文编号、标签、里程碑描述。
\end{tcolorbox}


% --- 常见错误与规避 ---

\begin{mistakebox}

\textbf{把AI生成的检索式直接拿去正式检索。}AI可能遗漏领域习惯用语（如circRNA领域常用的``sponge''）或使用过于宽泛的通配符（如proliferat*），检索式应在目标数据库中逐组件测试、确认代表性论文被覆盖后再使用。

\smallskip

\textbf{只看AI对摘要的分类就决定精读还是排除。}高度相关中可能混入方法学不严谨的论文，不相关中也可能遗漏摘要信息不完整的论文，必须从两级各随机抽检一定数量，确认分类边界无系统性偏差后，再批量信任其余分类结果。

\smallskip

\textbf{不读原文，直接拿AI生成的结构化摘要当引用来源。}AI复述文献时可能产生数值转录错误、语义强度偏移（``may suggest''变为``demonstrates''）或前提条件遗漏，凡计划引用的具体发现必须从原文亲自确认用词和数值。

\end{mistakebox}

\subsection*{小结}
本节从语义检索与关键词检索的技术差异入手，展示了PICOS框架如何将模糊研究兴趣转化为可执行的检索策略。三级分类法和六维度结构化精读帮助研究者从海量文献中高效锁定核心篇目，方法学横向对比防止将不同质量的研究同等对待。检索→精读→建库的完整流程即可形成闭环。

\bigskip
\subsection*{课后习题}
\begin{enumerate}
\setlength{\itemsep}{6pt}

\item 语义检索与关键词检索各有什么优劣势？分别在什么场景下应优先使用？

\item 基于你的研究问题完成PICOS拆解，生成检索式并在数据库中验证。PICOS框架在你的学科中适用性如何？如果不完全适用，你做了哪些调整？

\item 六维度结构化提取法包含哪六个维度？为什么AI的结构化摘要不能替代原文阅读？

\item 在三级分类初筛后，为什么必须从``高度相关''和``不相关''中各随机抽检一定数量的论文？如果不做这一步骤，可能出现什么问题？

\end{enumerate}
\newpage

% ============================================================
% 5.3 AI辅助课题设计与项目申报
% ============================================================
\section{AI辅助课题设计与项目申报}
\label{sec:5.3}

\begin{sectionkeypoints}
\item 理解AI在课题设计中的能力边界——它做的是模式匹配而非因果推理，生成的假设是启发而非结论
\item 用「已知-争议-空白」三角框架识别研究空白，从创新性、风险、可操作性、竞争程度四个维度评估其价值
\item 将研究空白转化为可检验的研究假设，每条假设均附带可证伪条件（明确写出什么结果支持、什么结果否定）
\item 在给定资源约束下设计可行的研究方案，包括实验框架、基线选择、评估指标和统计分析计划
\item 用四步检验法在分析数据前锁定分析方法、检查数据质量、确认方法时效性，避免p-hacking和方法保守偏向
\end{sectionkeypoints}

\subsubsection*{\textbf{•模式匹配与因果推理的边界}}

选题是三个范式交汇最密集的地方。经验驱动让研究者根据已有知识判断方向，数据驱动让研究者从文献数据中寻找空白，AI驱动则帮助研究者将这一过程系统化。但最终的判断——哪个方向值得做、哪个方案能够落地，仍然属于研究者。

到了选题这一步，研究者最希望AI做的事情是分析：哪里还没有人做过？我能提出什么假设？这个方案是否可行？

然而，AI执行的分析与研究者的分析在性质上完全不同。前者擅长的是模式匹配，即识别``这个领域已有的论文通常如何推演''。王工把30篇旋转机械故障诊断的论文摘要输入给AI后，AI能从中挑出反复出现的推论套路（``变工况下泛化能力不足''``迁移学习是热门方向''``图神经网络有潜力''），然后把它们重新组合成一个看起来很像学术假设的东西。成千上万的论文在讨论部分都用类似的句式提出未来研究建议，AI只是学会了这个语言模板。

真正的因果推理，AI无法完成。模式匹配回答的是``已有论文通常怎么推测''，因果推理回答的是``变量之间实际存在什么因果关系''。AI能生成格式正确的假设，但它不知道这个假设在真实世界中是否成立，不知道Transformer的自注意力机制是否确实比CNN的局部感受野更能捕捉故障信号的全局依赖。它只知道在训练数据里``Transformer在故障诊断中的应用''是一个增长中的讨论方向。

\subsubsection*{\textbf{•AI输出的正确定位：启发而非结论}}AI帮助研究者系统梳理，确保没有遗漏明显的方向和变量。但哪个方向值得投入时间、哪个方案在研究者自身的约束下确实可行，这些判断仍需研究者自己完成（图\ref{fig:ch5_05}）。

\begin{figure}[H]
\centering
\BookFigureImage{images_chapter_05/LLM模式匹配vs因果推理边界示意图.png}
\caption{LLM模式匹配与因果推理边界示意图}
\label{fig:ch5_05}
\end{figure}
\par

% --- 空白识别：从文献缝隙中找创新点 ---
\subsection{研究空白的系统识别}

王工读完30篇核心论文，积累了厚厚一叠笔记，但能够明确表述出来的只有一句模糊的印象：``大家都在单一工况下训练模型，一旦更换工况，模型性能就显著下降。''仅凭这种程度的模糊感觉，远不足以构成研究课题——而这恰恰是大部分硕士生开题前的真实状态。

\subsubsection*{\textbf{•「已知-争议-空白」三角框架}}从模糊感觉到清晰的研究空白，核心方法是将几十篇论文摘要提供给AI，让它按三个维度分类：

①\textbf{已知篇}——已有研究一致发现了什么（每点附支持该结论的论文数量范围），写作时一笔带过即可；

②\textbf{争议篇}——哪些结论在不同研究中存在分歧、各方依据是什么，这些争议恰好是研究者可以介入的学术对话；

③\textbf{空白篇}——分两类处理：很多论文都提到的``未来方向''（大家都能看到的，可能已经有人在做了），和只有一两篇提过的方向（更值得多加关注）。

拿到AI的空白分析后，用四个维度评估每个空白的价值：\textbf{创新性}（有多新）、\textbf{风险}（做不出来的概率多高）、\textbf{可操作性}（在研究者的约束下能不能做）、\textbf{竞争程度}（有多少团队在做）。AI能帮研究者列出并初评这些空白，但最终选择哪个方向，需要综合研究者自己的学科经验、资源约束和风险承受能力做判断。

以下通过王工的案例展示该框架的具体应用。

\begin{userbox}
我计划研究``基于深度学习的旋转机械故障诊断在变工况条件下的泛化性能提升''。以下是该领域30篇核心论文的摘要（已在之前的对话中逐批上传）。

请帮我做一个结构化的研究空白分析：

1.「已知篇」：已有研究一致发现什么？至少列出3点。每点附上支持该结论的论文数量范围
2.「争议篇」：哪些结论在不同研究中存在分歧？各方的核心依据是什么？
3.「空白篇」：哪些问题在上述论文的摘要中被明确标记为``未来研究方向''或``尚未研究''？哪些问题是这些论文的局限部分暴露出来但摘要未明确指出的潜在空白？

空白篇里，如果某个方向几乎每篇论文的``未来研究方向''都提到，它可能已经有人在做了；如果只有一两篇提过，倒更值得多加关注。
\end{userbox}

\begin{aibox}
\textbf{已知篇}

\begin{enumerate}[nosep,leftmargin=*]
\item CNN系列模型在单一工况下的轴承故障分类准确率可达98\%以上（25/30篇论文报告了>95\%的准确率）
\item 当测试工况（转速/负载）与训练工况不同时，模型性能显著下降（18/30篇论文报告了跨工况准确率下降10--30\%）
\item 时频域特征（小波变换、短时傅里叶变换STFT时频图）在变工况下的鲁棒性优于单一原始时域信号（12/30篇论文比较了两种方法）
\end{enumerate}

\smallskip
\textbf{争议篇}

\begin{enumerate}[nosep,leftmargin=*]
\item \textbf{迁移学习是否是变工况问题的最优解？}

支持方（12篇）：域对抗网络（DANN）、最大均值差异（MMD）等域适应方法在公开数据集上取得了最先进的跨工况结果

质疑方（5篇）：迁移学习需要目标域的部分标注数据，工业场景中往往不具备这个条件。这些论文主张基于信号处理的特征不变性方法更实用

\item \textbf{Transformer vs CNN在故障诊断中的优劣？}

支持Transformer方（8篇）：自注意力机制捕捉长距离时序依赖，理论上更适合非平稳振动信号

支持CNN方（14篇）：在中等规模数据集上CNN收敛更快、训练更稳定、对超参数选择的鲁棒性相对更强
\end{enumerate}

\smallskip
\textbf{空白篇}

\textit{很多论文都提到的：}
\begin{itemize}[nosep]
\item 需要在大规模工业现场数据（而非实验室公开数据集）上验证方法的实用性（22/30篇论文提到）→ 评估：方向正确但需要工业合作，资源门槛高
\item 轻量化模型以适应边缘计算部署（15/30篇论文提到）→ 评估：工程价值大的方向但对算法创新的空间可能有限
\end{itemize}

\textit{只有一两篇提过的：}
\begin{itemize}[nosep]
\item \textbf{物理约束融入深度学习模型}——将轴承的运动学和动力学方程作为正则化项约束网络训练（仅2篇论文提及，均为2023年后发表的会议论文）

→ 评估：高价值空白。目前故障诊断领域的深度学习模型是纯数据驱动的（``黑箱''），融入物理知识可能同时提升泛化能力和可解释性。该方向文献量极少，竞争不激烈。但技术门槛较高——需要研究者同时熟悉旋转机械动力学和深度学习两个领域

\item \textbf{跨设备类型的故障诊断泛化}——一个在轴承上训练好的模型能否直接用于齿轮故障诊断？（仅1篇论文在讨论中提到此方向）

→ 评估：极早期方向。理论上有吸引力（``通用故障诊断模型''），但存在较大的技术风险——轴承和齿轮的故障物理机制不同，共享特征空间是否真的存在尚未可知
\end{itemize}
\end{aibox}

\begin{critiquebox}
王工审视这份空白分析后，对``物理约束融入深度学习''这个仅有2篇论文涉及的方向产生了兴趣。他在Web of Science中搜索了``physics-informed deep learning bearing fault diagnosis''，发现过去两年内仅有5篇相关论文，且都来自机械工程和计算机科学的交叉团队。这个方向文献少、起步晚，但理论逻辑自洽：轴承的故障特征频率在理论上由转速、滚子数量、接触角等参数决定，这些物理关系可以作为先验知识约束网络的学习过程，而非让网络从零开始从数据中学习一切。

但他也意识到这个方向的风险：需要同时掌握旋转机械动力学建模和深度学习框架，且如果预实验发现物理约束对泛化性能的提升不显著，整个课题的核心论证将缺乏支撑。这个风险需要在假设推导阶段进一步评估。
\end{critiquebox}

\smallskip
同样的方法，小林和张博各自识别出的空白完全不同。学科差异在空白分析这一步体现得最为明显。

小林（环境健康）的空白分析节选：

\begin{itemize}[nosep]
\item \textbf{街景绿视率 vs NDVI对抑郁症状的预测能力直接对比}（仅1篇论文在局限中提到）→ 高价值空白。街景图像技术近年趋于成熟，且存在明确的科学问题——``人眼看到的绿色（绿视率）是否比卫星看到的绿色（NDVI）更直接地影响心理健康？''
\end{itemize}

张博（circRNA与胃癌）的空白分析节选：

\begin{itemize}[nosep]
\item \textbf{circRNA编码的短肽在胃癌微环境中的免疫调节功能}（仅1篇2022年首篇胃癌相关报道）→ 极高风险+极高创新的早期方向。文献量<5篇，竞争不激烈但实验体系需要从零搭建
\end{itemize}

\smallskip
纵观三个学科，可以发现同样的「已知-争议-空白」框架下，环境健康领域识别出的是``绿视率vs NDVI直接对比''这种测量学空白，机械工程领域识别出的是``物理约束融入深度学习''这种方法学空白，肿瘤生物学领域识别出的是``circRNA编码短肽''这种极高风险的早期方向。无论学科，空白分析的关键不是研究者``看到了什么空白''，而是研究者``如何评估每个空白的可行性''。创新性、风险、可操作性、竞争程度四个维度缺一不可。AI提供候选清单，研究者做出最终选择。

% --- 假设推导：把空白变成可证伪的命题 ---
\subsection{将研究空白转化为可检验假设}

空白是``没人做过什么''，假设是``如果我做了，会看到什么''。这一转化的关键不在于措辞，而在于研究者必须把``我想探索的方向''变成一句可以被数据证伪的陈述：\textbf{写得出否定条件，才是可检验的假设}。

\subsubsection*{\textbf{•可检验假设的三要素}}AI推导的假设格式规范、措辞学术，但它只是在模仿这个领域最常用的推论句式。把AI生成的假设转化为研究者自己的研究假设，需要逐条核查三个要素：

\textbf{可证伪性}——明确写出``什么实证结果会支持该假设''和``什么实证结果会否定该假设''。很多人在推导假设时只想着``我要证明它是对的''，忽略了``什么结果说明我错了''。能写出否定条件的假设，才是真正可检验的假设。
\textbf{已有间接支持}——每条假设至少引用1条已有论文的间接发现作为支撑。
\textbf{可操作性}——在研究者的实际约束下（时间、设备、数据、预算）是否可以检验。

此外还需要注意：AI给出的精确数字阈值（如王工H1中``准确率提升不到3个百分点''的否定条件、小林H1中β差值>0.05的判断标准）属于统计推断而非方法学标准，需要回到方法学文献中确认依据。

\begin{funbox}
\textbf{可证伪性：从波普尔到AI时代。}正文反复强调，``写得出否定条件，才是可检验的假设''。这一思想的源头可以追溯到20世纪最伟大的科学哲学家之一——卡尔·波普尔。波普尔在1934年出版的《科学发现的逻辑》中提出了一个颠覆性的观点：科学命题与伪科学的真正区别，不在于它能否被``证实''，而在于它能否被``证伪''。一个理论如果能被想象出的任何观察结果所支持而永远不会被推翻，那它就不是科学——例如弗洛伊德的精神分析，它可以解释任何人行为的任何可能走向，因而波普尔认为它不属于科学范畴。

波普尔的洞见时至今日仍然深刻影响着科研实践。正文中王工为每个假设明确写出``支持条件''和``否定条件''，正是可证伪性的具体操作化——无论数据最终支持还是否定假设，研究者手头都有内容可写。在AI时代，这个原则变得更为重要，因为AI本身不判断假设的学术价值，它只能基于统计模式生成``看起来像假设''的文本。判断一个假设是否可检验、是否值得做，仍然是人类研究者独有的能力。在这个意义上，可证伪性不仅是科学方法论的核心，也是人机协作边界的一条基线。
\end{funbox}

以下通过王工的案例展示从空白推导假设的具体方法。

\begin{userbox}
我确定的研究空白是：在旋转机械故障诊断的深度学习模型中融入轴承运动学和动力学物理约束，以期提升变工况下的泛化能力。

已有文献表明：(1) 纯数据驱动的深度学习模型在变工况下泛化性能显著下降（训练-测试工况不一致时准确率下降10-30\%）；(2) 物理信息神经网络（PINN）在其他工程领域（流体力学、固体力学）已初步验证了物理约束的正则化效果，但在故障诊断领域几乎空白；(3) 部分论文初步尝试了将故障特征频率作为辅助信息输入网络，但未将其作为硬约束融入损失函数。

请帮我从这个空白中推导出3-4个可检验的研究假设。每个假设需要满足以下标准：
1. 可证伪性：明确写出``什么实证结果会支持该假设''和``什么实证结果会否定该假设''
2. 已有支持：至少引用1条已有论文的间接发现作为该假设的间接支撑
3. 可操作性：在给定约束下（使用CWRU和Paderborn两个公开轴承数据集、NVIDIA RTX 3090级别的GPU算力、独立研究）是否可检验

假设的表述格式为标准的学术假设句式，并标注每个假设的理论来源。
\end{userbox}

\begin{aibox}
基于研究者的空白和约束，我推了三个假设：

\textbf{H1（核心）：物理约束提升泛化。}在损失函数里加轴承运动学方程作为正则化项，变工况准确率应显著高于纯数据驱动baseline。如果跨工况准确率提升不到3个百分点，或者训自己的工况反而下降，说明物理约束方向有问题或者加过头了。

\textbf{H2（辅助）：物理约束改善可解释性。}加约束后t-SNE图里同故障不同工况的特征应该聚得更紧。如果跟纯数据模型比没区别，说明约束对特征学习没起作用。

\textbf{H3（探索）：存在最优约束权重。}λ太小没用，太大模型学不到细节。在公开数据集上做网格搜索，如果λ跟性能是单调关系（越大越好或越差越好），说明物理约束方向本身需要重新考虑。

三个假设都在CWRU和Paderborn两个公开数据集上可检验，不需要额外数据。H1是主假设，直接回答``有没有用''；H2解释``为什么有用''；H3帮研究者调参。
\end{aibox}

\begin{critiquebox}
注意每个假设后面的``支持条件''和``否定条件''——这是这份清单最具价值之处。不管数据最后支持还是否定该假设，研究均有内容可报告。写不出否定条件的假设，不是真正可检验的假设。

但王工注意到一个问题：AI提出的H1以``跨工况准确率提升不到3个百分点''作为否定条件，却没有提供这个阈值的文献来源。3个百分点是AI的``合理推测''，不是基于方法学文献的标准。小林也碰到了类似情况——AI将她H1的β差值判断标准定为>0.05。AI给出的任何精确数值阈值，都只是基于训练数据中类似论文的统计推测，需要回到方法学文献中确认本领域公认的有实际意义的效应量标准，再据此调整。

\end{critiquebox}

\smallskip
\subsubsection*{\textbf{•可检验假设在不同学科中的典型形态}}同一套三要素框架，在不同学科中产出的假设形态完全不同。

小林（环境健康）的空白是``街景绿视率vs NDVI预测力直接对比''。AI帮她推导了三个假设，其中最核心的H3（双路径中介）提出了一个关键区分：绿视率更多地通过心理恢复感发挥作用，NDVI则更多地通过体力活动发挥作用。这一区分将研究问题从``哪个指标更强''的``测度差异''层面推进到了``为什么一个更强''的``机制差异''层面。她注意到风险：H3需要同时检验四个间接效应，n=680在统计检验力上可能处于临界状态，需要在讨论部分报告这一点。

张博（circRNA与胃癌）的空白是``circRNA编码短肽在胃癌微环境中的免疫调节''——极高风险、文献量<5篇。面对这种高风险方向，AI的策略是把H2坦率标注为探索性假设，并在否定条件里设了明确的止损线：如果预实验未检测到编码产物，就放弃免疫调节这个角度，转回更成熟的miRNA海绵机制。但他同时也留了出口：如果短肽确实存在但不通过PD-L1\footnote{PD-L1（Programmed Death-Ligand 1，程序性死亡配体1）：一种在肿瘤细胞和肿瘤微环境中表达的免疫检查点蛋白，通过与T细胞表面的PD-1受体结合抑制抗肿瘤免疫应答，是肿瘤免疫治疗的重要靶点。}发挥作用，还可以在RNA-seq\footnote{RNA-seq（RNA Sequencing，RNA测序）：利用高通量测序技术对细胞或组织中的RNA进行全转录组定量分析的方法，可同时检测mRNA、lncRNA、circRNA等多种RNA的表达水平。}数据里调整后续的功能验证方向。设止损线的同时留备选方向，是高风险课题的务实策略。


% --- 方案设计：给定约束下的最优路径 ---
\subsection{给定资源约束下的方案设计}

假设确定之后，接下来要回答的问题不是``理论上最好的方案是什么''，而是``研究者仅有单卡RTX 3090、两个公开数据集、八个月时间，能做什么''。

\subsubsection*{\textbf{•约束条件下方案设计的四个组件}}AI在设计方案时需要研究者明确给出实际约束（时间、设备、数据、预算），然后按四个组件输出：

\textbf{实验设计框架}——数据集划分策略、条件设置，需说明为什么该设计在约束下是最优选择及其核心局限；
\textbf{基线选择}——选哪些baseline进行对比，如果某个推荐超出约束条件，需提供备选方案；
\textbf{核心评估指标}——主指标和辅助指标各用什么，分别回答什么问题；
\textbf{统计分析计划}——每个假设对应什么统计方法，方法的前提假设是否满足。

拿到AI的方案后，重点审查两件事：AI推荐的数据集是否存在研究者自己领域里已知的局限（如CWRU被审稿人诟病为``玩具数据''），以及baseline比较是否公平（超参数调优程度不同可能导致结果不可比）。

\begin{userbox}
我的研究假设是：[H1物理约束泛化假设 + H2物理约束可解释性假设]。
我的实际约束是：独立研究（无企业合作数据）、GPU为单卡RTX 3090（24GB显存）、时间约8个月、可使用CWRU轴承数据集和Paderborn轴承数据集两个公开数据集。

请帮我设计一个在该约束下可行的研究方案：
1. 推荐实验设计框架（数据集划分策略、工况迁移任务的设置），说明为什么该设计在此约束下是最优选择，以及该设计的核心局限是什么
2. 模型选择建议：选择哪些baseline模型进行对比
3. 核心评估指标：使用哪些指标评估模型性能
4. 统计分析计划：针对每个假设，建议相应的统计分析方法

如果某个推荐超出了我的约束条件，明确指出并提供备选方案。
\end{userbox}

\begin{aibox}
给定约束下这个方案应该跑得通：

\textbf{实验设计：}k-fold跨工况迁移。选一个工况（比如1797 rpm、负载0）做训练集，其他几个工况做测试集。换3个不同的迁移方向各执行一次，每个方向训练时内部做5-fold CV。完全使用研究者现有的数据即可完成，不需要额外采集。唯一要注意的是CWRU和Paderborn都是实验室台架数据，讨论里得写清楚：结果可能高估工业场景表现。

\textbf{Baseline：}纯数据驱动用1D-CNN、LSTM、Transformer，当前方法用MMD domain adaptation和DANN。再加一组消融实验——只用故障特征频率那项物理约束、只用动力学那项，拆开看各自贡献。

\textbf{评估指标：}主指标用准确率和Macro F1，辅助看混淆矩阵（哪类故障跨工况最差）和t-SNE可视化（支撑H2）。

\textbf{统计：}H1用配对t检验（α=0.05），H2用Wilcoxon符号秩检验比类内距离，H3网格搜索λ画倒U曲线。
\end{aibox}

\begin{critiquebox}
王工仔细看了一遍，发现AI推荐CWRU数据集存在一个隐患。CWRU是故障诊断领域使用最广泛的公开数据集，但它的采集条件过于理想：单一轴承型号、实验室环境、无噪声干扰。现在越来越多的审稿人认为仅使用CWRU等于``在玩具数据上验证''。王工加上Paderborn数据集做第二验证，这个数据集包含不同轴承型号，实验条件也更接近工业现场。AI的建议里提到了Paderborn，但没有说明为什么要用两个数据集。这个判断需要王工自己完成：他对领域里关于数据集选择的最新讨论有清晰的了解。
\end{critiquebox}

\smallskip
\subsubsection*{\textbf{•约束条件下的方案设计：学科差异}}同一套``设计框架→数据方案→指标→统计''的模板，在不同学科中面临完全不同的约束类型。

小林（环境健康）希望进行纵向追踪，但硕士学制仅有两年。AI的方案给出了一个务实的折中：横截面作为主分析框架，用公开Landsat\footnote{Landsat：美国地质调查局（USGS）和美国国家航空航天局（NASA）联合运营的地球观测卫星系列，自1972年以来持续采集中等分辨率（30m）的多光谱地表影像，广泛用于植被监测和土地利用研究。} NDVI历史数据做回顾性辅助分析作为补充证据。回顾性分析不能替代纵向因果推断，但能让讨论部分的论证更扎实。她还确认了百度街景API在南京的覆盖率和更新频率满足课题需求。

张博（circRNA与胃癌）的约束最为严格：测序预算仅够20对样本，且无Ribo-seq\footnote{Ribo-seq（Ribosome Profiling，核糖体印记测序）：一种基于高通量测序的技术，通过捕获核糖体保护的mRNA片段来鉴定细胞内正在被翻译的RNA分子，可用于验证circRNA等非编码RNA是否具有编码蛋白质的能力。}设备需外送。AI设计了``三阶段递进''方案（发现→验证→功能），并给出了一个诚实的风险预案：如果阶段一未筛选到有翻译潜力的候选circRNA（可能性较高），则用同样20对样本的数据转向更成熟的miRNA海绵机制研究作为退路课题。这相当于为一个高风险课题设了止损线。


% --- 细节把关：样本量估算与方法时效性 ---
\subsection{样本量估算与方法时效性}

样本量估算是课题设计中特别容易出错的细节。AI能输出一个看似合理的数字，但它并不是在计算数学公式，而是在生成``在这个领域类似研究中看似合理的数值''。需注意：让AI提供框架和建议参数，研究者自己用G*Power\footnote{G*Power：免费的统计功效分析软件，用于计算研究所需的最小样本量，支持t检验、F检验、卡方检验等多种统计方法。}或者R的pwr包独立验算一遍。AI给出的数字仅作参考，以G*Power或R等专业工具的计算结果为准。

样本量在不同学科中的含义各不相同。王工的研究里样本量不是``需要多少被试''，而是``每个故障类别需要多少训练样本''。AI根据已有文献给出建议，但这个数字需要在公开数据集上实际运行一遍才能判断其可靠性。张博的细胞实验中，样本量则是``每组需要多少生物学重复（n=？）以确保统计检验力''，而这一数字取决于前期预实验的效应量估算——在效应量本身尚不确定的情况下，AI给出的样本量建议尤需谨慎对待。

另一个容易忽略的问题是AI推荐方法的系统性偏向：AI倾向于推荐更保守、更经典的方法。以中介效应检验为例，Baron \& Kenny在1986年提出的三步中介检验法\footnote{Baron, R. M., \& Kenny, D. A. (1986). The moderator-mediator variable distinction in social psychological research: Conceptual, strategic, and statistical considerations. \textit{Journal of Personality and Social Psychology}, 51(6), 1173--1182.}，被教材和几十年的论文反复引用，在训练数据中占据了绝对数量优势。但当前更受推荐的方法可能是基于Bootstrap\footnote{Bootstrap（自助法）：一种通过从原始样本中重复抽样（通常1000-5000次）来估计统计量分布的非参数统计方法。在间接效应检验中，Bootstrap置信区间法不依赖正态性假设，统计检验力通常优于传统的逐步检验法。}置信区间的中介效应检验——这一方法虽然更新且统计检验力更强，在训练数据中的出现频率却远不及经典方法。同样的偏向也出现在其他统计方法推荐中：AI更可能推荐t检验而非其非参数替代、更可能推荐逐步回归而非正则化方法。原因在于训练数据中经典方法的提及次数远多于新方法。

应对策略分两步。第一步，在提示词中加入时间限定，例如明确要求``优先推荐近10年该领域方法学文献推荐的方法''。第二步，拿到AI推荐的方法列表后，研究者需到近两年的方法学综述中做一次交叉验证，确认所推荐的方法是否仍为当前领域的最佳实践。这两步操作将AI从``方法推荐者''的角色重置为``方法候选列表生成者''。最终的方法选择权仍在研究者手中。该原则将在下一节的四步检验法中做进一步的操作化落实。


% --- 假设验证：让数据说话 ---
\subsection{四步检验法验证研究假设}

方案设计完成，实验执行完毕，数据收集到位。接下来不是急于撰写论文，而是先问自己：我手里的数据真的能回答我的假设吗？很多人在这个节点上会跳过检查，直接开始写作，但一旦审稿人发现分析方法与假设不匹配、数据质量存在隐患，或统计方法的选择已经过时，修改成本远高于提前检查。以下四步检验法旨在将这种风险降至最低。

\textbf{检查分析方案有没有漏洞。}王工在跑Transformer的变工况实验之前，让AI帮他审了一遍实验设计的统计方案：迁移方向的选择有没有系统性偏差？训练集和测试集的工况划分是否合理？配对t检验的前提假设（正态性、方差齐性）满足了吗？AI指出了两个问题——其中一个迁移方向工况差异太小（1797→1772 rpm，转速仅差25 rpm），可能测不出泛化能力的真实差距；配对t检验之前没做正态性检验，建议补一个Shapiro-Wilk\footnote{Shapiro-Wilk检验：一种用于检验数据是否服从正态分布的统计检验方法，原假设为数据来自正态分布。当p<0.05时拒绝原假设，表明数据显著偏离正态分布，此时应考虑使用非参数检验替代t检验。}检验。这两个问题都不致命，但如果不修，审稿人大概率会抓住。

\textbf{看结果有没有异常模式。}小林跑完中介模型后，AI帮她扫了一遍残差图和影响点诊断，发现三个受访者的Cook距离\footnote{Cook距离（Cook's Distance）：回归诊断中用于识别强影响点的统计量，衡量删除某个观测值后回归系数变化的程度。通常以Cook距离>1或>4/n（n为样本量）作为高影响点的筛查阈值。}远超正常范围。进一步查证后发现，这三个受访者同时缺失了体力活动数据，而中介模型恰好涉及体力活动变量。这不是统计方法的问题，是数据质量的问题。AI提醒她在论文里加一段缺失数据处理说明，并做一个完整的案例剔除敏感性分析，看看剔除这几个人后结论是否稳健。

\textbf{确认方法选择没有保守偏向。}前面讲过Baron \& Kenny三步法的问题：它可能比Bootstrap法的统计检验力低。同理，许多统计方法的选择都有时效性。研究者应该在拿到结果之前把分析方法确定下来，而不是跑完数据看哪个方法显著就用哪个。张博在分析自己的circRNA功能实验数据时，AI帮他检查了t检验的使用条件（每组n=6，正态性检验p>0.05，方差齐性检验p>0.10），确认t检验比Mann-Whitney U检验\footnote{Mann-Whitney U检验（也称Wilcoxon秩和检验）：一种非参数统计检验方法，用于比较两个独立样本是否来自同一分布。当数据不满足正态性假设时，该检验可作为独立样本t检验的替代方案。}更适合这个数据。理由有二：已知胃癌细胞系的circRNA表达数据在既往文献中符合正态分布，且样本量达到了统计检验力的最低要求。

\textbf{将验证结果与原始假设逐一比对。}数据支持了哪些假设？推翻了哪些？部分支持、部分否定的如何解释？AI在这里只能帮研究者整理结果，比如把每个假设的支持条件/否定条件和实际数据对照列出表格，但判断仍需研究者自己做。因为同一个结果，不同的人可能给出完全不同的解释。

\begin{critiquebox}
不少人在这个阶段会犯一个错误：数据不符合预期，就回去改假设。这是学术不端。假设验证的逻辑和提示词工程一样，不是一次就对的，需要反复调整。但调整的是实验设计和方法，不是假设本身。假设应该在实验开始之前就固定下来（哪怕只是在实验记录本上写一句``如果p>0.05则否定H1''），而不是看了数据再改。AI能帮研究者在这一步生成一份检查清单，但它不能替代研究者完成学术诚信这最后一道关。
\end{critiquebox}

\begin{figure}[H]
\centering
\BookFigureImage{images_chapter_05/假设验证四步法流程图.png}
\caption{假设验证四步法流程图}
\label{fig:ch5_11}
\end{figure}
\par

从空白到假设，从方案到验证，假设验证的完整流程如图\ref{fig:ch5_11}所示。完成这一阶段之后，研究者的结论就有了充分依据。进入下一节写作时，研究者对哪些结论可以确认、哪些需要审慎表述已心中有数。

\smallskip
\subsubsection*{\textbf{•从课题设计到项目申报}}本节讨论的五个环节——空白识别、假设推导、方案设计、假设验证和细节把关，恰好对应项目申报书（无论是学位论文开题报告还是基金申请书）中``立项依据''和``研究方案''两大部分的核心内容。

具体而言，这种对应关系是可操作的。空白识别的产出（「已知-争议-空白」三角框架的分析结果）可以直接转化为立项依据中的文献综述和问题提出：已知篇对应``国内外研究现状''的已有成果综述，争议篇对应``存在的科学问题''，空白篇对应``本研究的切入点''。假设推导的产出（可检验假设及其支持/否定条件）框定了``研究目标''和``关键科学问题''。方案设计的产出（实验框架、基线选择、评估指标、统计分析计划）填充了``技术路线''和``研究方法''。假设验证的预期结果对应了``可行性分析''，样本量估算则直接支撑``研究方案''中关于样本规模的数据。

在项目申报中使用AI时，需要注意三个层面的边界。第一，AI可以生成申报书各部分的结构框架和初稿，但不能替代研究者对研究意义的论证。第二，AI可以检查申报书中的逻辑断裂或表述不一致，但不能判断研究方案在给定预算和时间约束下是否确实可行。第三，AI生成的``创新点''陈述倾向于使用``首次''``率先''``突破''等措辞，必须由研究者在文献中逐一核实其真实性。

换言之，本节所学的AI辅助方法构成了项目申报中研究方案部分的AI辅助写作工具箱。关于基金申请书其他部分（研究基础、经费预算、团队成员介绍等）的AI辅助写作，详见5.7节传播转化部分。


% --- 提示词模板 ---
\subsection*{提示词模板}

\textbf{模板3-1：研究空白系统分析}

\begin{tcolorbox}[promptbox]
我计划研究 [研究主题]。以下是该领域 [N] 篇核心论文的摘要。

请做结构化研究空白分析：
1.「已知篇」：已有研究一致发现什么（至少3点，每点附支持该结论的论文比例）
2.「争议篇」：哪些结论存在分歧？各方核心依据是什么？
3.「空白篇」：分为两类——
   · 多篇论文的``未来研究方向''共同指出的空白（大家都能看到的，可能已经有人在抢了）
   · 只有一两篇论文提过、或者埋在别人论文``局限''段落里的缺口（值得多加关注）
   标注每个空白的评估（可操作性、竞争激烈程度、风险）
\end{tcolorbox}

\textbf{模板3-2：可检验假设的生成与评估}

\begin{tcolorbox}[promptbox]
我的研究空白是：[空白描述]。已有间接证据：[不超过100字]。实际约束：[约束条件]。

请从中推导3-4个可检验的研究假设。每个假设需附：
1. 理论来源（来自什么理论或哪篇论文的发现）
2. 可证伪性条件（什么结果支持、什么结果否定）
3. 可操作性评估（在当前约束下是否可检验）
4. 已有间接证据

按``核心竞争力''对假设排序：H1为最核心假设，H2/H3为支持性或探索性假设。
\end{tcolorbox}

\textbf{模板3-3：资源约束下的研究方案设计}

\begin{tcolorbox}[promptbox]
研究假设：[假设内容]。实际约束：[具体约束，包括时间/经费/设备/数据]。

请设计在该约束下可行的研究方案：
1. 推荐设计与数据方案及理由，说明核心局限
2. 核心变量/实验条件的设置方案
3. 评估指标与统计分析计划
4. 如果某个推荐超出约束，明确标注并提供备选
\end{tcolorbox}

\textbf{模板3-4：样本量估算框架}

\begin{tcolorbox}[promptbox]
研究设计：[类型]。核心分析方法：[方法]。预估效应量：[来源和大小]。
α=0.05，期望检验力=0.80。预测变量数=[K]。

请提供：
1. 最小所需样本量估算
2. G*Power对应检验类型和参数设置
3. 考虑15-20%无效样本率后，实际计划样本量

重要：此估算为参考框架，最终样本量需用G*Power独立验算。
\end{tcolorbox}

\textbf{模板3-5：假设验证前的分析方案检查}

\begin{tcolorbox}[promptbox]
我的研究假设是：[假设内容]。我的实验/调查已经完成，数据已清洗。
请帮我做以下验证前检查：
1. 分析方案与假设的对应关系——每个假设是否都有对应的统计检验方法
2. 统计方法的前提条件——正态性、方差齐性、样本量是否满足要求
3. 异常值和缺失数据的诊断——是否需要敏感性分析
4. 方法与领域标准的对比——推荐的统计方法是否与近两年文献一致
\end{tcolorbox}


% --- 常见错误与规避 ---

\begin{mistakebox}

\textbf{把AI生成的假设当成已经论证过的结论。}AI生成的假设是基于已有文献模式的统计拼凑，未经验证前只是需要检验的候选命题，不能作为论文引言中的立论基础。

\smallskip

\textbf{未核查AI推荐方法的时效性。}AI倾向于推荐训练数据中出现频率更高的经典方法（如Baron \& Kenny三步中介检验法），而非当前领域公认的最佳实践（如Bootstrap置信区间法）。拿到AI推荐的方法列表后，需到近两年方法学文献中交叉验证。

\smallskip

\textbf{样本量估算不亲自用专业工具独立验算。}AI输出的样本量数字是基于类似研究的统计推测而非精确计算，必须用G*Power或R的pwr包独立验算，以专业工具的计算结果为准。

\end{mistakebox}

\subsection*{小结}
本节在明确AI模式匹配与因果推理边界的基础上，用「已知-争议-空白」三角框架系统化地识别研究空白，从四个维度评估空白价值。可检验假设的三要素（可证伪性、已有支持、可操作性）和四步检验法帮助研究者将模糊想法转化为经得起检验的研究方案。AI负责广度——不遗漏，研究者负责深度——判断对。

\bigskip
\subsection*{课后习题}
\begin{enumerate}
\setlength{\itemsep}{6pt}

\item 推导研究假设时，为什么必须为每个假设明确写出``什么实证结果会否定它''？如果只写支持条件而不写否定条件，在研究设计和论文写作中会出现什么问题？

\item 在``已知-争议-空白''三角框架中，争议、独有空白和共识三者的价值排序是什么？为什么争议比共识更有价值？

\item AI在课题设计中进行的是模式匹配而非因果推理——请用一个具体的科研场景举例说明两者的区别，并说明这一区别对研究者意味着什么。

\item 四步检验法中，为什么要求``在分析数据之前先锁定分析方法''？这一步骤与避免p-hacking有什么关系？

\end{enumerate}
\newpage

% ============================================================
% 5.4 AI辅助论文写作
% ============================================================
\section{AI辅助论文写作}
\label{sec:5.4}

\begin{sectionkeypoints}
\item 理解AI学术写作的本质（统计模仿，没有意图，没有观点），并据此坚持``人写AI改''的铁律
\item 用漏斗式五段法生成引言大纲，每段明确功能（铺背景/综述已有/指缺口/亮问题/述方案）和信息密度（从粗到细）
\item 用方法报告规范（如STROBE）检查方法部分的遗漏项——AI在此任务上最可靠，因为只需模式匹配，不需领域判断
\item 用「逐句逻辑锚点检测法」找出论证链条中的因果跳跃、证据类型不匹配和推断过强
\item 让AI以目标期刊审稿人身份做投稿前预审，对审稿意见做三层分类处理（直接修改/查证后决定/仅供参考）
\end{sectionkeypoints}

\subsubsection*{\textbf{•AI写作的本质}}

经过前三节的工作，文献检索、课题设计与假设验证均已完成。研究者此刻手里有标签化的论文库、清洗好的数据库、注释完整的代码库。至此，写作所需的核心素材已经齐备。需要说明的是，上述描述将各阶段串联为线性流程，仅为叙述方便。实际科研过程通常需要在各阶段之间多次往返迭代：写作中可能发现需要补充文献，数据分析中可能发现假设需要修正，审稿后可能发现需要重做部分实验。研究者不应因反复迭代而感到挫败——这恰恰是科研工作的常态。

在让AI处理论文之前，需要先回顾5.1已经阐明的基本原理：AI学术写作的本质是统计模仿——它没有意图、没有论点、无法判断自己的论证是否成立，只是根据学术文本的统计规律逐词预测``这个位置最可能接什么''。

这直接决定了AI在论文写作中的能力边界。它最擅长的是\textbf{整理和组织}：将研究者已有的论点编排为清晰的逻辑链，找出论证链条上断开的位置，将同一层意思换几种表达方式。它最不擅长的是\textbf{贡献新观点和做判断}：每个看似新颖的观点，本质上都是训练数据中类似论点的统计拼凑；它只能判断``这个表述是否像学术论文里的典型论点''，而非``这个表述在学术上是否正确''。

由此得出本节的核心原则：\textbf{就论文正文而言，研究者写初稿，AI来修改，不可反过来。}研究者完成有判断、有立场的初稿，AI协助优化结构、检查逻辑、打磨表达。研究者负责``说什么''和``对不对''，AI负责``说得更清楚''。对于Cover Letter、审稿回复信、基金申请书等辅助性文书，AI可以协助起草初稿（详见\S5.6和\S5.7），但核心判断——接受哪些意见、强调哪些贡献、如何回应质疑——仍然必须由研究者亲自做出。

具体而言，人机协作在写作中可拆为三个层次：

\textbf{观点、判断、核心论证}：由研究者独立完成，AI不应参与；
\textbf{大纲、组织逻辑、多方案比较}：AI提出候选方案，研究者选择并调整；
\textbf{润色、格式、语言检查}：AI执行，研究者审核后确认。

如果把这个顺序颠倒过来——让AI出初稿，研究者做润色——结果往往是文字通顺但缺乏个人立场，更隐蔽的风险是：研究者会不自觉地信任一个``读起来很顺''的论证，降低对论证有效性的警惕。

\begin{funbox}
\textbf{为什么``读起来很顺''会骗过研究者？}心理学家发现了一种被称为``流畅性启发''（fluency heuristic）的认知偏差：人们倾向于认为认知加工越流畅的信息，越可能是真实的。一段文字如果字体清晰、措辞优雅、逻辑通顺，大脑在处理它时毫不费力，这种轻松感会被无意识地转化为信任感——``读起来这么顺，应该没问题吧？''

AI生成的学术文本恰恰是流畅性的极致：术语精准、句式工整、结构规范。正是这种完美流畅，构成了它最大的隐蔽风险。研究者在审读AI生成的文本时，认知负荷极低，流畅性启发在意识层面以下持续作用，使人不自觉地跳过对论证有效性的严格要求。这也是为什么本章反复强调``人写AI改''而非反过来——研究者自己写出的初稿可能有各种粗糙之处，但这些粗糙恰恰反映了思考的真实轨迹。修改自己的文字时，研究者知道每一处判断的来源，流畅性启发的影响被降到最低；而在审读AI写出的完美初稿时，研究者面对的是一个光滑如镜的表面，很难看出哪里有裂痕。
\end{funbox}

本节以小林为主例，展示从数据分析结果到讨论稿的完整打磨过程。王工和张博的写作场景作为对照穿插其中。论点链逻辑检测的完整流程如图\ref{fig:ch5_06}所示。

% --- 引言写作：漏斗式五步法 ---
\subsection{引言写作的漏斗式结构}

一篇论文的引言遵循的是漏斗结构：从大背景逐步收窄到具体的研究问题。新手写引言时常见的困难不是不知道自己的研究问题是什么，而是不熟悉收窄的节奏：每一段该放多少背景、在哪一步亮出文献缺口、研究问题放在什么位置最有力度。

AI在此处的合理用法是：研究者将自己的研究问题和3-5篇最相关文献的摘要提供给AI，让它生成一个漏斗式引言的大纲，每段标注``功能''（铺背景、综述已有、指缺口、亮问题、述方案）和``信息密度''（从粗到细）。然后研究者自己往框架里填入具体内容和学术判断。

\begin{userbox}
以下是我的研究信息和5篇最相关文献的摘要。请为我的论文引言生成一个漏斗式大纲：
- 研究背景：城市绿地与居民心理健康的关系
- 已有研究缺口：已有文献分别用NDVI和街景绿视率测量绿地，但鲜有在同一个研究中比较两者的相对预测力；亚洲高密度城市的证据极少
- 本研究目标：在南京社区居民中同时比较街景绿视率与NDVI对抑郁症状的预测力，并检验心理恢复感的中介效应
- 研究设计：横截面问卷调查，n=680

请提供一个五段式引言大纲，每段标注：
1. 该段的功能（铺背景/综述已有/指缺口/亮问题/述方案）
2. 该段的核心论点（一句话）
3. 需要引用什么类型的文献（不需要给出具体文献）
\end{userbox}

\begin{aibox}
\textbf{第1段「铺背景」：}城市化与心理健康的大背景→绿地作为潜在保护因素。
功能：建立研究议题的公共健康重要性。信息密度：粗——不需要在引言第一段就定义NDVI，先让读者理解``为什么这个问题重要''。

\textbf{第2段「综述已有」：}已有绿地-抑郁研究的两个主要发现方向：(1) NDVI与抑郁的负相关已有多国证据；(2) 机制上体力活动和心理恢复是两条主要路径。
功能：总结已知。信息密度：中——列出关键结论和代表研究，不用穷举。

\textbf{第3段「指缺口」：}两个缺口：(1) 不同绿地测量指标（NDVI vs 绿视率）的预测力缺乏直接比较——已有两篇论文结论不一致；(2) 现有证据几乎全部来自欧美城市，亚洲高密度城市语境中绿地的作用模式可能不同。
功能：确立本研究的必要性——不是``没人研究过''，而是``已有的结论之间有张力需要解决''。

\textbf{第4段「亮问题」：}直接用一句话提出研究问题：``本研究旨在南京社区居民中同时比较街景绿视率与NDVI对抑郁症状的预测力，并检验心理恢复感的中介效应。''
功能：明确告诉审稿人``这篇论文要回答什么问题''。

\textbf{第5段「述方案」：}简要说明横截面设计、n=680、南京单城市、一手问卷+街景+卫星数据。
功能：让审稿人知道``这个研究是怎么做的''，为方法部分做过渡。不需要写具体统计方法。
\end{aibox}

\begin{critiquebox}
AI生成的大纲是骨架，具体内容需要研究者自己填充。第1段和第2段的``宽背景''最容易写成泛泛而谈：AI只标注了``建立议题重要性''这一功能，但用什么事实、什么数据来建立这种重要性，只有研究者知道。第3段``指缺口''是引言中最关键的段落：AI只能基于研究者提供的摘要来判断缺口所在，它不知道研究者没有给它的那些文献。在定稿引言之前，需要回头确认：研究者宣称的``缺口''，最近两年是否已有人填补？

另外，学术界的引言并非只有这一种写法。有的期刊（尤其是Nature/Science系列）引言极短，一段之内完成``背景→缺口→问题''；有的领域（如人文学科）引言更长，更偏重叙事。拿到AI生成的漏斗式大纲后，应找一篇目标期刊近期发表的论文，对照检查引言段落结构是否与AI输出的一致。
\end{critiquebox}

% --- 方法部分：AI帮研究者做报告清单检查 ---
\subsection{方法部分的报告规范检查}

方法部分在学术写作中有一个特殊之处：它的内容是最``客观''的——用了什么数据、做了什么实验、跑了什么统计——但它的写作要求恰恰是最不容出错的。方法部分的读者（尤其是审稿人）在阅读时手持一把``这个领域公认应该报告的信息''的标尺，发现漏报即会产生质疑。

AI在方法部分的最佳用法不是替研究者起草，而是充当检查清单：研究者将写好的方法初稿交给AI，让它对照目标领域的方法报告标准（如STROBE\footnote{STROBE（Strengthening the Reporting of Observational Studies in Epidemiology，加强流行病学观察性研究报告规范）：用于横截面研究、病例对照研究和队列研究报告质量的国际指南，包含22个核心条目。}、CONSORT\footnote{CONSORT（Consolidated Standards of Reporting Trials，临床试验报告统一标准）：用于随机对照试验（RCT）报告质量的国际规范，包含25项核心条目和流程图模板。}、PRISMA等）逐项核对。

\begin{userbox}
以下是我论文``方法''部分的初稿。请对照STROBE横截面研究报告规范，逐项检查是否有遗漏：
- 研究设计与设置
- 参与者（纳排标准、招募方式）
- 变量（暴露、结局、协变量的操作性定义和测量工具）
- 统计方法（样本量估算、主分析、亚组分析、缺失数据处理）

如有遗漏，指出缺少了STROBE清单中的哪一项，并标注信息缺口的位置。
\end{userbox}

\begin{aibox}
逐项检查结果：

\textbf{已报告：}研究设计（横截面）、场所（南京）、参与者纳排标准、样本量（n=680，标注了问卷回收率和有效率）、PHQ-9量表信度（Cronbach's α=0.84）、绿地指标计算方法（GVI：百度街景API+语义分割，300m缓冲区；NDVI：Landsat 8，500m缓冲区）

\textbf{遗漏项：}
\begin{enumerate}[nosep]
\item \textbf{缺失数据处理方法}（STROBE第12项）：方法部分未说明缺失数据如何处理——是完整案例分析法（complete case analysis）还是多重插补（multiple imputation）？应按实际处理方式补充
\item \textbf{样本量估算依据}（STROBE第10项）：如果本研究在数据收集前做了样本量估算，需报告估算方法和参数。如果未做事前估算，报告事后检验力分析的结果
\item \textbf{敏感性分析}（STROBE第17项）：方法部分未提及任何敏感性分析
\end{enumerate}
\end{aibox}

\begin{critiquebox}
AI对照报告规范做检查是它最可靠的能力之一：这件事不需要领域判断，只需要模式匹配——``论文里是否写了X''与``规范里要求了X''的逐项比对。但研究者需要注意：AI有时会将``漏写''误判为``方法有问题''，它分不清是研究者做了但忘记写，还是根本没有做。拿到检查结果后，每一条都应先自行判断：这个项目我做了吗？如果做了，补写上去；如果没做，是否有合理理由？有理由则写入局限部分。

另一个提醒：领域特定的方法报告规范，AI可能知道也可能不知道。如果研究者所在的领域有专门的报告规范（如地理信息科学的数据质量报告标准、工程领域的实验验证最佳实践），最好将规范的关键条目一并提供给AI，让它对照检查，不要默认AI``应该知道''这些规范。
\end{critiquebox}

% --- 结果部分：区分描述与解读 ---
\subsection{结果部分的描述与解读边界}

论文各节中，结果部分是AI最不宜参与撰写的环节：它的核心任务就是诚实地报告数据，不需要任何外部知识。但这一环节看似简单，新手却容易犯一个错误：在结果部分就开始解读数据，把本该放在讨论中的判断提前说了出来。

AI在此处的用法很窄，但效果明确：帮助研究者识别结果部分中混杂在数据描述之间的解读性语句。

\begin{userbox}
以下是我论文``结果''部分的初稿。请逐句检查：哪些句子超出了``客观报告数据''的边界，进入了``解读数据含义''的范畴？标注出这些句子，并给出修改建议——要么挪到讨论，要么改为纯描述。
\end{userbox}

\begin{aibox}
\textbf{句3：``绿视率的β系数（-0.22）大于NDVI的β系数（-0.14），表明人眼感知的绿化比卫星测量的绿化覆盖对心理健康的影响更直接。''}

→ \textbf{检测：后半句超出了结果边界。}``表明……更直接''是对β系数差异的解读（为什么一个更大？），属于讨论范畴。建议：结果部分只写``绿视率与PHQ-9的标准化回归系数为β=-0.22（95\%CI: -0.31, -0.13, p<0.001），NDVI的标准化回归系数为β=-0.14（95\%CI: -0.23, -0.05, p=0.002），两者差异检验p=0.04''，差异大小的解读留给讨论。

\textbf{句5：``社区社会凝聚力的调节效应不显著，意味着在南京这样的中国城市中社会因素可能不是绿地-健康关联的核心机制。''}

→ \textbf{检测：整句超出了结果边界。}``意味着……''纯属讨论。结果部分应仅报告交互项的统计检验结果（β=0.03, p=0.34），不做任何解读。
\end{aibox}

\begin{critiquebox}
结果和讨论的边界只有一条：结果部分回答``数据是什么''，讨论部分回答``数据意味着什么''。新手最容易在报告完一个统计量之后顺手加上一句``这表明……''——``这表明''正是越界的信号词。拿到AI的标注后，不必直接删除越界句：如果解读方向正确，将其移至讨论部分作为论证起点；如果解读本身有误或措辞过强，则在讨论部分以更审慎的措辞重写。
\end{critiquebox}

\medskip

% --- 摘要写作：信息密度对标 ---
\subsection{摘要的信息核查与对标}

摘要是整篇论文最后撰写的内容，但不少人写到此时已经疲惫，便草率拼凑，仓促收尾。摘要是论文被阅读和检索的入口：多数数据库只索引标题和摘要，审稿人决定是否接审也主要依据摘要。一篇摘要若缺少效应量、置信区间、样本特征等核心信息，等于让读者和审稿人在缺乏关键判断依据的情况下做决定。

AI在此处可以做的事是：将研究者写好的摘要与目标期刊最近发表的3篇论文的摘要放在一起，做信息密度对标——检查研究者的摘要是否遗漏了该期刊摘要通常包含的关键信息项（如效应量、置信区间、样本特征），措辞强度是否与该期刊的惯例相符，以及摘要长度是否在期刊字数限制之内。研究者根据对标结果自行调整，不要直接让AI修改摘要：摘要的每一句都承载着研究者对研究贡献的自我定位，这个判断AI做不了。

以下通过小林的案例展示摘要对标的具体操作。她写完摘要初稿后，将摘要与目标期刊的三篇已发表论文摘要一并提交给AI，让AI逐项比对。

\begin{userbox}
以下是我的论文摘要初稿，以及目标期刊\textit{Landscape and Urban Planning}近两年发表的三篇论文的摘要（已粘贴在下方）。

请做信息密度对标分析：
1. 列出该期刊摘要通常包含的关键信息项（如研究背景、研究设计、样本特征、效应量、置信区间、局限性等），逐项检查我的摘要是否包含该项
2. 将我的摘要与三篇参考摘要进行措辞强度对比：我的摘要是否存在过度断言（如将关联表述为因果、省略限定语）或过度保守（如过度使用``可能''``初步''而弱化了确实显著的发现）
3. 检查摘要长度是否在期刊规定的字数限制之内

输出格式：先用表格列出信息项核对结果（包含/缺失/不完整），再给出措辞强度的逐句对比，最后给出修改优先级排序（必须修改/建议修改/可选修改）。
\end{userbox}

\begin{aibox}
\textbf{信息项核对结果（节选）：}

\begin{formalTable}[H]
\begin{tabular}{p{4.5cm} p{3.0cm} p{5.0cm}}
\toprule
\textbf{信息项} & \textbf{状态} & \textbf{说明} \\
\midrule
研究背景与问题重要性 & 包含 & 第1句已建立 \\
研究设计类型 & 包含 & ``横截面问卷调查'' \\
样本量与人群特征 & 不完整 & 仅写了``n=680，南京社区居民''，缺少年龄范围、性别比例、纳排标准简述 \\
核心效应量（β系数） & 包含 & 绿视率β=-0.22，NDVI β=-0.14 \\
置信区间 & 缺失 & 全文仅报告了p值，未报告任何95\%CI \\
中介效应占比 & 包含 & ``间接效应占62\%'' \\
研究局限 & 缺失 & 三篇参考摘要均在末句提及1-2个关键局限，本文摘要完全没有 \\
\bottomrule
\end{tabular}
\end{formalTable}

\smallskip
\textbf{措辞强度对比（节选）：}

原文：``绿视率对抑郁症状的保护效应显著强于NDVI''→ 问题：三篇参考摘要中，比较两个暴露指标的结论均使用了``was more strongly associated with''而非``的保护效应显著强于''。前者是关联措辞，后者含因果暗示。建议改为``绿视率与抑郁症状的负相关强于NDVI（β差值检验p=0.04）''。

\textbf{修改优先级：}
\begin{itemize}[nosep]
\item 必须修改：补充95\%置信区间（参考摘要均包含）；末句补充1-2个关键局限
\item 建议修改：补充样本人口学特征简述；将因果暗示措辞改为关联措辞
\item 可选修改：摘要长度目前为248词（期刊限250词），已接近上限，可酌情精简背景句
\end{itemize}
\end{aibox}

\begin{critiquebox}
摘要的信息密度对标效果较好，因为AI在做的事情本质上仍然是模式匹配：将研究者的摘要与三篇已发表摘要逐项比对``有没有X''。但有三点需要注意。

第一，AI的``必须修改''和``建议修改''分级是基于它在训练数据中学到的统计频率，不是基于该期刊的真实审稿标准。期刊的摘要信息要求有时在``投稿须知''中有明确规定——在采纳AI的建议之前，先去期刊官网逐条核对投稿须知中关于摘要的具体要求。

第二，AI对措辞强度的判断有时会过度谨慎。上面AI建议将``保护效应''改为``负相关''，这个方向是对的（横截面设计确实不宜使用因果措辞），但在某些情况下AI可能将领域内公认的合理措辞也标记为``过度推断''。最终措辞强度的决定权在研究者手中——尤其是当研究者对自身研究设计能支撑的推论强度有清晰判断时。

第三，摘要的修改务必在全文定稿之后进行。摘要是对论文最精炼的概括——如果在引言和讨论尚未定稿时就先改摘要，摘要中的措辞可能和正文产生出入。
\end{critiquebox}

% --- 结构先行：三种可选讨论结构方案 ---
\subsection{讨论部分的三种论证结构}

小林的数据分析已全部完成：回归表、中介模型、亚组分析，文件夹里存储了大量图表。但她面对空白的Word文档，二十分钟未能写出一个字。不是写不出句子，而是不知道该从哪个发现开始讲。将一组零散结果组织成一条有逻辑的论证线索，与进行数据分析是完全不同的认知任务。

\subsubsection*{\textbf{•讨论的三种论证结构}}AI在此处的核心价值不是替研究者写讨论，而是生成几个可选的结构方案，由研究者判断哪个方案最适合组织自己的论证。常见的三种结构各有适用场景：

\begin{itemize}[nosep]
\item \textbf{方案A「逐发现解读式」}：按发现顺序逐一解读。优点：结构清晰，审稿人容易跟踪。风险：如果发现之间逻辑关系不强，读起来像流水账。适用：发现数量少（2-3个）且各自独立的场景。
\item \textbf{方案B「先确定后意外式」}：先讲与已有文献一致的发现（确认性），再讲与预期不符的发现（意外性），最后寻找统一的解释框架。优点：审稿人容易判断研究者的贡献层次，意外发现有充足讨论空间。风险：如果意外发现的原因解释不够有力，会成为审稿人攻击的焦点。适用：至少有一个发现与已有文献或研究假设方向不一致。
\item \textbf{方案C「理论驱动式」}：以一条主线理论贯穿所有发现。优点：论证有深度，逻辑链条强。风险：如果理论覆盖力不够，部分发现会显得``牵强附会''。适用：有明确的、成熟的中心理论框架。
\end{itemize}

选定了结构后，下一步不是把整段讨论交给AI扩写，而是按照结构框架\textbf{自己逐段写出初稿}，然后用下一节的逻辑检测法检查论证链条。研究者负责每个发现的相关文献和判断，AI负责检查论证链是否断裂。

以下通过小林的案例展示讨论结构的选定方法。

\begin{userbox}
以下是我论文``讨论''部分需要覆盖的核心发现（五句话概括，每句包含核心效应量和统计显著性状态）：
1. 街景绿视率与PHQ-9抑郁评分的负相关（β=-0.22, 95\%CI -0.31, -0.13, p<0.001）
2. NDVI与PHQ-9的负相关弱于绿视率（β=-0.14, 95\%CI -0.23, -0.05, p=0.002）
3. 心理恢复感在绿视率-抑郁关联中的中介效应显著（间接效应占比62\%），而体力活动的中介效应不显著（p=0.21）
4. 社区社会凝聚力对绿视率-抑郁关联的调节效应不显著（交互项p=0.34）——这与我们参考的两篇国际文献的方向不一致
5. 性别亚组分析发现绿视率的保护效应在女性中更强（交互项p=0.03）

请为``讨论''部分提出3种可选的结构方案：
- 方案A「逐发现解读式」：按发现1→2→3的顺序逐一解读
- 方案B「先确定后意外式」：先讲与已有文献一致的发现，再讲与预期不符的发现
- 方案C「理论驱动式」：以注意力恢复理论为主线串联所有发现

每种方案标注：优点、风险、适用条件。
\end{userbox}

\begin{aibox}
\textbf{方案A「逐发现解读式」：}按发现1→2→3→4→5的顺序逐一解读。优点：结构清晰，审稿人容易跟踪每个发现与对应文献的对话。风险：发现4（调节效应不显著）和发现5（性别差异）之间缺乏内在逻辑联系，逐条解读可能使讨论后半段显得松散。适用条件：各发现之间彼此独立，不需要统一的解释框架。

\textbf{方案B「先确定后意外式」：}第一段解读与已有文献一致的发现（发现1-3），第二段聚焦意外发现（发现4：社会凝聚力调节效应不显著），第三段讨论亚组差异（发现5），第四段提出统一的解释框架。优点：审稿人容易识别研究的贡献层次，意外发现获得了充分的讨论空间。风险：如果对意外发现的解释不够有力，这一段会成为审稿人集中质疑的目标。适用条件：至少有一个发现与已有文献或研究假设方向不一致。——推荐方案，最匹配小林的情况。

\textbf{方案C「理论驱动式」：}以注意力恢复理论（ART）为主线串联所有发现。第一段阐述ART的理论预测，后续段落将每个发现与ART的具体命题对照。优点：整篇讨论有一条清晰的论证主线，逻辑链条强。风险：发现4（调节效应）和发现5（性别差异）与ART核心命题的关联较弱，可能显得``牵强附会''。适用条件：有明确的、成熟的中心理论框架，且所有发现均可被该框架覆盖。
\end{aibox}

\begin{critiquebox}
小林选择了方案B，因为她的发现4（社会凝聚力的调节效应不显著）与所参考的两篇国际文献方向不一致：这个``意外''恰好可以成为讨论部分最有价值的内容。选定了结构之后，她的下一步是按照结构框架自己逐段写出初稿，再用下一节的方法对初稿进行逻辑检测。
\end{critiquebox}

\medskip
王工和张博遇到的写作难题不同，但解法遵循相同的逻辑。

王工的故障诊断论文在``结果与讨论''部分面临的结构问题是：需要同时报告单一工况下的模型性能（10个工况×3个模型的对比表格庞大）和跨工况泛化性能（3个迁移方向×3个模型）。AI给出的结构建议是将``单一工况结果''压缩为一个摘要表格附在正文中，将``跨工况泛化结果''作为核心数据图放在正文中央，因为后者才是本文的核心创新贡献。这个建议与王工导师的看法一致。

张博的circRNA论文讨论部分的主要挑战则是：需要将circRNA-miRNA-靶基因轴放在胃癌信号通路全景中解读，而非仅仅报告``又发现了一个circRNA''。AI帮他生成了一个信号通路级的解读框架，将circ\_0008039→miR-375→YAP1轴与Hippo\footnote{Hippo信号通路：一条高度保守的细胞内信号转导通路，通过调控细胞增殖和凋亡来控制器官大小。在胃癌等多种癌症中，Hippo通路的核心效应因子YAP1常被异常激活，促进肿瘤发生发展。}通路在胃癌中的已知功能衔接起来；但AI无法判断哪些通路在胃癌中的文献支持最充分、哪些关联属于首次暗示。这些判断需要张博逐一查文献确认。


% --- 逻辑检测：逐句扫描论证链 ---
\subsection{论证链的逐句逻辑检测}

学术写作中有一类问题，比语法错误或措辞不当更难被发现：研究者自己看不出来的逻辑跳跃。从一个数据发现直接跳到一个因果结论，中间缺失了一步研究者自认为``不言自明''的推理，但读者未必能跟上这个跳跃。

AI在此处有一个天然优势：研究者容易对自己的论证产生``熟悉性盲区''，AI则不受此影响。研究者觉得理所当然的推理步骤，AI能迅速识别出推理链条中的缺失环节。

\begin{userbox}
以下是我论文``讨论''部分的初稿。请进行逐句逻辑锚点检测：

对每一句，标注：
- 该句的断言类型：[E]证据陈述（报告数据或已有文献）/ [C]因果断言（A导致B）/ [I]推断（从证据推导结论）/ [S]推测（超出当前证据范围的猜想）
- 如果该句是因果断言或推断，检查该句的断言在前文中是否有直接证据支撑。如果支撑不足，指出缺少什么类型的证据

输出格式：原句编号 + 断言类型 + 检测结果（``有支撑''/``支撑不足，需要补充[xxx]''）

不要修改句子本身，只做标注。
\end{userbox}

\begin{aibox}
\textbf{句4：[C]} ``绿视率的保护效应主要通过心理恢复而非体力活动来实现''

→ \textbf{检测结果：支撑不足。}前置句（句1-3）仅报告了中介模型的系数比较结果，未引用已有文献中关于``恢复体验的独立作用''的理论依据。建议在此句之前或之后补充：
\begin{enumerate}[nosep,label=(\alph*)]
\item 注意力恢复理论（ART）的核心观点——自然环境的``软性吸引''不依赖于体力消耗；
\item 引用已有文献中报告``恢复体验''独立于``体力活动''的中介效应的相关研究（如Dzhambov et al.\ 2018，该文经核实确实存在）。
\end{enumerate}

\smallskip
\textbf{句7：[C]} ``社区社会凝聚力的调节效应不显著，说明在中国的城市社区中，绿地与心理健康的关系不依赖于邻里关系质量''

→ \textbf{检测结果：逻辑跳跃。}``调节效应不显著''不能直接推导出``不依赖''。统计学上不显著的结果不能等于``零效应''（可能是样本量不足、量表测量精度不够、或社会凝聚力的操作性定义与西方不同）。建议修改为更谨慎的表述：

``在本研究样本中未观察到社会凝聚力的显著调节效应。这一结果可能反映了 (a) 中国城市社区的社会互动模式与西方不同，(b) 本研究的社区凝聚力量表可能未充分捕捉中国社区中与绿地使用相关的社会互动维度，(c) 样本量（n=680）可能不足以检测中等以下强度的交互效应。''
\end{aibox}

\begin{critiquebox}
小林拿到标注结果后，看到句7的检测报告时才意识到一个问题：她下意识将``交互项不显著''解读为``不存在调节效应''。这是统计基础课上反复强调的基本概念：不显著的p值只能说明``没有足够证据拒绝零假设''，不能等同于``零假设为真''。她自己在写作过程中完全没有察觉这个跳跃。AI的逻辑检测不会对研究者的论证产生这种盲区。
\end{critiquebox}

\medskip
\subsubsection*{\textbf{•逻辑检测在不同学科中的典型断点}}同一套E/C/I/S标注法，在不同学科中暴露的逻辑断点类型各不相同。

王工（机械故障诊断）的讨论中出现了\textbf{「性能→机制」习惯性跳跃}：他将``Transformer跨工况准确率更高''直接归因于``自注意力机制捕捉了全局依赖''。AI的标注明确划分了边界：数据能支持的是``它更好''，不能直接支持的是``为什么更好''。机制解释需要额外的验证手段（t-SNE可视化、attention权重分析等）。这个提醒促使他将讨论中多处类似措辞均做了降级处理。

张博（circRNA与胃癌）的讨论中出现了两种不同的断点类型：一是\textbf{证据类型不匹配}：引用文献中报告的是``细胞系''数据，他在论文中写成了``组织''数据，两种实验系统的结论不可直接等同；二是\textbf{「相关→标志物」推理跳跃}：p=0.02的生存关联不能直接推出``可作为预后标志物''，标志物需要额外的诊断效能评估（灵敏度、特异度、AUC\footnote{AUC（Area Under the Curve，曲线下面积）：接收者操作特征（ROC）曲线下的面积，用于评估诊断或预后模型的区分能力。AUC取值范围0.5-1.0，越接近1表示模型的判别能力越强，通常>0.7为可接受，>0.8为良好。}值、独立于TNM\footnote{TNM分期（Tumor-Node-Metastasis，肿瘤-淋巴结-转移分期系统）：国际通用的恶性肿瘤解剖学分期标准，T描述原发肿瘤的大小和侵犯范围，N描述区域淋巴结转移情况，M描述远处转移的有无。}分期的预后价值）。张博补充了AUC和多因素分析的要求后，讨论段的学术规范性显著提升。

三个人的逐句检测显示，这套操作的完整逻辑流是一致的：原文→逐句标注断言类型→检查支撑→修改。AI的逻辑检测的真正价值在于：不是因为它比研究者更聪明，而是它不会像研究者那样，在长时间的写作和反复修改中对自身的论证产生盲区（图\ref{fig:ch5_06}）。

\begin{figure}[H]
\centering
\BookFigureImage{images_chapter_05/论点链逻辑检测流程图.png}
\caption{论点链逻辑检测流程图}
\label{fig:ch5_06}
\end{figure}
\par


% --- 预审模拟：让AI扮演严苛审稿人 ---
\subsection{投稿前的AI模拟审稿}

初稿完成后，不要急于投稿。先让AI切换成审稿人的视角，对论文进行一轮系统性挑错。这是投入产出比极高的关键步骤。

\begin{userbox}
请以目标期刊的严苛审稿人身份，对以下论文全文进行批判性评审。每条批评需要：
1. 指向具体的段落或句子
2. 解释为什么这是一个问题（引用该领域的方法学标准或理论期望）
3. 建议具体的修改方案

评审维度：方法论稳健性、论证逻辑（是否有超出数据支持的过度推断、是否有未考虑的替代解释）、文献覆盖（是否遗漏了重要的已有研究）、表述规范

不需要写礼貌性的``本文也有以下优点''部分，直接给批评意见。
\end{userbox}

AI扮演审稿人之所以有效，是因为它的训练数据中既有海量的论文正文，也有海量的同行评审意见。当研究者让它切换身份，它调用的语言模式也随之切换：从``作者写论文''的模式，切换到``审稿人挑毛病''的模式。本质上，研究者是在利用AI训练数据中本就存在的两套话语体系，完成一次投稿前的低成本预审。

但这个操作有一个重要局限：AI不是真正的领域专家。它提出的批评经常是``方向正确，但诊断不准''。例如，研究者只是漏写了一个拟合指标，AI可能将其描述为``方法选择存在系统性问题''。因此，AI的审稿意见应视为``漏洞检测器''，而非``诊断报告书''。每条批评，都需回到文献中逐条验证。

AI审稿可能一次提出二十多条意见。不必因其数量而焦虑，也不必逐条照改。一个实用的策略是按``修改成本''和``修订必要性''两个维度分类，优先处理``高必要性+低成本''的项目。

\medskip
回到本章的三位研究者。王工也在投稿前让AI扮演了一次审稿人，目标期刊是IEEE Transactions on Industrial Electronics。

\begin{userbox}
请以IEEE Transactions on Industrial Electronics的审稿人身份，对以下论文全文进行批判性评审。每条批评需要：
1. 指向具体的段落或句子
2. 解释为什么这是一个问题（引用该领域的方法学标准或理论期望）
3. 建议具体的修改方案

评审维度：方法论稳健性、论证逻辑、文献覆盖、表述规范。
不需要写优点部分，直接给批评意见。
\end{userbox}

\begin{aibox}
主要问题（必须修改）：

1. 方法部分（第3.3节）：使用了t-SNE进行特征可视化，但未报告perplexity参数设置。t-SNE结果高度依赖perplexity的选择（van der Maaten \& Hinton, 2008），不同perplexity值可能导致完全不同的聚类形态。按该领域方法学标准，须报告perplexity值和选择依据。建议在3.3节补充``t-SNE使用的perplexity=30，选择依据为样本量的平方根近似值''。

2. 方法部分（第3.2节）：消融实验仅比较了``有物理约束''和``无物理约束''两种条件，但未区分物理约束中不同组成部分的独立贡献——故障特征频率项和动力学方程项是两种性质不同的约束，它们的贡献可能不同。建议增加一组消融：仅用故障特征频率项、仅用动力学方程项，分别与全约束和无约束baseline对比。

3. 实验部分（第4.1节）：仅使用CWRU数据集验证，未在更大规模或更接近工业条件的数据集（如Paderborn）上测试。该方向近年论文已普遍使用多数据集验证。建议补充Paderborn数据集的实验结果，或在讨论部分将仅用CWRU作为关键局限明确写出。

次要问题：讨论部分第5.2段将DANN和MMD方法的性能与本文方法对比时，未说明这些baseline方法的超参数调优过程——不同方法的调优程度不同，可能导致比较不公平。
\end{aibox}

\begin{critiquebox}
王工发现AI挑出的两个主要问题他此前确实没有注意到：t-SNE的perplexity参数他设了但忘记在论文中写出，消融实验实际上也跑了但``审稿人''在阅读论文时无法获知。这两个问题都属于``做了但未写''，AI审稿恰好帮他补上了这类漏洞。关于CWRU单一数据集的局限，他此前已有意识地在方案中加入了Paderborn，但AI审稿的建议促使他将Paderborn从``补充验证''提升为``核心验证数据集''，在两个数据集上均报告完整结果。
\end{critiquebox}

\medskip
张博用同样的方法让他领域里的AI审稿人审了一轮。

\begin{userbox}
请以Molecular Cancer的审稿人身份，对以下论文全文进行批判性评审。每条批评需指向具体段落或句子，解释为什么是问题，建议修改方案。评审维度：方法论稳健性、论证逻辑、文献覆盖、表述规范。直接给批评意见。
\end{userbox}

\begin{aibox}
主要问题（必须修改）：

1. 引言部分（第2段）：描述了circRNA在胃癌中的研究现状，引用了12篇文献。但引言中``circRNA在胃癌中的临床转化前景乐观''的表述为单向判断，建议补充引用该领域内对临床转化前景持不同评估立场的研究，使引言呈现更平衡的学术讨论。

2. 方法部分（第2.4节）：荧光素酶报告实验验证了circRNA与miR-375的直接结合，但仅使用了野生型结合位点——未包含突变型对照（miR-375结合位点突变后的荧光素酶活性应恢复至对照组水平）。突变型对照是排除荧光素酶活性变化来自非特异性结合的标准方法，缺失该对照会使miRNA靶标验证的证据不完整。建议补充突变型对照实验或解释为何未包含。

3. 讨论部分（第6段）：将circ\_0008039→miR-375→YAP1轴描述为``胃癌治疗的新靶点''，但全文仅提供了体外细胞系实验和裸鼠皮下成瘤模型的证据——未进行原位胃癌模型或PDX（人源肿瘤异种移植模型）的验证，也未评估该靶点的正常组织毒性。措辞建议降级为``可能成为胃癌治疗的候选靶点，但需要进一步在原位模型和正常组织毒性评估中验证''。
\end{aibox}

\begin{critiquebox}
\textbf{重大警告：AI审稿人会编造文献。}张博对第一条评语最为警觉：AI审稿指出他引言中``漏引了''一篇对立观点的文献。他立即去PubMed仔细检索，发现AI指出的那篇文献\textbf{根本不存在}。这个问题的机制与\S5.1节讨论的文献幻觉完全相同：AI不是在``检索''文献，而是在``生成''符合审稿话语模式的文本。处理方式也一样：AI审稿意见中每一条涉及具体文献的指控或推荐，都必须研究者在PubMed/Web of Science中逐条核实，确认文献真实存在后再决定是否引用。

回到张博的例子：虽然AI推荐的这篇文献并不存在，但它指出的方向是正确的——他的引言确实需要呈现更平衡的学术讨论。他将措辞从单向的``前景乐观''调整为``虽然已有初步诊断潜力，但关于临床转化前景的评估存在不同立场''，并在PubMed中自行找到了两篇确实存在的、对circRNA临床转化持谨慎态度的综述来支撑这一平衡表述。第三条关于``新靶点''措辞的批评也引起了他的警觉：讨论部分的措辞在反复修改中容易不自觉地升级强度，需要在每轮修改后进行自查。
\end{critiquebox}

\begin{funbox}
\textbf{学术期刊对AI写作的政策演变。}研究者在使用AI辅助写作时，除了遵循正文强调的``人写AI改''原则，还需要了解学术出版界对AI使用的正式立场。这一立场在过去数年间经历了快速演变：

2019--2022年，主流态度是``禁止AI作为作者''——AI不能出现在作者列表中，因为它无法对论文内容承担责任。2023年，随着ChatGPT的爆发，各大学术出版商迅速更新政策。多家主要出版商（包括Springer Nature、Elsevier、IEEE等）要求作者透明披露生成式AI的使用情况，并强调AI不得列为作者、研究者对论文的完整性和原创性负有最终责任。到了2024--2025年，政策的焦点从``是否可用''转向了``如何透明披露''——越来越多的期刊要求作者在提交时填写AI使用声明表，说明在研究的哪些环节使用了何种AI工具。具体披露位置（方法部分、致谢或专门的AI声明栏）和格式要求因刊而异，投稿前务必以目标期刊最新政策为准。

对于本书的读者，一个实用的底线是：在投稿前，务必查看目标期刊官网上的AI使用政策页面，确认该刊对AI辅助写作、数据分析和图表生成的具体要求。大部分期刊目前允许AI用于语言润色和结构优化，但要求作者在致谢或方法部分做出声明。未经声明而大量使用AI生成内容，在某些期刊已被视为学术不端的范畴。
\end{funbox}


% --- 提示词模板 ---
\subsection*{提示词模板}

\textbf{模板4-1：讨论部分结构方案生成}

\begin{tcolorbox}[promptbox]
以下是我的核心研究发现（5句话概括，每句附带效应量和统计显著性）：
[发现1-5]

请为``讨论''部分提出3种可选结构方案：
- 方案A「逐发现解读式」
- 方案B「先确定后意外式」
- 方案C「理论驱动式」
每种标注：优点、风险、适用条件。不对内容本身做任何修改。
\end{tcolorbox}

\textbf{模板4-2：论点链逐句逻辑检测}

\begin{tcolorbox}[promptbox]
以下是我论文 [某部分] 的初稿。请逐句进行逻辑锚点检测：

对每一句标注：
- 断言类型：[E]证据陈述 / [C]因果断言 / [I]推断 / [S]推测
- 如果该句是因果断言或推断，检查前文是否有直接证据支撑。
  支撑不足时，指出缺少什么类型的证据

不要修改句子本身，只做标注。输出格式为逐句检测表。
\end{tcolorbox}

\textbf{模板4-3：审稿人视角的批判性模拟}

\begin{tcolorbox}[promptbox]
请以 [目标期刊名] 的审稿人身份，对以下论文全文进行批判性评审。

每条批评指向具体段落或句子，解释为什么这是问题，建议修改方案。

评审维度：
1. 方法论稳健性（研究设计、测量、分析的可质疑点）
2. 论证逻辑（超出数据支持的过度推断、未考虑的替代解释）
3. 文献覆盖（遗漏的重要研究，尤其是与结论方向不一致的研究）
4. 表述规范（不清晰、不一致、易误读的表述）

不写优点部分，直接给批评意见。
\end{tcolorbox}


% --- 常见错误与规避 ---

\begin{mistakebox}

\textbf{让AI扩写初稿，结果AI在扩写过程中``补''入了假文献。}AI在扩写文本时会自动填补它认为``应该出现''的内容，其中可能包括不存在的引用。研究者自己写出包含真实文献的初稿，AI仅做结构和表达层面的修改，是防止这一问题的根本方法。

\smallskip

\textbf{多轮AI润色后，论文的个人学术风格被稀释。}每经过一轮AI润色，文本都会向训练数据中``平均化学术表达''靠近一步。五轮以上的润色可能使不同研究者的论文读起来高度相似。建议将AI润色控制在2-3轮以内，每轮之后研究者自行通读并恢复被AI改掉的个性化表达。

\smallskip

\textbf{对AI审稿意见全盘接受，不区分层级。}AI审稿意见中混杂着漏报检测（可靠）、方法质疑（需要文献验证）和创新评判（不可靠）。不加区分地逐条照改，既可能将正确的表述改成错误的，也可能在非关键问题上耗费过多时间。应遵循三层分类法：直接改的、查后再定的、仅供参考的。

\end{mistakebox}

\subsection*{小结}
本节以``人写AI改''为核心原则，依次展示了引言漏斗式五段法、方法报告规范检查、结果描述与解读边界的区分、讨论三种论证结构的选择、逐句逻辑锚点检测和投稿前AI模拟审稿。AI在写作中最有价值的能力是帮助研究者走出``熟悉性盲区''，但它不能替代研究者做核心论证和学术判断。

\bigskip
\subsection*{课后习题}
\begin{enumerate}
\setlength{\itemsep}{6pt}

\item 为什么学术写作应遵循``人写AI改''而非``AI写人改''？如果颠倒过来，会产生哪些风险？

\item 取自己论文的一段讨论，用E/C/I/S逐句逻辑检测法检查论证链，找出至少一处逻辑跳跃（如``相关→因果''``不显著→零效应''），并给出修正版本。

\item 漏斗式五段法引言中，每一段的功能分别是什么？五段中哪一段最关键的段落，为什么？

\item AI模拟审稿意见最常见的局限是什么？收到AI审稿意见后，应如何分层处理？

\end{enumerate}
\newpage

% ============================================================
% 5.5 AI辅助语言润色与翻译
% ============================================================
\section{AI辅助语言润色与翻译}
\label{sec:5.5}

\begin{sectionkeypoints}
\item 理解学术翻译的三个对等层级（词汇对等、概念对等、功能对等），知道AI在各层的可靠性差异
\item 用分步翻译法（语义准确→语域对齐→风格对标）替代一步到位翻译，每次只追一个目标
\item 在翻译启动前构建并锁定全文术语表，防止同一术语在不同段落出现不同译法
\item 识别并修复语义强度漂移——学术翻译中最隐蔽的系统性错误，通过限定语审计和回译检验来捕捉
\end{sectionkeypoints}

\subsubsection*{\textbf{•学术语言打磨的两个场景与三个对等层级}}

5.4节产出的初稿，在投稿前通常还需要经历两个环节。一是将中文初稿译为英文，二是对英文初稿（无论直接撰写还是翻译而来）进行语言润色。两个场景的操作存在重叠——其核心挑战都落在学术文本的对等转换上——因此本节将它们合并讨论，以翻译为主线展示操作方法，润色场景下的差异在相应步骤中标注。

学术翻译和文学翻译遵循不同的标准。文学翻译的偏差通常仅影响阅读体验；学术翻译的错误则可能直接影响审稿人对论文质量的判断。这类问题可以按层级来理解：有些错误较为显眼，有些则相当隐蔽——审稿人通读后只觉得``这篇论文的英语不够地道''，但未必能指出具体问题出在哪里。

最基本的层级是词汇对等。``抑郁症状''译为``depressive symptoms''，``归一化植被指数''译为``Normalized Difference Vegetation Index''。AI在这一层的表现基本可靠，尤其是在自然科学和医学这类术语高度标准化的领域。在润色场景中，词汇对等通常不是问题（英文初稿中术语错误相对少见），但在翻译场景中，术语误译一旦发生，后续的语域对齐和风格对标会将其``加固''得更深，修正成本成倍增加。

往上一级是概念对等，情形更为复杂。中国环境健康研究中常说的``心理恢复感''，按字面译为``psychological recovery feeling''，英文读者会判断作者不熟悉该领域的学术用语，因为英文学术语境中这一概念的标准表达是``perceived restorativeness''。概念对等要求研究者了解目标学术领域对每个概念的标准术语。AI在这一层具有优势，其训练数据中包含了大量英文文献的标准表达，能够为研究者提供术语选择建议。

最复杂的层级是功能对等，即整个段落在目标学术文化中的表达规范。中文论文讨论部分常写``结果表明……可能与……存在关联''，语气含蓄；英文期刊讨论部分的修辞习惯则更为直接——``This finding suggests that……''。功能对等不只涉及语言转换，还需调整论述强度、确定程度以及作者声音的可见度。AI在功能对等上最容易出现偏差，因为它的默认策略是保持原文的语义强度不变。当两种学术文化的默认修辞强度存在差异时，``保持不变''本身就是一种偏移。这一问题在润色和翻译中都会出现，5.5.3节将专门讨论。语义强度漂移的发生机制如图\ref{fig:ch5_07}所示。

% --- 翻译三步法：语义→语域→风格 ---
\subsection{学术翻译的分步操作法}

不少研究者将翻译视为一步到位的操作——输入中文，输出英文，即可使用。但如果仔细审读一篇由AI一次性翻译完成的英文论文，通常会发现它各项要求都触及了，却没有一项做到位：语义大致正确但关键句的措辞不够精准，语域接近学术英语但读起来并不自然，同一术语在不同段落出现了不同译法。

这一现象的根源在于翻译过程中有多个目标在同时竞争。AI在一轮生成中需要同时兼顾语义准确、语域合适、术语一致、句式流畅，结果便是每项仅达到中等水平。将翻译过程拆分开来，一次只专注于一个目标，能有效改善这一状况。对于英文润色场景，第一步（语义准确翻译）可以跳过，直接从第二步（语域对齐）或第三步（风格对标）开始。

\begin{userbox}
【第一步·语义准确翻译】

请将以下中文段落翻译为英文。这一步只追求语义准确性——术语正确、句义完整、不添加不删除任何实质性信息。不需要追求语言风格的流畅或学术化。输出时标注：每一句是否在语义上与中文严格对应。如果有无法确定的术语，在输出中标注``[term: 中文术语]''以提醒后续步骤。
\end{userbox}

\begin{aibox}
英文翻译（语义准确版）：

...（英文译文）...

术语标注：
- ``心理恢复感'' → [term: perceived restorativeness —— 见下文说明]
- ``街景绿视率'' → [term: street-level green view index (GVI)]
\end{aibox}

\begin{critiquebox}
第一步只追求准确性，不追求风格。遇到无法确定的术语，标注出来留给后续步骤处理，不要猜测一个填入。术语是后续所有修改的基础——第一步的术语错误会在语域对齐和风格对标两个步骤中被反复使用，修正成本远高于在第一步标注清楚后统一处理。
\end{critiquebox}

\medskip
\textbf{第二步：语域对齐。}

\begin{userbox}
【第二步·语域对齐】

以下是一段英文学术翻译的初始版本。请将此译文与目标期刊（Landscape and Urban Planning）的语域对齐。具体要求：
- 将口语化表达替换为正式的学术表达
- 将过长的从句拆分为学术写作可接受的长度
- 调整被动语态和主动语态的比例以匹配目标期刊的风格
- 保留第一步中的术语选择和句义，不在此步骤修改术语
- 输出时标注三到五处研究者做出了风格调整的位置及其调整原因
\end{userbox}

\begin{aibox}
【初始译文（第一步语义准确版）】
``Green space in cities can make people feel less depressed. We tested this idea using street-level images and satellite NDVI data in Nanjing. The results show that what people actually see (street greenness) matters more for their mental health than what satellites measure from above. The reason seems to be that seeing greenery helps people recover from mental fatigue --- physical activity does not seem to be the main pathway.''

【语域对齐版】
``Urban green space exposure has been associated with lower levels of depressive symptoms. The present study examined this association in a cross-sectional sample of 680 community-dwelling adults in Nanjing, China, using both street-level green view index (GVI) and satellite-derived Normalized Difference Vegetation Index (NDVI). The results indicate that GVI was more strongly associated with depressive symptom scores than NDVI (β = -0.22 vs. -0.14, p = 0.04 for the difference). Mediation analysis suggests that perceived restorativeness accounted for 62\% of the GVI-depression association, while physical activity did not emerge as a significant mediator (p = 0.21).''

\smallskip
\textbf{关键调整及原因：}
\begin{enumerate}[nosep]
\item ``can make people feel less depressed'' → ``has been associated with lower levels of depressive symptoms''：学术写作中相关性陈述不应使用使役结构（make），替换为被动+关联术语
\item ``We tested this idea'' → ``The present study examined this association''：Landscape and Urban Planning期刊以第三人称叙述为主（第一人称使用率<30\%），且``test this idea''过于口语化
\item ``what people actually see'' → ``street-level green view index (GVI)''：概念对等——英文学术界对GVI有标准术语，口语化描述必须替换
\item 补充了效应量（β = -0.22 vs. -0.14）和置信区间：学术讨论段必须包含精确数据，而这些数据在初始译文中被省略了
\item ``The reason seems to be'' → ``Mediation analysis suggests that''：将口语化的因果暗示替换为精确的统计方法指向
\end{enumerate}
\end{aibox}

\begin{critiquebox}
语域对齐这一步有一个容易被忽略的效果：它不仅帮研究者将口语化表达替换为正式学术表达，还可能暴露初始译文中遗漏的信息。上面这个例子中，第一步的译文省略了效应量和置信区间——这并非翻译错误，而是第一步只追踪语义准确、未追踪学术写作对信息密度的要求。AI在语域对齐时补充了这些数据。但这也正是研究者需要警惕的地方：AI补充的数据是否准确，必须回到原始分析结果中逐条核对。
\end{critiquebox}

\medskip
\textbf{第三步：风格对标。}

\begin{userbox}
【第三步·风格对标】

以下是我的英文学术论文译文，以及目标期刊Landscape and Urban Planning一篇已发表论文的``讨论''段落。请将两者在语言风格上进行对标分析——列出我的译文与已发表论文在以下维度上的差异：句长分布、从句密度、被动语态使用频率、模糊限制语（hedging）的使用频率和类型、第一人称代词的使用。然后据此调整我的译文。
\end{userbox}

\begin{aibox}
\textbf{风格对标分析：}

\begin{formalTable}[H]
\begin{tabular}{p{3.5cm} p{3.5cm} p{3.5cm}}
\toprule
\textbf{维度} & \textbf{研究者的译文} & \textbf{目标期刊已发表论文（Dzhambov et al.\ 2018）}\\
\midrule
平均句长 & 28.4词/句 & 22.1词/句 \\
从句密度 & 2.1个从句/句 & 1.4个从句/句 \\
被动语态比例 & 18\% & 42\% \\
Hedging频率 & 1.2个/100词 & 2.7个/100词 \\
第一人称 & 3次（we） & 0次 \\
\bottomrule
\end{tabular}
\end{formalTable}

\smallskip
\textbf{主要差距及调整建议：}
\begin{enumerate}[nosep]
\item \textbf{句式过长：}研究者的平均句长比目标期刊高约29\%。建议将超过35词的长句拆分为2-3句，特别是包含``whereas''或``while''的并列对比句——目标期刊偏好将对比发现分为独立两句陈述。
\item \textbf{被动语态偏低：}该期刊方法/结果部分被动语态通常占40-55\%，研究者的18\%显著偏低。调整时不要全局堆被动——在描述``本研究做了什么''（方法）和``前人发现了什么''（文献引用）时提高被动比例，在陈述``研究者的核心发现''和``研究者的判断''时保留主动。
\item \textbf{Hedging明显不足：}这是个需要重点关注的差距。研究者的译文每100词仅1.2个限定语，目标期刊讨论部分为2.7个。建议在每个因果推断或推广性陈述前至少加一个限定词——may、suggest、potentially、in part、appear to。
\item \textbf{第一人称：}该期刊几乎不使用第一人称（抽样段均未出现we/I）。建议将3处we替换为``the present study''或``this study''。
\end{enumerate}
\end{aibox}

\begin{critiquebox}
风格对标的价值不在于数字本身（句长22词与28词的差异），而在于它让研究者注意到平时不易察觉的语言习惯——例如讨论部分下意识使用了``we suggest''，但目标期刊从不使用第一人称；或者hedging密度仅为目标期刊已发表论文的一半。需要注意：对标分析给出的是该期刊的平均水平，不代表研究者必须逐条照搬。如果研究者的论证本身需要较复杂的句式（例如报告一个多重调节中介模型），将长句强制拆短反而会损害表达精度。AI提供参考数据，最终的调整方向由研究者自己判断。

此外，如果研究者使用AI生成风格对标的统计数据（如上文表格中的数字），需注意一个原则：AI给出的精确数值（如``平均句长22.1词''）可能是从训练数据中推算出来的，并非逐句统计那篇论文的结果。拿到数据后应先用常识进行合理性判断：如果AI报告一篇Landscape and Urban Planning论文讨论段的被动语态占80\%，研究者应当存疑；42\%则在合理范围内。如果对标数据将对修改方向产生决定性影响，建议研究者自行抽样统计确认。
\end{critiquebox}

分步法的核心逻辑是每次只让AI追踪一个目标。长论文无需每部分都执行全部三步——``方法''和``结果''部分术语密度高但风格差异较小，执行前两步即可；``引言''和``讨论''部分对风格和语气最为敏感，三步全部执行收益最大。对于英文润色场景，跳过第一步（语义转换），从语域对齐或风格对标直接开始。

\subsubsection*{\textbf{•英文润色的特殊考量}}

对于英文初稿已经完成、仅需语言润色的场景，上述分步法中的第一步可以省略，但润色场景有两个特有的风险需要警惕。

一是``均值回归''效应。每经过一轮AI润色，文本都会向训练数据中的平均化学术表达靠近一步。五轮以上的深度润色可能导致不同研究者的论文在语言层面高度趋同，个人学术风格被稀释。建议将AI润色控制在二至三轮以内，每轮之后研究者自行通读并恢复被AI过度修改的个性化表达。

二是润色边界模糊。翻译场景中，原文语义是天然锚点——偏离即错误。润色场景中，AI对``improve clarity''的理解可能超出研究者预期：将研究者有意使用的审慎措辞替换为更强烈的断言句式，或在``优化句式''的过程中改变原句的论证逻辑顺序。提示词中应明确划定边界：``保留所有专业术语和量化数据不变；仅在语法错误、明显冗余和不符合学术写作规范的句式上进行修改；任何涉及论证方向或结论强度的改动，在修改处标注[changed]等待研究者确认。''

润色完毕后，限定语审计和回译检验（详见5.5.3节）同样适用——润色场景下的语义强度漂移风险并不低于翻译场景。


% --- 术语锁定：翻译前的必备工序 ---
\subsection{翻译前的术语表构建}

用AI翻译长论文，最容易被忽略的问题是术语一致性。同一篇论文中，``绿视率''可能在第2段被译为``green view index''，第4段变为``street-level greenness''，第6段又变为``GVI''\footnote{GVI（Green View Index，绿视率）：基于街景图像计算的、反映人眼视角中绿色植被所占比例的指标，通常通过语义分割算法从全景街景图像中提取。}——审稿人会以为研究者使用了三个不同的暴露指标。这种现象的常见原因是分批次翻译：每次调用AI的时间不同，模型可能``遗忘''此前选定的译法，对同一术语给出不同的翻译。

解决这一问题最可靠的方法是：翻译开始之前，让AI生成一份全文术语表，之后每次翻译时将术语表附在提示词最前面，锁定术语选择。

\begin{userbox}
以下是我的论文中文全文。请从中提取约20个核心专业术语，生成中英对照术语表。每个术语的翻译依据以下优先序：
1. 如果目标领域已有标准的英文术语，使用标准术语（标注出处）
2. 如果目标领域没有统一标准，选择在英文学术文献中最常见的翻译（标注``基于文献使用频率选择''）
3. 如果是中国文化特有的概念（如``事业单位''``户籍''``双一流''），先给出直译再给出带简要解释的意译，供研究者选择

该术语表生成后，将在后续所有翻译中使用。请在本次输出中确认：``本术语表的术语选择已在全文范围内被锁定，后续翻译将保持术语一致。''
\end{userbox}

术语表生成后，工作尚未完成：逐条核查，确认每个英文术语确实是目标学术领域使用的标准表达。确认完毕后将术语表锁定，此后每次分段翻译时将术语表附在提示词最前面。

另一个需要注意的细节是中国文化特有概念的翻译——``事业单位''``双一流''``新农合''这类术语在英文中既没有词汇对等也没有概念对等。AI的默认策略是直译，但直译往往使英文读者无法理解其具体含义。研究者需要提前确定统一的处理策略：直译后加括号注释？采用功能对等翻译？还是保留拼音并附加脚注？关键在于整篇论文从头到尾使用同一种策略，避免在不同段落中对同一概念采用不同处理方式。


% --- 强度审计：语义漂移的捕捉与修复 ---
\subsection{翻译中的措辞偏移与修复}

语义强度漂移是学术翻译中最隐蔽的系统性错误。其表现是：原文中一个语气审慎的陈述，翻译之后变成了语气确定的断言。这不是译者有意为之，而是两种语言表达认识论审慎程度（epistemic caution）的方式不同——中文常通过上下文和语气词来传达不确定性，英文则依赖专门的模糊限制词汇（hedging devices）。

以下是一个典型的例子。

\textbf{中文原文：}``植被覆盖度越高的地方，居民的心理状态明显更好。''这句话描述的是一种关联，并未声称因果。

\textbf{AI英译：}``Vegetation coverage improves mental health.''译文的语法结构已经构成了因果断言。

关联漂移为因果，语义强度上升了一级。

反向的漂移同样存在。\textbf{英文原文：}``The association may be partly explained by……''作者使用了``may''和``partly''双重限定。\textbf{AI中译：}``该关联可以通过……来解释''——两个限定词全部丢失。

这种现象之所以是系统性的，根源在于5.1所述的逐词预测机制：面对两个候选选项——例如英文中``may suggest''与``suggests''——模型倾向于选择目标语言中出现频率更高的表达。英文学术语料中，更直接、更简洁的表达（``suggests''比``may suggest''短）在绝对数量上占优。因此AI存在一个微弱但持续的方向性偏差：选择更常见、更直接、限定更少的表达。一段文字包含十几个句子，每句发生一次微小的偏移，整篇论文的语义强度便系统性地高于原文。

针对语义强度漂移，有三项可以落地的预防措施：

\begin{itemize}
\setlength{\itemsep}{4pt}
\item \textbf{翻译时在提示词中明确要求保留限定语。}提示词中写明：``保留原文中所有模糊限制语——'可能''或许``在一定条件下''部分地倾向于'——不要为了流畅性而删除或弱化它们。``
\item \textbf{译文完成后进行限定语审计。}将英文译文中所有hedging词汇（may、might、possibly、potentially、in part、tend to、suggest、indicate、appear to等）逐一列出，与原文中的限定语数量进行对比。如果英文版hedging词汇明显少于中文原文，逐一回溯检查丢失的限定语。
\item \textbf{对核心论述句进行回译检验。}从讨论和结论部分选取5至10个核心句子，将英文译文回译为中文，与原文对照。回译中出现的语义强度变化，在正译中必然已经发生。
\end{itemize}

三种防范手段的完整流程如图\ref{fig:ch5_07}所示。

\begin{figure}[H]
\centering
\BookFigureImage{images_chapter_05/语义强度漂移示意图.png}
\caption{语义强度漂移示意图}
\label{fig:ch5_07}
\end{figure}
\par

\medskip
语义强度漂移并非仅限于上述学科，而是在各领域中普遍存在。化学论文翻译中，实验方法部分的产率85.6\%在译文中绝对不能变为86\%——小数点后一位的偏差看似微小，在学术上不可接受。专利翻译中，``所述方法包括''``其特征在于''这类表述有固定的法律译法，翻译失范可能直接改变权利要求的保护范围。经济学政策建议中，``建议''与``强烈建议''仅差两个字，但在翻译中必须严格保持原有的力度区分。归结起来是一条原则：无论学科，语义强度不能发生漂移。

\begin{funbox}
\textbf{为什么中文比英文更"含蓄"？——Hedging的跨语言差异。}正文描述了翻译中语义强度漂移的现象，其语言学根源在于中文和英文表达不确定性的方式存在深层差异。

英文学术写作拥有一套高度形式化的模糊限制语（hedging devices）：情态动词（may, might, could）、认知动词（suggest, indicate, appear to）、副词（potentially, presumably, arguably）以及句式层面的限定手段（``It is possible that...''）。这套语言工具经过了数百年学术写作传统的打磨，已经成为英文论文中标记科学论断不确定性程度的标准手段。中文学术写作同样表达审慎，但手段不同——更多依赖上下文语境、程度副词（"可能""似乎""在一定程度上"）以及句末语气词，而非固定的语法标记。

当AI执行中译英时，中文原文通过语境和语气词传达的审慎态度，往往缺乏可以直接映射的词汇锚点。AI的默认策略是选择英文中出现频率最高的对应句式——而更简洁、更直接的表达在英文语料中恰好在数量上占优。于是，中文中微妙的审慎便在逐词预测的过程中被系统性地"稀释"了。理解这一跨语言差异，研究者就能更主动地在翻译提示词中要求AI保留和补充限定语，而非依赖AI默认的"流畅翻译"。
\end{funbox}


% --- 提示词模板 ---
\subsection*{提示词模板}

\textbf{模板5-1：学术翻译三步法}

\begin{tcolorbox}[promptbox]
【第一步·语义准确翻译】
请将以下中文段落翻译为英文，仅追求语义准确性。不追求风格流畅。保留所有限定语。
如果某个术语无法确定标准英文翻译，标注``[term: 中文术语]''。

【第二步·语域对齐】
目标期刊：[期刊名]。请将此初始译文与目标期刊进行语域对齐：
- 口语→正式学术表达
- 过长从句拆分
- 调整被动/主动比例
保留第一步的术语选择。标注3-5处关键调整及原因。

【第三步·风格对标】
以下是我的译文和目标期刊一篇已发表论文的段落。请做语言风格对标分析
（句长分布/从句密度/被动频率/hedging频率/第一人称使用），然后据此微调译文。
\end{tcolorbox}

\textbf{模板5-2：术语表生成与锁定}

\begin{tcolorbox}[promptbox]
以下是我论文的中文全文。请提取 [N] 个核心专业术语，生成中英对照术语表。

翻译优先序：
1. 目标领域已有标准英文术语→使用标准术语并标注出处
2. 无统一标准→选择英文文献中最常见翻译并标注``基于使用频率''
3. 中国文化特有概念→给出直译+意译两个方案供选择

确认：本术语表将在后续所有翻译中被锁定，保持术语一致。
\end{tcolorbox}

\textbf{模板5-3：Hedging审计}

\begin{tcolorbox}[promptbox]
以下是我论文英译文中的关键论述段落。请：
1. 列出译文中所有包含模糊限制语（hedging）的句子及使用的hedging词汇
2. 逐句检查是否有：关联→因果、可能→确定、部分→全部的论述强度变化
3. 如有强度变化，给出修正版本
\end{tcolorbox}


% --- 常见错误与规避 ---

\begin{mistakebox}

\textbf{翻译或润色后未进行限定语审计，导致语义强度系统性漂移。}AI翻译存在微弱但持续的方向性偏差——倾向于选择更直接、限定更少的表达。一段文字中每句发生一次微小偏移，整篇论文的语义强度便系统性地高于原文。翻译和润色完成后必须进行hedging审计，对核心论述句执行回译检验。

\smallskip

\textbf{分段翻译时未锁定术语表，导致同一术语在论文不同部分出现不同译法。}分批次调用AI翻译时，模型可能在不同会话中为同一中文术语选择不同的英文对应词。读者（尤其是审稿人）会认为这些是不同的概念或指标，直接影响论文的可理解性。翻译启动前生成全文术语表并在每次翻译时将其附在提示词最前面，是解决这一问题的标准做法。

\smallskip

\textbf{中国文化特有概念在不同段落中采用了不同的翻译策略。}一些段落采用直译，另一些段落使用拼音加脚注，同一篇论文对``事业单位''等术语的处理方式前后不一致。研究者应在翻译启动前确定统一的策略——直译加注释、功能对等翻译、还是拼音加脚注——并适用于全文。

\end{mistakebox}

\subsection*{小结}
本节从学术翻译的三个对等层级出发，展示了分步翻译法（语义准确→语域对齐→风格对标）的操作流程。翻译启动前的术语表锁定和完成后的限定语审计是两项关键工序，后者专门针对学术翻译中最隐蔽的系统性错误——语义强度漂移。润色场景同样适用这些防范措施。

\bigskip
\subsection*{课后习题}
\begin{enumerate}
\setlength{\itemsep}{6pt}

\item 学术翻译中为什么要对AI译文做限定语审计？什么样的限定语变化最值得警惕？

\item 学术翻译的三个对等层级分别是什么？AI在各层的可靠性有何差异？

\item 翻译长论文时，为什么要在翻译启动前先构建并锁定术语表？如果不锁定，可能出现什么问题？

\item 什么是语义强度漂移？为什么说它是学术翻译中"最隐蔽"的系统性错误？回译检验如何帮助发现这一问题？

\end{enumerate}
\newpage

% ============================================================
% 5.6 AI辅助自我审稿与投稿准备
% ============================================================
\section{AI辅助自我审稿与投稿准备}
\label{sec:5.6}

\begin{sectionkeypoints}
\item 理解同行评审的三层逻辑（技术把关、实质评价、创新判断），知道AI在各层的可信度和参与边界
\item 用三层分类法处理AI审稿意见：可直接修改的（漏报指标、格式）、需核实后再决定的（方法质疑、论证强度）、仅作参考的（创新评分）
\item 用多维匹配策略选择目标期刊：综合评估主题、方法、数据量级、创新度和审稿周期五个维度的匹配度
\item 回复审稿意见时坚持战略性判断——AI起草回复习惯全盘接受，在哪些问题上让步、哪些问题上坚持，只有研究者能做决定
\end{sectionkeypoints}

\subsubsection*{\textbf{•同行评审的三层逻辑与AI能力边界}}

投稿前的最后一轮检查——前文构建的论文库、数据库和代码库在此环节发挥关键作用：AI审稿指出研究者漏引了某篇文献，论文库中按标签即可定位；AI质疑数据质量，数据库中的数据字典和清洗记录即是最直接的辩护材料。关于投稿前让AI扮演审稿人的操作方法，5.4已做了详细演示。在此基础上，要有效利用AI审稿，还需要理解同行评审的三个层次，以便对AI的审稿意见做出分级处理。

第一层是技术把关：方法是否完整报告、数据分析是否正确、引用是否齐全、格式是否规范。这些问题存在明确的判断标准，AI在这一层具有参考价值。模型在训练过程中接触了大量学术论文，习得了学术写作的规范模式，能够帮助研究者识别论文中偏离常规规范的部分。

第二层是实质评价：研究问题是否成立、证据是否足以支撑结论、贡献是否达到目标期刊的门槛。这一层需要的不是格式层面的判断，而是在领域知识框架中做出评估。AI在这一层的能力已接近边界：它能够指出论文中缺少了哪些要素，但无法判断缺少这些要素之后论文是否仍然达到发表标准。

第三层是创新判断：论文是否为领域带来了新的知识、是否会改变现有认知。在这一层，AI不具备可靠的判断能力。它无法识别``什么是新的''——其知识停留在训练截止日期，无法知晓截止日期之后是否有类似研究出现。

操作原则：\textbf{将AI的审稿意见分为三个层级，分别采取不同的处理策略（图\ref{fig:ch5_08}）。}

\begin{figure}[H]
\centering
\BookFigureImage{images_chapter_05/同行评审三层逻辑与AI能力边界图.png}
\caption{同行评审三层逻辑与AI能力边界图}
\label{fig:ch5_08}
\end{figure}
\par

% --- 模拟审稿：投稿前的系统性漏洞检测 ---
\subsection{投稿前的系统性自我审稿}

关于投稿前让AI扮演审稿人的具体操作方法（提示词模板和完整对话示例），5.4已做了详细演示，此处不再重复。本节聚焦于一个5.4未展开的关键问题：收到AI审稿意见后如何分级处理。不应立即逐条修改，而是首先将意见归入以下三个层级：

\textbf{第一层：可直接修改的。}包括漏报的统计指标、前后不一致的术语、明显漏引的文献（在数据库中确认后补入）、格式和引用不规范之处。AI对这类问题的判断通常可靠，经研究者确认后即可修改。

\textbf{第二层：需核实后再决定的。}包括对方法选择的质疑（例如``为什么使用横截面设计而非纵向设计''——若纵向设计超出了研究者的资源约束，在论文中解释清楚即可，无需修改研究设计）、对论证强度的质疑、替代统计方法的建议。这一层的问题需要在文献中核实后做出判断。

\textbf{第三层：仅作参考，不宜作为修改依据的。}包括创新性评分、贡献评分，以及``该研究是否值得做''这类价值判断。此类判断即使在真实的同行评审中，不同审稿人之间也经常存在分歧，AI的判断更不应作为修改的依据。

\begin{critiquebox}
AI审稿意见最常见的问题是：\textbf{问题定位准确，但原因诊断有偏差。}例如，AI指出``方法部分结构方程模型拟合指标报告不完整''——这一判断的方向是正确的，研究者确实遗漏了SRMR（标准化残差均方根，结构方程模型拟合指标之一，通常要求SRMR<0.08表示模型拟合可接受）的报告。但AI将原因诊断为``指标选择依据不足，方法论可能存在问题''，而实际情况是模型拟合良好，仅是撰写时遗漏了该指标的汇报。AI审稿的盲区在于：它能判断哪些内容偏离了报告规范，却无法辨别这种偏离的原因——是遗漏了必要信息，还是在研究设计上确实存在缺陷。因此，建议按以下流程处理AI的审稿意见：将每一条批评定位到原文的具体位置；技术报告类问题先检查是否确实存在信息遗漏；方法论层面的质疑回到文献中核实；论证逻辑类问题由研究者逐条分析判断。
\end{critiquebox}


% --- 期刊匹配：多维度评估与三级候选 ---
\subsection{期刊选择的多维匹配策略}

选择目标期刊时，不应仅关注影响因子。需要综合评估以下维度：

\begin{itemize}[nosep]
\item 论文主题与期刊定位的匹配程度；
\item 所使用的研究方法在该刊是否受认可；
\item 数据量级是否达到该刊的通常标准；
\item 创新程度是否满足该刊的门槛要求；
\item 审稿周期是否在研究者的时间约束之内。
\end{itemize}

期刊选择不当，论文可能在编辑初审阶段即被退回（desk reject），无法进入外审环节，审稿意见也无从获取。

\begin{userbox}
我有一篇待投稿的论文，特征如下：
- 研究主题：街景绿视率与城市居民抑郁症状的关联及心理恢复感的中介效应
- 研究方法：横截面问卷调查，n=680，单城市（南京），使用多元回归+并行多重中介模型
- 数据类型：一手问卷调查+街景图像+NDVI卫星数据
- 创新程度自评：中等偏上——已有文献中绿视率与抑郁的直接研究较少，且比较了绿视率与NDVI的相对效应，但研究设计为横截面
- 对审稿周期的期望：中等（毕业前需要接收函，但尚有约10个月缓冲期）

请推荐三个层级的候选期刊：
- Tier 1（高竞争，匹配度高，难度大）：2-3个
- Tier 2（中等竞争，匹配度最好，建议主攻）：2-3个
- Tier 3（接受率高，审稿快，保底选择）：1-2个

每个期刊标注：定位与偏好、与本文匹配度的逐维评估（主题/方法/数据量级/创新度各1-5分）、审稿周期的估计范围。
\end{userbox}

期刊推荐返回后，去JCR或Scopus\footnote{Scopus：爱思唯尔（Elsevier）维护的多学科摘要与引文数据库，覆盖自然科学、社会科学、医学和工程技术等领域，提供期刊的CiteScore、SJR（SCImago Journal Rank）和SNIP（Source Normalized Impact per Paper）等学术影响力指标。}核实影响因子和分区——期刊排名每年更新，AI的数据是快照。AI还可能遗漏近两年新创刊的期刊或新兴Open Access期刊（训练数据中不包含这些信息）。一个补充做法：在Web of Science中按主题检索最近两年该主题的论文主要发表在哪些期刊，与AI推荐的清单进行交叉对比。

期刊选择有一个基本原则：\textbf{匹配度比影响因子更重要。}若论文的数据量有限、创新程度一般，勉强投向匹配度低的高影响因子期刊，被编辑初审直接退稿的风险较高。因为编辑在初审阶段主要依据主题匹配度和创新性做判断，匹配度不足的稿件往往到不了外审环节，审稿意见也无从获取。选择匹配度高的期刊，可以降低因选刊不当而被编辑直接退稿的风险，并提高进入外审的可能性。论文能否进入外审还取决于创新性、方法质量和编辑判断等多种因素，匹配度并非唯一决定因素。根据审稿意见修改后更有可能获得理想的发表结果。


% --- 投稿实操：Cover Letter、审稿人与回复信 ---
\subsection{投稿信、审稿人与回复信}

Cover Letter是编辑阅读的第一份材料，需要在短短一两分钟内让编辑理解四个核心问题：论文的研究内容是什么？为什么适合在该期刊发表？创新贡献是什么？是否存在一稿多投等投稿合规问题？

AI可以协助起草Cover Letter初稿，但以下事项必须由研究者亲自完成：主编姓名需到期刊官网查询确认（不可使用AI生成的可能过时的信息）；全文语气需进行调整，避免模板化的生硬感；声明部分需逐条核对，确保符合目标期刊的具体要求。

推荐审稿人是另一个需要研究者亲自把关的环节。AI可以根据研究方向推荐候选人，但它不了解研究者的人际网络和利益冲突关系——包括合作者、同机构同事、近期共同发表者、导师的前学生等。若推荐了存在利益冲突的审稿人，可能构成学术伦理问题。研究者必须对AI推荐的每一位候选人进行逐人核查。

收到审稿意见后，不宜立即着手修改。先将所有意见按主题归类（方法类、论证逻辑类、文献覆盖类、表述规范类），并标注出不同审稿人之间意见相左之处。逐条回复的结构包含三个要素：感谢审稿人的具体建议、说明研究者做了哪些修改、指出修改在论文中的具体位置。

AI可以协助起草回复初稿，但有一项判断无法替代研究者：在哪类问题上做出让步，在哪类问题上坚持原有立场。AI在起草回复时倾向于全盘接受审稿人的所有意见，但实际审稿过程中部分意见未必合理或可行。这一战略性判断必须由研究者亲自完成。

% --- 提示词模板 ---
\subsection*{提示词模板}

\textbf{模板6-1：投稿前模拟审稿}

\begin{tcolorbox}[promptbox]
请以 [目标期刊名] 的审稿人身份对以下论文进行结构化评审。

维度：
1. 原创性与贡献（1-10分+简短理由）
2. 方法论稳健性（逐一指出可质疑的点）
3. 论证逻辑（证据支撑充分/不充分的结论、替代解释）
4. 文献覆盖（遗漏的重要研究）
5. 可读性与结构

输出格式：按``主要问题''和``次要问题''分别列出。每个问题指向具体段落或句子。
\end{tcolorbox}

\textbf{模板6-2：目标期刊匹配分析}

\begin{tcolorbox}[promptbox]
待投稿论文特征：[主题]、[方法类型与样本量]、[数据类型与范围]、[创新程度自评]、[审稿周期偏好]

请推荐Tier 1/Tier 2/Tier 3三级候选期刊，每个期刊标注主题匹配度、方法匹配度、数据量级匹配度、创新度匹配度（各1-5分）和综合判断。
\end{tcolorbox}

\textbf{模板6-3：审稿意见分类与回复}

\begin{tcolorbox}[promptbox]
以下是我收到的审稿意见全文。请将每条意见按类别分类并标注修改难度：
- 方法论类 / 论证逻辑类 / 文献覆盖类 / 表述规范类

对审稿人之间存在的矛盾意见做出标记。输出修改优先序建议表格。
\end{tcolorbox}


% --- 常见错误与规避 ---

\begin{mistakebox}

\textbf{未加甄别地接受AI对方法论的批评。}AI审稿常将``漏写''误判为``方法有问题''——技术报告类问题先检查是否确实遗漏了必要信息，方法论层面的质疑需回到文献中核实后再决定。

\smallskip

\textbf{推荐审稿人时未核查利益冲突。}AI不了解研究者的人际网络，可能推荐存在利益冲突的候选人（合作者、同机构同事、导师的前学生等），必须逐人核查以避免学术伦理问题。

\smallskip

\textbf{回复审稿意见时在所有问题上让步，未区分哪些意见值得坚持。}AI起草回复时习惯全盘接受所有意见，但实际审稿中部分意见未必合理或可行，在哪些问题上让步、在哪些问题上坚持，只能由研究者亲自判断。

\end{mistakebox}

\subsection*{小结}
本节将同行评审分为技术把关、实质评价和创新判断三个层次，明确了AI在各层的可信度。投稿前让AI扮演审稿人进行预审，意见采用三层分类法处理。期刊选择遵循``匹配度优先于影响因子''的原则，回复审稿意见时``在哪些问题上让步、在哪些问题上坚持''的判断只属于研究者。

\bigskip
\subsection*{课后习题}
\begin{enumerate}
\setlength{\itemsep}{6pt}

\item AI模拟审稿的意见应如何分层处理？哪些意见可直接采纳、哪些必须核实后再决定、哪些仅作参考？

\item 选择目标期刊时，为什么匹配度比影响因子更重要？

\item AI在起草审稿回复时有什么系统性偏向？在哪些问题上研究者不应让步？

\item 同行评审的三个层次分别是什么？AI在各层的可信度如何？

\end{enumerate}
\newpage

% ============================================================
% 5.7 AI辅助多模态信息处理与内容创作
% ============================================================
\section{AI辅助多模态信息处理与内容创作}
\label{sec:5.7}

\begin{sectionkeypoints}
\item 理解知识转化中的三条规律（信息压缩律、话语转换律、传播伦理律），明确在传播中信息可以压缩，但科学结论的性质、方向和确定性不可改变
\item 区分三种传播形态的信息组织逻辑：海报是浏览式（依赖视觉层级），PPT是听觉线性（依赖三段重复），科普是注意力竞争（结论先行）
\item 用禁忌词清单和``更准确地说''段落在科普转化中守住科学底线——相关不能变为因果，可能性不能变为确定性
\item 理解基金申请书与论文的写作逻辑差异：论文陈述``已经完成了什么''，基金论证``将要做什么以及为什么能够完成''
\end{sectionkeypoints}

\subsubsection*{\textbf{•知识转化中的信息压缩与传播伦理}}

论文投稿之后，传播工作随之展开。研究者在论文库和数据库中积累的结构化信息——经过5.2的三库体系整理——此时可以直接调取使用，无须反复翻阅原始数据。同一项研究成果需要转化为学术海报、口头报告PPT、科普文章和基金申请书等多种形态，每种形态面对的受众不同，信息筛选的标准也不同。处理不当，论文中最核心的科学发现可能在传播过程中被曲解。

在这一转化过程中，有三条规律贯穿始终。

\textbf{规律一：信息压缩律。}从一篇完整论文到一篇千字科普文章，绝大部分信息在压缩过程中被舍弃。最先消失的往往是方法细节和限定条件：它们在论文中是不可或缺的组成部分，但在面向公众的传播中阅读门槛最高。核心结论最容易在压缩后保留下来，然而一旦脱离了原文中的限定条件，它便可能构成误导。问题的关键不在于压缩了多少信息，而在于压缩掉了哪些信息、保留了哪些信息。

\textbf{规律二：话语转换律。}话语体系每经过一次转换，精度便随之下降一个层级。``β=-0.18, p<0.01''变为``存在负面影响''，再变为``有害''——每次转换都伴随着信息精度的损失。恰当的传播实践会在每一步转换中补充新的限定语加以对冲，以控制精度损失的程度。

\textbf{规律三：传播伦理律。}有一条底线不可逾越：科学结论的性质不能被改变。相关不能变为因果，方向不能反转，可能性不能变为确定性。AI在辅助撰写传播文案时，统计上天然倾向于生成更吸引人而非更谨慎的表述，这是因为训练数据中传播类文本本身即比学术文本更为直接。因此，传播环节的审校需要比写作环节更为严格（图\ref{fig:ch5_09}）。

\begin{figure}[H]
\centering
\BookFigureImage{images_chapter_05/知识转化三规律示意图.png}
\caption{知识转化三规律示意图}
\label{fig:ch5_09}
\end{figure}
\par

% --- 多形态输出：海报、PPT与科普文章 ---
\subsection{从论文到海报与PPT}

同一项研究在不同场合面对不同受众，需要采用完全不同的叙事方式。小林关于城市绿地与抑郁的研究，现在需要同时制作三种形态：一张学术海报、一份口头报告PPT、一篇科普文章。

\textbf{学术海报}与论文的本质区别在于信息接收方式。论文支持线性深度阅读：读者可以随时停下来思考，也可以回看前面的内容。海报则是浏览式阅读：一个人站在海报前通常不超过三分钟，视线移动路径未必线性。因此，海报的设计应遵循三项原则：以视觉中心吸引注意力，以分区布局支持快速扫描，以关键词确保核心信息被记住。

AI无法直接生成完成的海报，但可以输出实用的设计方案，研究者依据该方案在PowerPoint或Illustrator中制作即可。以小林的论文为例，AI给出的分区方案是：标题和作者栏占顶部15\%；研究问题和背景置于左上区域（占20\%高度、40\%宽度）；方法部分置于左下区域；核心结果置于中间和右侧上半区域（视觉锚点区域）；结论和意义置于右下区域。设计底线为标题不小于72pt，正文不小于28pt，确保1.5米外可以看清。

\textbf{口头报告PPT}的信息组织遵循一个核心原则：先告知听众将要讲什么，然后展开讲述，最后再次告知听众刚才讲了什么。这一``三段重复''原则的根源在于听众无法像阅读论文一样回看前面的幻灯片，因此每一页只承载一个核心信息。以8分钟报告为例，典型的逐页脚本为：标题页、问题引入、方法简介、核心结果（分2页呈现）、讨论与贡献、局限与未来方向、致谢页。

海报、PPT和科普文章这三种形态各自对应不同的信息组织逻辑：海报是浏览式，依赖视觉层级；PPT是听觉线性，依赖重复强化记忆；科普是注意力竞争，需要结论先行。有一条通用底线适用于所有三种形态：\textbf{所有数据图必须从原始分析软件中导出}，不得使用AI生成的图表。AI绘制的图表可能出现数据刻度失真、误差线漏标或坐标轴标注错误等问题。


% --- 传播转化：科普写作与基金申请 ---
\subsection{科普写作与基金申请}

\subsubsection*{\textbf{•科普写作}}

将规律一至规律三落实到具体操作层面，即为科普转化的核心方法。以小林的研究为例，其论文的核心发现为``街景绿视率与居民抑郁症状存在横截面负相关（β=-0.22, p<0.001），心理恢复感中介了62\%的效应''。一个不合格的科普版本会将此写成``家门口的绿色视野能缓解抑郁——看绿色让人更快乐''。在这个转化中，横截面相关的性质被改写为因果关系，横截面设计的限定条件全部丢失，中介机制被简化成了单一因果链条。

科普转化有两条硬性规则。其一，建立并遵守禁忌词清单：不得使用``证明''``证实''``导致''``引发''``引起''``改善''等词汇，代之以``发现''``关联''``相关''``提示''``可能''。这一要求必须在提示词中明确写明，因为AI从学术语体切换至科普语体时，在统计上自然偏向更直接、更少限定的表达方式。其二，在科普文章末尾附加一段``更准确地说''，用2至3句话还原研究的准确科学表述，包括研究设计的核心局限。这一段既是对准确性的兜底，也是对读者的坦诚。

\subsubsection*{\textbf{•基金申请书辅助写作}}

基金申请书与论文的写作逻辑存在根本差异。论文陈述``已经完成了什么''，基金论证``将要做什么以及为什么能够完成''。这一差异决定了AI在基金申请书写作中的角色：它能够帮助研究者从已有成果中提取和整理素材，但核心论证框架必须由研究者自己搭建，因为AI不了解研究者过往的实验条件、团队能力和学科经验。

具体而言，AI在基金申请书写作中可以提供以下几类辅助。其中，研究方案设计（空白识别→假设推导→方案设计→细节把关）的完整AI辅助流程已在5.3中详细讨论，该节内容可直接用于构建基金申请书的``研究方案''模块。本节补充5.3未涉及的其他申请书模块。

\textbf{研究基础素材提取。}将已发表论文摘要提供给AI，由其从中识别与拟研究主题直接相关的发现、可迁移的技术能力、以及已发表论文中的局限——分别对应可行性、延续性和必要性论证。素材经研究者核实后纳入``研究基础''和``立项依据''部分。对应的提示词模板见本节末尾模板7-4。

\textbf{经费预算草稿的生成。}AI可以根据研究者提供的实验方案框架，生成经费预算的初始版本，涵盖设备费、材料费、测试化验加工费、差旅费、出版/文献/信息传播/知识产权事务费等常规科目。但AI不了解研究者所在机构的财务规定（如设备采购限额、差旅标准）和当前市场价格，预算表必须在核实所有单价和限额之后方可使用。

\textbf{团队成员介绍的格式化。}将团队成员的简历信息提供给AI，由AI将其整理为基金申请书要求的叙事格式，突出与项目相关的研究经历和代表性成果。研究者需逐条核实AI改写后的成果清单，防止AI在格式化过程中误改论文标题、期刊名称或发表年份。

\textbf{非技术性段落的规范化。}基金申请书中部分非技术性段落——如``研究意义''``国内外研究现状概述''——的初稿可以由AI辅助生成，但研究者必须对其中涉及学科判断的表述逐一核实，并补充AI因训练数据截止而无法涵盖的最新进展。

需要特别强调的是，\textbf{技术路线的撰写应由研究者亲自完成}。技术路线涉及对实验条件、团队能力和学科经验的综合判断，AI在这些方面不具备有效知识。关于研究方案设计的AI辅助方法已在5.3中详细讨论，此处不赘述。

此外，基金申请书中出现大段与公开模板高度相似的文本可能涉及科研诚信问题，所有AI辅助生成的内容均需经研究者实质性修改后方可提交。

% --- 提示词模板 ---
\subsection*{提示词模板}

\textbf{模板7-1：学术海报设计方案}

\begin{tcolorbox}[promptbox]
我需要在 [会议类型] 展示学术海报。论文信息：
- 研究问题：[一句]  - 核心发现：[≤3条，含数据]
- 创新贡献：[一句]  - 受众：[描述]  - 规格：[尺寸，横/竖]

请提供完整设计方案：分区布局、视觉中心建议、配色方案、文字上限、标题方案。
\end{tcolorbox}

\textbf{模板7-2：口头报告逐页PPT脚本}

\begin{tcolorbox}[promptbox]
我需要做 [N] 分钟口头报告。听众：[描述]。核心信息：[问题、发现、贡献各一句]。

请设计逐页PPT脚本，每页标注：核心信息（仅一个）、视觉建议、时长分配、文字上限。
\end{tcolorbox}

\textbf{模板7-3：论文→科普文章转化}

\begin{tcolorbox}[promptbox]
请将以下学术论文转化为面向 [受众] 的科普文章，[N] 字以内。

要求：
1. 生活化语言，保留科学内核
2. 每个被简化的结论后附加简短的限定条件
3. 不得使用：证明、证实、导致、引发、引起、改善——改用：发现、关联、可能、提示
4. 文末加一段``更准确地说''还原研究的准确科学表述
\end{tcolorbox}

\textbf{模板7-4：基金申请研究基础推导与辅助写作}

\begin{tcolorbox}[promptbox]
以下是我已发表的 [N] 篇论文摘要。拟研究主题：[主题]。

请帮助完成以下任务：
1. 研究基础提取：从已有论文中识别以下三类信息——
   a. 与拟研究主题直接相关的已有发现（支持可行性）
   b. 已有方法中可以迁移或升级的技术能力（展示研究积累）
   c. 已发表论文中明确指出的局限（为拟研究的必要性提供依据）
2. 经费预算草稿：根据以下实验方案框架生成经费预算初始版本，涵盖设备费、材料费、测试费、差旅费、出版费等常规科目。注意：无法确定的单价标注``[待核实]''
3. 团队成员介绍格式化：将以下简历信息整理为基金申请格式，突出与项目相关的经历和成果

输出格式：分三个模块输出。每个模块末尾标注``以上内容需经研究者逐条核实后使用''。
\end{tcolorbox}

使用前提：实验方案框架已确定（参考5.3方案设计部分）。AI输出的经费预算中所有单价和限额均需按所在机构财务规定核实。


% --- 常见错误与规避 ---

\begin{mistakebox}

\textbf{科普文本将``相关''表述为``因果''。}科普转化中必须遵守禁忌词清单，用``发现''``关联''``可能''``提示''替代``证明''``导致''``改善''等因果暗示词汇，并在文末附加``更准确地说''段落还原科学表述。

\smallskip

\textbf{PPT和海报中直接使用了AI生成的数据图。}AI绘制的图表可能出现数据刻度失真、误差线漏标或坐标轴标注错误，所有数据图必须从原始分析软件（如R、Python）中导出。

\smallskip

\textbf{基金申请书的技术路线和模板化文本直接使用AI生成的内容。}技术路线涉及对实验条件和团队能力的综合判断，AI不具备有效知识；AI生成的模板化文本未经实质性修改即提交，可能涉及科研诚信问题。

\end{mistakebox}

\subsection*{小结}
本节涵盖了三类信息传播形态（海报、PPT、科普文章）和基金申请书的辅助写作。知识转化中三条规律贯穿始终：信息可以压缩但科学结论的性质不可改变，话语转换必然伴随精度损失，传播伦理的底线不容突破。科普转化中的禁忌词清单和``更准确地说''段落，是守住底线的两道关键防线。

\bigskip
\subsection*{课后习题}
\begin{enumerate}
\setlength{\itemsep}{6pt}

\item 知识转化中的三条传播规律是什么？为什么``相关不能变因果''是科普写作的底线？

\item 学术海报（浏览式）、口头报告PPT（听觉线性）和科普文章（注意力竞争）各自的信息组织逻辑有什么不同？这一差异如何影响内容的筛选和呈现？

\item 基金申请书与论文的写作逻辑有何根本差异？技术路线部分为什么不应由AI辅助撰写？

\end{enumerate}
\newpage

% ============================================================
% 全章回顾
% ============================================================
\section*{全章回顾}

小林、王工、张博，分属环境健康、机械工程、肿瘤生物学三个学科，表面上毫无交集，但在用AI做科研这件事上，他们遵循的是同一套流程。本章沿这条线索逐节展开：

\begin{enumerate}[label=\arabic*.]
\setlength{\itemsep}{3pt}
\item \textbf{AI工作原理与领域导航（5.1）。}进入新领域的第一天，用AI快速构建知识框架：知识树、核心期刊、研究团队、领域争议。学科不同，操作逻辑相同。
\item \textbf{AI辅助文献检索与阅读（5.2）。}语义检索超越关键词匹配的局限，PICOS框架将模糊兴趣转化为可执行检索，结构化摘要帮助研究者从大量文献中快速筛选核心篇目。阅读完成后，借助AI将论文库、数据库、代码库建立起来，检索与精读的成果转化为可随时调用的知识体系，后续写作和审稿的每一节均受益于这些积累。
\item \textbf{AI辅助课题设计与项目申报（5.3）。}系统识别研究空白：研究者密集的方向创新空间有限，少有涉足的方向风险较高但一旦突破便是独到贡献。推导假设时每条均需附带可证伪条件，方案设计完成后先验证分析方法的可靠性再动笔撰写。该流程的产出直接对接开题报告和基金申请书中``立项依据''与``研究方案''两大部分。
\item \textbf{AI辅助论文写作（5.4）。}从引言到摘要，每一部分都明确区分AI可以辅助的环节与AI不应介入的边界。核心原则是研究者撰写初稿、AI协助修改，不可颠倒。逐句逻辑检测帮助研究者发现自身的推理盲区，AI模拟审稿在投稿前尽可能修复可修复的问题。
\item \textbf{AI辅助语言润色与翻译（5.5）。}分步翻译，一次只追踪一个目标：先语义准确，再语域对齐，最后风格对标。翻译启动前以术语表统一全文译法，完成后通过限定语审计防止语义强度在翻译过程中系统性上移。
\item \textbf{AI辅助自我审稿与投稿准备（5.6）。}投稿前让AI扮演审稿人进行预审，但意见须分层处理：技术报告类可直接修改，实质评价类需核实后再决定，创新判断仅作参考。选择目标期刊时匹配度比影响因子更重要；回复审稿意见时，在哪些问题上让步、在哪些问题上坚持，由研究者做出最终判断。
\item \textbf{AI辅助多模态信息处理与内容创作（5.7）。}论文转化为海报、PPT、科普文章、基金申请书等多种形态。信息可以压缩，但科学结论的性质、方向和确定性不可改变。
\end{enumerate}

\begin{figure}[H]
\centering
\BookFigureImage{images_chapter_05/全章时间线图.png}
\caption{全章时间线图}
\label{fig:ch5_12}
\end{figure}
\par

本章从开题到投稿的完整时间线如图\ref{fig:ch5_12}所示。七节学完，研究者已拥有三库（论文库、数据库、代码库），验证过的假设有了底气，论文有了经过逻辑检测和AI预审的初稿。AI在这一路上做的事情是加速那些重复性最高、最耗时间的环节——文献初筛、方案草图、逻辑检查、语言润色——但每一步的判断和决策仍然在研究者手里。从开题到投稿的总周期取决于研究者的学科：计算和文本分析类的课题，AI加速后也许两三个月可以投出一篇会议论文；需要湿实验或大规模数据采集的课题，周期以年计，AI帮研究者缩短的是文献检索和写作润色等环节，不是实验本身。七节串起来就一句话：\textbf{AI是效率放大器和系统性建议提供者，但研究者的判断无可替代。}


\medskip
\textbf{关于引用来源的说明。}本章中涉及的学术论文引用分为两类：(1) 已核实真实文献（共3篇）——Baron \& Kenny (1986)、van der Maaten \& Hinton (2008)、Dzhambov et al.\ (2018)，均已在正文中以脚注形式给出完整题录信息，并经联网检索验证；(2) 本书教学场景中提及的其他研究者，均为各自领域的真实学者，其理论贡献已在正文中说明。读者在实际研究中应始终遵循本章反复强调的原则：AI推荐的每一篇文献，均需在数据库中独立核实。